% 第七批：原 PDF 物理页354起。
\chapter{镜像定理}\label{chap:mirror-theorem}\index{镜像定理|hyperpage}
\SourcePage{354}
本章旨在证明镜像定理的几个版本，特别着重于五次三维簇。还将涉及与证明技巧有关的一些有趣主题。如第一、二章所指出的，证明需要几个此前尚未讨论的新概念。

\hyperref[chap:quantum-cohomology]{第八章}说明，“镜像定理”有多种表述方式。\hyperref[sec:mirror-symmetry-revisited]{第~\ref*{sec:mirror-symmetry-revisited} 节}的版本用 Hodge 结构变分陈述。这里采取另一种观点，介绍两种密切相关的证明方法，两者都以稳定映射空间上的等变交理论为基础。历史上较早的方法由 Givental 提出 \SourceCite{Givental2}，其特点是使用\hyperref[chap:quantum-differential-equations]{第十章}定义的引力关联数。这些关联数导出\hyperref[subsec:giventals-j-function]{第~\ref*{subsec:giventals-j-function} 小节}定义的 Givental $J$ 函数；我们将看到，$J$ 在 Givental 版本的镜像定理中起重要作用。Givental 的方法可推广到环面完全交 \SourceCite{Givental4}、齐性空间 \SourceCite{Kim2}，并可部分推广到高亏格 \SourceCite{Givental5}。Lian、Liu 和 Yau 的方法 \SourceCite{LLY} 则着重于线性与非线性 $\sigma$ 模型之间的相互作用。在数学上，非线性 $\sigma$ 模型用稳定映射模空间的交理论理解，而线性 $\sigma$ 模型用射影空间 $N_d$ 上的交理论理解；$N_d$ 参数化 $\mathbb P^1$ 上次数 $d$ 的齐次式元组。主要结果用 $N_d$ 的等变交理论表述。这种方法也适用于某些非紧簇。我们先介绍 \SourceCite{LLY} 的方法，因为展开它所需的技术细节较少。

\hyperref[sec:quintic-mirror-theorem]{第~\ref*{sec:quintic-mirror-theorem} 节}概述 \SourceCite{LLY} 对五次三维簇镜像定理的证明，定理在这里陈述为定理~\ref{thm:lly-quintic-mirror}。由\hyperref[subsec:quintic-mirror-symmetry]{第~\ref*{subsec:quintic-mirror-symmetry} 小节}，它蕴含\hyperref[chap:quantum-cohomology]{第八章}用 Hodge 理论表述的镜像对称。在 \SourceCite{LLY} 的方法中，等变版本的 Gromov--Witten 不变量被编码为射影空间 $N_d$ 上的一列等变上同调类 $\widehat Q_d$。我们将证明，它们与等变上同调类 $\widehat P_d$ 有关；后者联系于给出镜像族周期的超几何函数。证明依靠直接比较 $\widehat P_d$ 与 $\widehat Q_d$。为此使用Euler 数据（定义~\ref{def:omega-euler-data}）及其关联关系（定义~\ref{def:linked-euler-data}）这一概念，它本身以乘性上同调类的粘合引理（引理~\ref{lem:euler-data-gluing}）为基础；还要使用镜像变换（定义~\ref{def:euler-data-mirror-transformation}）。也将解释此方法如何适用于\hyperref[subsec:multiple-cover-localization]{第~\ref*{subsec:multiple-cover-localization} 小节}中的多重覆盖计算。

\hyperref[sec:givental-mirror-approach]{第~\ref*{sec:givental-mirror-approach} 节}转向 Givental 的镜像定理方法。这里着重讨论完全交 $X\subset\mathbb P^n$，将其表示为 $\mathcal V=\bigoplus_{i=1}^{\ell}\mathcal O_{\mathbb P^n}(a_i)$（$a_i>0$）的一个截面的零点集。在此情形下，定义两个取形式上同调值的函数 $J_{\mathcal V}$ 和 $I_{\mathcal V}$。函数 $J_{\mathcal V}$ 与该完全交的 Givental $J$ 函数密切相关，而 $I_{\mathcal V}$
\SourcePage{355}
就是\hyperref[subsec:nef-gkz]{第~\ref*{subsec:nef-gkz} 小节}研究的函数。在\hyperref[subsec:givental-nef-projective-mirror]{第~\ref*{subsec:givental-nef-projective-mirror} 小节}，陈述定理~\ref{thm:givental-projective-mirror}，说明当 $\sum_{i=1}^{\ell}a_i\le n+1$ 时，如何用 $I_{\mathcal V}$ 计算 $J_{\mathcal V}$。对于 $\mathcal V=\mathcal O_{\mathbb P^4}(5)$，此镜像定理蕴含五次三维簇的镜像对称。更一般地，在\hyperref[subsec:quantum-hyperplane-section-principle]{第~\ref*{subsec:quantum-hyperplane-section-principle} 小节}考虑完全交 $X\subset Y$，它由 $Y$ 上一个向量丛 $\mathcal V$ 的截面消失所定义，其中 $\mathcal V$ 是线丛的直和。这里将给出 $J_{\mathcal V}$ 的一般定义，并讨论猜想上可以计算 $J_{\mathcal V}$ 的\emph{量子超平面截面原理}。然后证明定理~\ref{thm:givental-projective-mirror}：先在\hyperref[subsec:givental-projective-space-proof]{第~\ref*{subsec:givental-projective-space-proof} 小节}证明 $\mathcal V=0$ 的情形，再在\hyperref[subsec:givental-general-bundle-proof]{第~\ref*{subsec:givental-general-bundle-proof} 小节}证明一般 $\mathcal V$ 的情形。证明使用等变上同调中的局部化，并将它应用于 $J_{\mathcal V}$、$I_{\mathcal V}$ 的等变版本。当 $\mathcal V=0$ 时，\hyperref[chap:localization]{第九章}的局部化公式（推论~\ref{cor:orbifold-localization-integration}）导出 $J_{\mathcal V}$ 局部化的递推关系，将其显式求和便得到 $J_{\mathcal V}$。对于一般 $\mathcal V$，所得关系尚不是递推关系，因此需要更多工作。本章最后讨论 Givental 版本的环面完全交镜像定理。

\section{五次三维簇的镜像定理}\label{sec:quintic-mirror-theorem}
本节遵循 \SourceCite{LLY}，概述五次三维簇镜像定理的证明。
\subsection{定理的陈述}\label{subsec:quintic-mirror-theorem-statement}
设 $V\subset\mathbb P^4$ 为一般的五次三维簇。先回顾 Gromov--Witten 不变量 $\langle I_{0,0,d}\rangle$。固定正整数 $d$，考虑映射
\[
\begin{tikzcd}[column sep=large,row sep=large]
\overline M_{0,1}(\mathbb P^4,d)\arrow[r,"e_1"]\arrow[d,"\pi_1"']&\mathbb P^4\\
\overline M_{0,0}(\mathbb P^4,d)&
\end{tikzcd}
\]
其中 $\pi_1$ 忘掉标记点，$e_1$ 为通常的\IndexLocation{indexloc:0276:355}赋值映射。令 $\mathcal V=\mathcal O_{\mathbb P^4}(5)$，并令 $\mathcal V_d=\pi_{1*}e_1^*(\mathcal V)$，它是 $\overline M_{0,0}(\mathbb P^4,d)$ 上的向量丛。容易验证，它与例题~\ref{ex:quintic-virtual-class} 中定义的丛一致。于是上述 Gromov--Witten 不变量为
\[
\langle I_{0,0,d}\rangle=\int_{\overline M_{0,0}(\mathbb P^4,d)}c_{5d+1}(\mathcal V_d).
\]
\hyperref[chap:gromov-witten]{第七章}用 $N_d$ 表示它，但这里改用 $K_d$（与 \SourceCite{LLY} 一致），因为下面的 $N_d$ 将有另一含义。此外，$c_{5d+1}(\mathcal V_d)$ 为 $\mathcal V_d$ 的 Euler 类，记为 $\operatorname{Euler}(\mathcal V_d)$。用这一记号，上式写成
\[
K_d=\int_{\overline M_{0,0}(\mathbb P^4,d)}\operatorname{Euler}(\mathcal V_d).
\]
注意，这里将积分写在模空间 $\overline M_{0,0}(\mathbb P^4,d)$ 上，而不是代数叠 $\overline{\mathcal M}_{0,0}(\mathbb P^4,d)$ 上。这是因为，虽然 $\mathcal V_d$ 仅是代数叠上的向量丛，但它的 Euler 类在粗模空间上给出定义良好的有理零维闭链，因而上述积分完全有意义。以下常滥用记号，将 $\mathcal V_d$ 当作 $\overline M_{0,0}(\mathbb P^4,d)$ 上的向量丛。

\SourcePage{356}
由\hyperref[chap:gromov-witten]{第七章}还知道，$K_d=\sum_{k\mid d}n_{d/k}k^{-3}$，其中 $n_d$ 是\hyperref[subsec:instanton-numbers]{第~\ref*{subsec:instanton-numbers} 小节}讨论的瞬子数。此外，由\hyperref[subsec:quintic-gw-potential]{第~\ref*{subsec:quintic-gw-potential} 小节}，五次三维簇的 Gromov--Witten 势为
\[
\Phi(T)=\frac56T^3+\sum_{d>0}K_de^{dT},
\]
所以 $\Phi'''$ 为 A 模型关联函数。\hyperref[sec:quintic-mirror-theorem]{第~\ref*{sec:quintic-mirror-theorem} 节}的目标是证明：按照镜像对称的规定，$\Phi$ 由五次镜像族的 Hodge 理论决定，从而证明\hyperref[chap:quintic]{第二章}中给出的 $n_d$ 公式。

为陈述 \SourceCite{LLY} 所证明的结果，考虑取 $H^*(\mathbb P^4,\mathbb C)$ 中形式值的表达式
\[
e^{-Ht/\hbar}\sum_{d\ge0}
\frac{\prod_{m=0}^{5d}(5H-m\hbar)}{\prod_{m=1}^d(H-m\hbar)^5}e^{dt},
\]
其中 $H$ 为超平面类。按 $H$ 的幂展开，得到
\[
5H\left(y_0-y_1\frac H\hbar+y_2\frac{H^2}{\hbar^2}-y_3\frac{H^3}{\hbar^3}\right),
\]
其中 $y_i=y_i(t)$ 是一些函数。在\hyperref[chap:moduli]{第六章}中已经看到，$y_i$ 是五次镜像周期所满足的超几何微分方程
\[
\left[\left(\frac d{dt}\right)^4-5e^t
\left(5\frac d{dt}+1\right)\left(5\frac d{dt}+2\right)
\left(5\frac d{dt}+3\right)\left(5\frac d{dt}+4\right)\right]y=0
\]
的一组解基，用 Frobenius 方法即可验证。例题~\ref{ex:quintic-hypergeometric-mirror-map} 的后续讨论解释了这一点。注意，$y_i$ 与\hyperref[subsec:nef-gkz]{第~\ref*{subsec:nef-gkz} 小节}的函数 $I$ 的系数一致；此函数对应于集合 $\mathcal A=\Delta\cap M$，其中 $\Delta$ 是集合 \eqref{eq:quintic-polar-vertices} 的凸包，该集合由 $\mathbb P^4$ 标准扇的本原整数生成元组成。由命题~\ref{prop:nef-i-gkz}，知道 $y_i$ 满足与 $\mathcal A$ 对应的 GKZ 方程组。

令 $\Psi(t)=y_1/y_0$。\hyperref[subsec:givental-ij]{第~\ref*{subsec:givental-ij} 小节}已指出，镜像映射由 $q=\exp(\Psi(t))$ 定义。于是，\SourceCite{LLY} 的\emph{镜像定理}陈述如下。
\begin{theorem}\label{thm:lly-quintic-mirror}
若 $\Phi$ 为五次三维簇的 Gromov--Witten 势，则
\[
\Phi(\Psi(t))=\frac52\left(\frac{y_1}{y_0}\frac{y_2}{y_0}-\frac{y_3}{y_0}\right).
\]
\end{theorem}
这个五次簇\IndexLocation{indexloc:0788:356}镜像猜想的表述最早见于 \SourceCite{Kontsevich3}。例题~\ref{ex:givental-quintic-mirror-potential} 将说明，定理~\ref{thm:lly-quintic-mirror} 可由 \SourceCite{Givental2} 的结果推出，不过第一个完整而详细的证明见于 \SourceCite{LLY}。\hyperref[subsec:euler-data]{第~\ref*{subsec:euler-data} 小节}解释了，此版本的镜像定理如何导出“经典”镜像对称 \eqref{eq:quintic-mirror-identity}。还应回顾命题~\ref{prop:quintic-mirror-potential}：上述定理等式的右端，是五次镜像的 Hodge 结构变分的“势函数”。如\hyperref[subsec:quintic-mirror-symmetry]{第~\ref*{subsec:quintic-mirror-symmetry} 小节}所指出，定理~\ref{thm:lly-quintic-mirror} 意味着，\IndexLocation{indexloc:0759:356}五次镜像与五次三维簇构成定义~\ref{def:mathematical-mirror-pair} 意义下的数学镜像对。

\SourcePage{357}
本节余下部分用于证明定理~\ref{thm:lly-quintic-mirror}；\hyperref[sec:givental-mirror-approach]{第~\ref*{sec:givental-mirror-approach} 节}将概述此定理的第二种证明。
\subsection{稳定模空间的线性模型}\label{subsec:stable-moduli-linear-models}\index{稳定映射!线性模型|hyperpage}
下面引入稳定映射空间的一个“线性化”版本，它在后续讨论中起关键作用。在物理上，这类似于从非线性 $\sigma$ 模型转到线性 $\sigma$ 模型。先对 $\mathbb P^n$ 中次数 $\ell$ 的超曲面建立理论，再特化到 $\mathbb P^4$ 中的五次三维簇。

先紧化次数 $d$ 的映射 $\mathbb P^1\to\mathbb P^n$ 所组成的空间。通常的模空间 $\overline M_{0,0}(\mathbb P^n,d)$ 并不实现这种紧化，因为它的一般元素只能在相差源空间重新参数化的意义下确定这样的映射。因此改用空间
\[
M_d=\overline M_{0,0}(\mathbb P^1\times\mathbb P^n,(1,d)).
\]
它的一般元素是一个 $\mathbb P^1\to\mathbb P^1\times\mathbb P^n$ 的映射，像的次数为 $(1,d)$。关键在于，重新参数化源空间不会改变其像，而这个像是一个次数 $d$ 的映射 $\mathbb P^1\to\mathbb P^n$ 的图像。因此，$M_d$ 确实紧化了这些映射组成的空间。

第二种紧化如下。一个次数 $d$ 的映射 $\mathbb P^1\to\mathbb P^n$ 可由 $\mathbb P^1$ 上的 $n+1$ 个次数 $d$ 的齐次多项式 $(\alpha_0,\ldots,\alpha_n)$ 表示。这导出射影空间
\[
N_d=\mathbb P\bigl(H^0(\mathbb P^1,\mathcal O_{\mathbb P^1}(d))^{n+1}\bigr),
\]
其中 $\mathbb P$ 表示一维子空间组成的射影空间。\hyperref[sec:gauged-linear-sigma-models]{附录 B.5} 在 $n=4$ 的情形下解释了 $N_d$ 在线性 $\sigma$ 模型中的作用。

$M_d$ 和 $N_d$ 均有自然的环面作用。以下大量使用第 \hyperref[sec:localization-theorem]{\ref*{sec:localization-theorem}}、\hyperref[sec:stable-map-localization]{\ref*{sec:stable-map-localization}} 节，读者继续之前可以先复习。命题~\ref{prop:atiyah-bott-localization} 及其推论、例题~\ref{ex:projective-equivariant-localization} 和\hyperref[subsec:kontsevich-localization]{第~\ref*{subsec:kontsevich-localization} 小节}尤其重要。

令 $G=\mathbb C^*\times T$ 为环面，其中 $T=(\mathbb C^*)^{n+1}$。先描述要使用的 $\mathbb P^1\times\mathbb P^n$ 上的 $G$ 作用。令 $t\in\mathbb C^*$、$(t_0,\ldots,t_n)\in T$，则
\[
\begin{aligned}
&(t,t_0,\ldots,t_n)\cdot((w_0,w_1),(x_0,\ldots,x_n))\\
&\qquad=((t^{-1}w_0,w_1),(t_0^{-1}x_0,\ldots,t_n^{-1}x_n)).
\end{aligned}
\]
如例题~\ref{ex:projective-equivariant-localization} 一样，选择逆元是为了让经由 $G$ 作用拉回所诱导的齐次式上的作用具有正指数。

接着描述 $G$ 在 $M_d=\overline M_{0,0}(\mathbb P^1\times\mathbb P^n,(1,d))$ 上的作用，它由 $\mathbb P^1\times\mathbb P^n$ 上的 $G$ 作用诱导。这与\hyperref[sec:stable-map-localization]{第~\ref*{sec:stable-map-localization} 节}研究的 $(\mathbb C^*)^{r+1}$ 在 $\overline M_{0,n}(\mathbb P^r,d)$ 上的作用密切相关。回顾，我们已仔细描述了固定点轨迹的各分支。这里关心的是 $G$ 在 $M_d$ 上作用的固定点轨迹的分支；可得到类似描述，只需作以下修改：
\begin{itemize}
\item 一个 $G$ 不动稳定映射的收缩分支必须映到 $\mathbb P^1\times\mathbb P^n$ 的某个 $G$ 不动点。
\item 非收缩分支必须覆盖 $\mathbb P^1\times\mathbb P^n$ 中的某条 $G$ 不变曲线。这些曲线形如 $\mathbb P^1\times\{q_i\}$ 和 $\{p\}\times\ell$，其中 $q_i\in\mathbb P^n$ 为 $T$ 不动点，$p$ 为 $0$ 或 $\infty$，$\ell\subset\mathbb P^n$ 为坐标直线。
\item 若 $\sum'$ 表示对覆盖 $\mathbb P^1\times\{q_i\}$ 的分支求和，$\sum''$ 表示对覆盖 $\{p\}\times\ell$ 的分支求和，则 $\sum'd_i=1$、$\sum''d_i=d$，其中 $d_i$ 为各覆盖的次数。
\end{itemize}
\SourcePage{358}
这些 $G$ 固定点轨迹的分支，与引理~\ref{lem:graph-space-descendant-factorization} 的证明中出现的分支类似。

最后描述 $G$ 如何作用于 $N_d$。此处作用为
\[
\begin{aligned}
&(t,t_0,\ldots,t_n)\cdot(\alpha_0(w_0,w_1),\ldots,\alpha_n(w_0,w_1))\\
&\qquad=(t_0^{-1}\alpha_0(t^{-1}w_0,w_1),\ldots,
 t_n^{-1}\alpha_n(t^{-1}w_0,w_1)).
\end{aligned}
\]
这给出射影空间 $N_d$ 上的一个线性化 $G$ 作用。$N_d$ 的 $G$ 不动点为
\[
p_{i,r}=(0,\ldots,0,w_0^rw_1^{d-r},0,\ldots,0),
\]
其中非零单项式位于第 $i$ 个位置。以 $i_{p_{i,r}}$ 表示 $p_{i,r}$ 到 $N_d$ 的包含映射。

如\hyperref[subsec:equivariant-cohomology]{第~\ref*{subsec:equivariant-cohomology} 小节}一样，令 $\lambda_0,\ldots,\lambda_n$ 生成 $H^*(BT)$，并类似地引入 $H^*(B\mathbb C^*)$ 的生成元 $\hbar$。\SourceCite{LLY} 将这个生成元记为 $\alpha$，但这里采用 \SourceCite{Givental2} 的记号，以使全章记号一致（不过 \SourceCite{Givental2} 中 $\hbar$ 的符号不同，对应于 $\mathbb C^*$ 作用的不同约定）。

可用 $\lambda_i$、$\hbar$ 描述等变上同调 $H_G^*(N_d)$，答案类似于例题~\ref{ex:projective-equivariant-localization}。将 $p_{i,r}$ 视为 $H^0(\mathbb P^1,\mathcal O_{\mathbb P^1}(d))^{n+1}$ 中的向量，就使 $G$ 在此空间上的作用对角化。由 $G$ 在 $N_d$ 上的作用，作用于 $p_{i,r}$ 的权重显然为 $-(\lambda_i+r\hbar)$。因此，只要将权重 $\{\lambda_i\}$ 换成 $\{\lambda_i+r\hbar\}$，就得到例题~\ref{ex:projective-equivariant-localization} 中的同样公式。特别地，$(N_d)_G$ 是 $BG$ 上丛 $\bigoplus_{i,r}\mathcal O(-\lambda_i-r\hbar)$ 的射影化。若 $\kappa$ 为 $N_d$ 的等变超平面类，则与式~\eqref{eq:projective-equivariant-relation} 对应的关系为
\begin{equation}\label{eq:linear-model-equivariant-cohomology}
\prod_{i,r}(\kappa-\lambda_i-r\hbar)=0
\end{equation}
（在 $H_G^*(N_d)$ 中）。此外，作为式~\eqref{eq:projective-fixed-dual-class} 的特殊情形，令
\begin{equation}\label{eq:linear-model-fixed-point-class}
\phi_{p_{i,r}}=\prod_{(j,s)\ne(i,r)}\bigl(\kappa-(\lambda_j+s\hbar)\bigr),
\end{equation}
因为 $\lambda_j+s\hbar$ 是 $N_d$ 线性化 $G$ 作用的相应权重参数。回顾，将非零权重取逆后，$\phi_{p_{i,r}}$ 构成 $\mathcal R_GH_G^*(N_d)$ 作为 $\mathcal R_G$ 向量空间的一组基。

注意，每个类 $Q\in H_G^*(N_d)$ 都可由 $\kappa,\hbar,\lambda_i$ 的多项式表示。上述讨论结合式~\eqref{eq:projective-hyperplane-fixed-restriction} 表明，类 $i_{p_{i,r}}^*(Q)$ 就是在 $\kappa$ 中代入 $\lambda_i+r\hbar$ 后所得的 $H^*(BG)=\mathbb C[\hbar,\lambda_i]$ 中的元素。因此，也可合理地以 $Q(\lambda_i+r\hbar)$ 表示 $i_{p_{i,r}}^*(Q)$。由关系~\eqref{eq:linear-model-equivariant-cohomology}，此记号定义良好。当 $Q$ 属于上同调环的某个适当局部化（交换代数意义下）时，也使用类似记号。

作为局部化定理应用于这一情形的简单例子，证明
\begin{equation}\label{eq:linear-model-restriction-integral}
Q(\lambda_i+r\hbar)=i_{p_{i,r}}^*(Q)
=\int_{(N_d)_G}\phi_{p_{i,r}}\cup Q,
\end{equation}
其中 $\int_{(N_d)_G}$ 为推论~\ref{cor:localization-integration} 之前定义的等变积分。为证明此式，先注意，若 $N_{i,r}$ 为 $p_{i,r}$ 的法丛，$\operatorname{Euler}_G(N_{i,r})$ 为其等变 Euler 类，则由式~\eqref{eq:projective-fixed-euler}，
\SourcePage{359}
\begin{equation}\label{eq:linear-model-normal-euler}
\operatorname{Euler}_G(N_{i,r})
=\prod_{(j,s)\ne(i,r)}\bigl(\lambda_i+r\hbar-(\lambda_j+s\hbar)\bigr).
\end{equation}
此外，将式~\eqref{eq:projective-dual-class-restriction} 适配到 $N_d$，容易得到
\begin{equation}\label{eq:linear-model-fixed-class-restriction}
i_{p_{j,s}}^*(\phi_{p_{i,r}})=
\begin{cases}\operatorname{Euler}_G(N_{i,r}),&(j,s)=(i,r),\\0,&\text{其他情形}.
\end{cases}
\end{equation}
现在应用推论~\ref{cor:localization-integration}，得到
\[
\int_{(N_d)_G}\phi_{p_{i,r}}\cup Q
=\frac{i_{p_{i,r}}^*(\phi_{p_{i,r}}\cup Q)}{\operatorname{Euler}_G(N_{i,r})}
=i_{p_{i,r}}^*(Q).
\]

应指出，\hyperref[chap:localization]{第九章}中的局部化定理，使用的是\IndexLocation{indexloc:0247:359}等变上同调在交换代数意义下的局部化。如\hyperref[subsec:atiyah-bott-localization]{第~\ref*{subsec:atiyah-bott-localization} 小节}一样，以 $\mathcal R_T$ 表示 $H^*(BT)\simeq\mathbb C[\lambda_i]$ 的分式域，以 $\mathcal R_TH_G^*(N_d)$ 表示张量积 $H_G^*(N_d)\otimes_{H^*(BT)}\mathcal R_T$。这是将 $\lambda_i$ 的非零多项式取逆后得到的局部化。类似地，若 $\mathcal R_G$ 为 $H^*(BG)\simeq\mathbb C[\lambda_i,\hbar]$ 的分式域，则得到局部化 $\mathcal R_GH_G^*(N_d)$，其中 $\lambda_i$ 和 $\hbar$ 的非零多项式均取逆。

现在断言，存在自然的 $G$ 等变态射 $\varphi:M_d\to N_d$，描述如下。考虑双次数 $(1,d)$ 的映射 $f:\mathbb P^1\to\mathbb P^1\times\mathbb P^n$，它表示 $M_d$ 的一个元素。必要时重新参数化源空间 $\mathbb P^1$，可用以下形式的映射表示 $f\in M_d$：
\begin{equation}\label{eq:graph-space-standard-parametrization}
(w_0,w_1)\longmapsto\bigl((w_1,w_0),
(\alpha_0(w_0,w_1),\ldots,\alpha_n(w_0,w_1))\bigr),
\end{equation}
其中 $\alpha_i$ 是 $\mathbb P^1$ 坐标 $(w_0,w_1)$ 中次数 $d$ 的齐次式。对于这样的 $f$，定义 $\varphi(f)=(\alpha_0,\ldots,\alpha_n)$。概括地说，一个参数化的次数 $d$ 的映射，可由 $\mathbb P^1$ 上 $n+1$ 个无公因子的次数 $d$ 的齐次式模去共同非零标量倍来表示。

之所以看似奇怪地将 $(w_0,w_1)$ 换为 $(w_1,w_0)$，是因为要求 $\varphi$ 为 $\mathbb C^*$ 等变映射。对于 $t\in\mathbb C^*$，$t\cdot f$ 为
\[
(w_0,w_1)\longmapsto\bigl((t^{-1}w_1,w_0),(\alpha_0,\ldots,\alpha_n)\bigr)
=\bigl((w_1,tw_0),(\alpha_0,\ldots,\alpha_n)\bigr).
\]
为计算 $\varphi(t\cdot f)$，须将 $t\cdot f$ 写为标准形式~\eqref{eq:graph-space-standard-parametrization}：用 $t^{-1}w_0$ 替换 $w_0$ 即可。结果为
\[
\varphi(t\cdot f)
=(\alpha_0(t^{-1}w_0,w_1),\ldots,\alpha_n(t^{-1}w_0,w_1)),
\]
它等于 $t\cdot\varphi(f)$，如所断言。

对于任意稳定映射 $f$，$\varphi(f)=(\alpha_0,\ldots,\alpha_n)$ 的描述来自对收缩分支的研究。令 $p_1:\mathbb P^1\times\mathbb P^n\to\mathbb P^1$、$p_2:\mathbb P^1\times\mathbb P^n\to\mathbb P^n$ 为两个投影，并假设 $p_1\circ f$ 将一个分支收缩到点 $p\in\mathbb P^1$，而 $p_2\circ f$ 在此分支上的次数为 $d'$。此时可证明，$\alpha_i$ 有公因子 $L(w_0,w_1)^{d'}$，其中 $L$ 是在按式~\eqref{eq:graph-space-standard-parametrization} 的坐标交换对应于 $p$ 的点处消失的线性函数。下面定理~\ref{thm:geometric-euler-data-gluing} 的证明将使用此事实。$\varphi$ 的更明确公式见 \SourceCite{Givental2} 或 \SourceCite{LLY}。这些文献也证明了 $\varphi$ 是 $M_d$ 上的态射（\SourceCite{LLY} 所给的证明出自 J. Li）。\SourceCite{Givental2} 实际构造的是映射 $L_d\to N_d$，其中 $L_d$ 与 $M_d$ 密切相关；将在引理~\ref{lem:complete-intersection-double-construction} 的证明中定义 $L_d$。

\SourcePage{360}
\subsection{Euler 数据}\label{subsec:lly-euler-data}\index{Euler 数据|hyperpage}
固定整数 $\ell\ge1$，令 $\mathcal V=\mathcal O_{\mathbb P^n}(\ell)$。如第~\ref{subsec:quintic-mirror-theorem-statement} 小节一样，有图
\begin{equation}\label{eq:euler-data-universal-curve}
\begin{tikzcd}[column sep=large,row sep=large]
\overline M_{0,1}(\mathbb P^n,d)\arrow[r,"e_1"]\arrow[d,"\pi_1"']&\mathbb P^n\\
\overline M_{0,0}(\mathbb P^n,d)&
\end{tikzcd}
\end{equation}
并定义 $\mathcal V_d=\pi_{1*}e_1^*(\mathcal V)$。这是代数叠 $\overline{\mathcal M}_{0,0}(\mathbb P^n,d)$ 上的丛，不过与第~\ref{subsec:quintic-mirror-theorem-statement} 小节一样，滥用记号，将它视为 $\overline M_{0,0}(\mathbb P^n,d)$ 上的丛。

现在定义几个重要的等变上同调类。考虑自然投影 $\pi:M_d\to\overline M_{0,0}(\mathbb P^n,d)$，定义
\[
\chi_d^{\mathcal V}=\operatorname{Euler}_G(\pi^*(\mathcal V_d))\in H_G^*(M_d).
\]
利用上述映射 $\varphi:M_d\to N_d$，定义
\begin{equation}\label{eq:geometric-euler-data}
\widehat Q_d=\varphi_!(\chi_d^{\mathcal V})\in H_G^*(N_d),
\end{equation}
即 $\chi_d^{\mathcal V}$ 的等变推前。可将 $\widehat Q_d$ 看作 $\chi_d^{\mathcal V}$ 的“线性化”。接着定义
\begin{equation}\label{eq:hypergeometric-euler-data}
\widehat P_d=\prod_{m=0}^{\ell d}(\ell\kappa-m\hbar)\in H_G^*(N_d).
\end{equation}
$\widehat P_d$、$\widehat Q_d$ 的定义均对应于丛 $\mathcal V=\mathcal O_{\mathbb P^n}(\ell)$。在 $\mathcal V=\mathcal O_{\mathbb P^4}(5)$ 的特殊情形下，这两个等变上同调类将在定理~\ref{thm:lly-quintic-mirror} 的证明中起核心作用。

\SourceCite{LLY} 对比这里广得多的一类丛 $\mathcal V$ 定义了 $\widehat P_d$ 和 $\widehat Q_d$。若将 $\mathcal O_{\mathbb P^n}(\ell)$ 看作与超曲面相关联的丛，则 \SourceCite{LLY} 所考虑的丛，不但包含与 $\mathbb P^n$ 中完全交相关联的丛，也包含\emph{凸凹丛}。下面\hyperref[subsec:hypergeometric-mirror-transformations]{第~\ref*{subsec:hypergeometric-mirror-transformations} 小节}末的评论中将定义凸凹丛。

按照目前的构造，丛 $\mathcal V=\mathcal O_{\mathbb P^n}(\ell)$ 导出无穷列丛 $\{\mathcal V_d\}_{d=1}^{\infty}$。也定义 $d=0$ 的情形：以 $\mathcal V_0$ 表示 $\mathbb P^n$ 上的丛 $\mathcal V$。须注意，虽然这些丛位于不同空间上，它们之间仍有非平凡关系。为说明这一点，取 $0<r<d$，用拉回图定义 $\overline M(r,d-r)$：
\begin{equation}\label{eq:euler-data-gluing-fiber-product}
\begin{tikzcd}[column sep=large,row sep=large]
\overline M(r,d-r)\arrow[r,"p_1"]\arrow[d,"p_2"']&\overline M_{0,1}(\mathbb P^n,r)\arrow[d,"e_1"]\\
\overline M_{0,1}(\mathbb P^n,d-r)\arrow[r,"e_1"']&\mathbb P^n
\end{tikzcd}
\end{equation}
其中 $e_1$ 为通常的赋值映射，$p_i$ 为投影。可将 $\overline M(r,d-r)$ 看作由单标记稳定映射对 $((f_1,C_1,p_1),\allowbreak(f_2,C_2,p_2))$ 组成，它们满足 $f_1(p_1)=f_2(p_2)$。于是有图
\begin{equation}\label{eq:euler-data-gluing-map}
\begin{tikzcd}[column sep=large,row sep=large]
\overline M(r,d-r)\arrow[r,"\psi"]\arrow[d,"e_1"']&\overline M_{0,0}(\mathbb P^n,d)\\
\mathbb P^n&
\end{tikzcd}
\end{equation}
\SourcePage{361}
其中 $e_1$ 为赋值映射（由 $\overline M(r,d-r)$ 的定义可知它定义良好）。水平映射将 $((f_1,C_1,p_1),\allowbreak(f_2,C_2,p_2))$ 映到无标记稳定映射 $(f,C)$；这里 $C$ 由 $C_1\cup C_2$ 识别 $p_1$ 和 $p_2$ 后得到，且 $f|_{C_i}=f_i$（同样，由 $\overline M(r,d-r)$ 的定义可知它定义良好）。

当 $r=0$ 或 $r=d$ 时，须修改 $\overline M(r,d-r)$ 的定义。由于 $\overline M_{0,1}(\mathbb P^n,0)$ 为空，在图~\eqref{eq:euler-data-gluing-fiber-product} 中以恒等映射 $1_{\mathbb P^n}:\mathbb P^n\to\mathbb P^n$ 替换 $e_1:\overline M_{0,1}(\mathbb P^n,0)\to\mathbb P^n$。对于 $r=0$，得到 $\overline M(0,d)=\overline M_{0,1}(\mathbb P^n,d)$；在图~\eqref{eq:euler-data-gluing-fiber-product} 中，$p_1$ 识别为赋值映射 $e_1$，$p_2$ 为恒等映射。此外，图~\eqref{eq:euler-data-gluing-map} 中的 $\psi$ 识别为忘掉标记点的映射 $\pi_1:\overline M_{0,1}(\mathbb P^n,d)\to\overline M_{0,0}(\mathbb P^n,d)$。$r=d$ 的情形类似。

在这些准备下，得到以下\emph{粘合引理}。
\begin{lemma}\label{lem:euler-data-gluing}\index{Euler 数据!粘合引理|hyperpage}
若 $b$ 是向量丛的任意乘性等变特征类（例如等变 Euler 类），则在 $\overline M(r,d-r)$ 上有
\[
e_1^*b(\mathcal V)\,\psi^*b(\mathcal V_d)
=p_1^*\pi_1^*b(\mathcal V_r)\,p_2^*\pi_1^*b(\mathcal V_{d-r}).
\]
当 $r=0$ 或 $r=d$ 时，为使该式有意义，所有出现的 $\pi_1^*b(\mathcal V_0)$ 都须解释为 $b(\mathcal V_0)=b(\mathcal V)$。
\end{lemma}
\begin{proof}
当 $r=0,d$ 时，引理显然成立。设 $0<r<d$，并设 $((f_1,C_1,p_1),\allowbreak(f_2,C_2,p_2))$ 如上映到 $(f,C)$。有正合列
\[
0\longrightarrow f^*\mathcal V\longrightarrow
f_1^*\mathcal V\oplus f_2^*\mathcal V\longrightarrow\mathcal V|_p\longrightarrow0,
\qquad p=f_1(p_1)=f_2(p_2).
\]
于是得到正合列
\[
0\longrightarrow H^0(C,f^*\mathcal V)\longrightarrow
H^0(C_1,f_1^*\mathcal V)\oplus H^0(C_2,f_2^*\mathcal V)
\longrightarrow\mathcal V|_p\longrightarrow0.
\]
由此即可得到引理。
\end{proof}

粘合引理给出等变 Euler 类 $\operatorname{Euler}_G(\mathcal V_d)\in H_G^*(\overline M_{0,0}(\mathbb P^n,d))$ 之间的恒等式。由于 $\widehat Q_d\in H_G^*(N_d)$ 用 $\operatorname{Euler}_G(\mathcal V_d)$ 定义，自然会问这些类是否满足类似恒等式。$H_G^*(N_d)$ 中的类容易研究：固定点集很简单，仅由孤立点 $p_{i,r}$ 组成。因此，由\IndexLocation{indexloc:0739:361}局部化定理，类 $\widehat Q_d$ 被限制 $i_{p_{i,r}}^*(\widehat Q_d)$ 唯一确定。这提示我们寻找这些限制所满足的、类似于引理~\ref{lem:euler-data-gluing} 的恒等式。

下面的定义~\ref{def:omega-euler-data} 将使这一点明确。不过，在陈述定义之前，需要映射 $\overline{\phantom\alpha}:N_d\to N_d$，定义为
\[
(\alpha_0(w_0,w_1),\ldots,\alpha_n(w_0,w_1))
\longmapsto(\alpha_0(w_1,w_0),\ldots,\alpha_n(w_1,w_0)).
\]
此映射显然 $T$ 等变，但不 $\mathbb C^*$ 等变。\SourceCite{LLY} 证明，配合 $\mathbb C^*$ 参数反转，它在等变上同调上给出的变换 $\overline{\phantom\alpha}:H_G^*(N_d)\to H_G^*(N_d)$ 满足
\begin{equation}\label{eq:euler-data-involution}
\overline\kappa=\kappa-d\hbar,\qquad\overline\hbar=-\hbar,
\qquad\overline{\lambda_i}=\lambda_i.
\end{equation}
还回顾例题~\ref{ex:projective-equivariant-localization}：$G$（或 $T$）在 $\mathbb P^n$ 上作用的不动点记为 $q_i$。

\SourcePage{362}
\begin{definition}\label{def:omega-euler-data}\index{Euler 数据!定义|hyperpage}
设 $\Omega\in\mathcal R_TH_T^*(\mathbb P^n)$ 为可逆类。一列 $Q=\{Q_d\}$，其中 $Q_d\in\mathcal R_TH_G^*(N_d)$（$d>0$），若对所有 $r=0,\ldots,d$ 及 $i=0,\ldots,n$ 均有
\[
i_{q_i}^*(\Omega)i_{p_{i,r}}^*(Q_d)
=\overline{i_{p_{i,0}}^*(Q_r)}\,i_{p_{i,0}}^*(Q_{d-r}),
\]
则称它为一组 \emph{$\Omega$--Euler 数据}，其中 $Q_0=\Omega$。所有 $\Omega$--Euler 数据组成的集合记为 $\mathcal A^\Omega$。
\end{definition}
前面已指出，$i_{p_{i,r}}^*(Q_d)=Q_d(\lambda_i+r\hbar)$；由式~\eqref{eq:projective-polynomial-fixed-restriction}，有 $i_{q_i}^*(\Omega)=\Omega(\lambda_i)$。因此，上式可写成
\[
\Omega(\lambda_i)Q_d(\lambda_i+r\hbar)
=\overline{Q_r(\lambda_i)}\,Q_{d-r}(\lambda_i).
\]
注意，由于 $\Omega\in\mathcal R_TH_T^*(\mathbb P^n)$，有 $\overline\Omega=\Omega$，故 $r=0$ 时此式显然成立。下面将证明，由式~\eqref{eq:geometric-euler-data} 和~\eqref{eq:hypergeometric-euler-data} 定义的序列 $\widehat Q=\{\widehat Q_d\}$、$\widehat P=\{\widehat P_d\}$ 均为 $\Omega=\operatorname{Euler}_T(\mathcal V)=\operatorname{Euler}_T(\mathcal O_{\mathbb P^n}(\ell))=\ell p$ 的 Euler 数据，其中 $p$ 为等变超平面类。采用“Euler 数据”这一名称，部分是因为 $\widehat Q$ 用等变 Euler 类定义。

先给出 $\widehat P$ 是一组 $\ell p$--Euler 数据的简单证明。
\SourceExampleNumber{11.1.3}{1}
\begin{example}\label{ex:hypergeometric-euler-data-gluing}
先注意，$\Omega=\ell p$ 在 $\mathcal R_TH_T^*(\mathbb P^n)$ 中可逆。由局部化（命题~\ref{prop:atiyah-bott-localization}）即可看出这一点，因为限制 $i_{q_i}^*(\Omega)=\Omega(\lambda_i)=\ell\lambda_i$ 均可逆。还回顾式~\eqref{eq:hypergeometric-euler-data}，有 $\widehat P_d=\prod_{m=0}^{\ell d}(\ell\kappa-m\hbar)$。于是，$\widehat P=\{\widehat P_d\}$ 为 $\ell p$--Euler 数据，等价于恒等式
\[
\ell\lambda_i\prod_{m=0}^{\ell d}\bigl(\ell(\lambda_i+r\hbar)-m\hbar\bigr)
=\prod_{m=0}^{\ell r}(\ell\lambda_i+m\hbar)
\prod_{m=0}^{\ell(d-r)}(\ell\lambda_i-m\hbar),
\]
而此式立即可验证。
\end{example}
证明 $\widehat Q$ 是 $\ell p$--Euler 数据需要更多工作。
\begin{theorem}\label{thm:geometric-euler-data-gluing}
序列 $\widehat Q=\{\widehat Q_d\}=\{\varphi_!\chi_d^{\mathcal V}\}$ 是一组 $\ell p$--Euler 数据。
\end{theorem}
\begin{proof}
须证明，对所有 $r,i$，有
\begin{equation}\label{eq:geometric-euler-data-gluing-target}
\ell\lambda_i\widehat Q_d(\lambda_i+r\hbar)
=\overline{\widehat Q_r(\lambda_i)}\,\widehat Q_{d-r}(\lambda_i).
\end{equation}
先计算 $i_{p_{i,r}}^*(\widehat Q_d)$，它为
\begin{equation}\label{eq:geometric-euler-data-localization-integral}
\widehat Q_d(\lambda_i+r\hbar)
=\int_{(N_d)_G}\phi_{p_{i,r}}\widehat Q_d
=\int_{(M_d)_G}\varphi^*(\phi_{p_{i,r}})\chi_d^{\mathcal V}.
\end{equation}
第一个等号来自式~\eqref{eq:linear-model-restriction-integral}，第二个等号来自 $\widehat Q_d$ 的定义和投影公式。策略是用推论~\ref{cor:localization-integration} 的局部化公式计算式~\eqref{eq:geometric-euler-data-localization-integral} 右端的积分。

先确定 $M_d$ 中哪些 $G$ 不动分支对此积分有贡献。令 $i_\Gamma:\overline M_\Gamma\subset M_d$ 表示与图 $\Gamma$ 对应的分支。将推论~\ref{cor:localization-integration} 应用于此积分时，$\overline M_\Gamma$ 的贡献项含因子
\[
i_\Gamma^*(\varphi^*(\phi_{p_{i,r}}))
=(\varphi|_{\overline M_\Gamma})^*\bigl(i_{\varphi(\overline M_\Gamma)}^*(\phi_{p_{i,r}})\bigr).
\]
\SourcePage{363}
由于 $\varphi(\overline M_\Gamma)$ 必为某个不动点 $p_{j,s}$，由式~\eqref{eq:linear-model-fixed-point-class} 中 $\phi_{p_{i,r}}$ 的定义，当 $(j,s)\ne(i,r)$ 时，相应项为零。

剩下的是确定在 $\varphi$ 下映到 $p_{i,r}$ 的 $G$ 固定点集分支。为此使用式~\eqref{eq:euler-data-gluing-map} 中 $\psi$ 的一个变体。给定 $\overline M(r,d-r)$ 的一点，它由满足 $f_1(p_1)=f_2(p_2)=q\in\mathbb P^n$ 的 $f_j:(C_j,p_j)\to\mathbb P^n$ 表示。定义 $C=C_0\cup C_1\cup C_2$，其中 $C_0=\mathbb P^1$，将 $0\in\mathbb P^1$ 与 $p_1\in C_1$ 粘合，将 $\infty\in\mathbb P^1$ 与 $p_2\in C_2$ 粘合。然后定义 $f:C\to\mathbb P^1\times\mathbb P^n$ 为
\[
f|_{C_0}(z)=(z,q),\qquad f|_{C_1}(z)=(0,f_1(z)),
\qquad f|_{C_2}(z)=(\infty,f_2(z)).
\]
于是 $f\in M_d$，故这一过程定义映射
\begin{equation}\label{eq:euler-data-graph-gluing-map}
\widetilde\psi:\overline M(r,d-r)\longrightarrow M_d.
\end{equation}
令 $\{F_k\}$ 为 $\overline M_{0,1}(\mathbb P^n,k)$ 的 $T$ 不动分支中、标记点映到 $q_i$ 的分支集合。于是 $F_r\times F_{d-r}\subset\overline M(r,d-r)$；将其与 $\widetilde\psi$ 下的像识别，便写作 $F_r\times F_{d-r}\subset M_d$。

断言：遍历所有可能的 $F_r$ 和 $F_{d-r}$ 所得到的 $F_r\times F_{d-r}$，恰好就是在 $\varphi$ 下映到 $p_{i,r}$ 的 $G$ 固定点集分支。为证明此断言，取 $f\in F_r\times F_{d-r}$。按定义，$p_1\circ f$ 收缩 $C_1$、$C_2$，而 $p_2\circ f$ 在 $C_1$、$C_2$ 上的次数分别为 $r$、$d-r$。于是，第~\ref{subsec:stable-moduli-linear-models} 小节对 $\varphi(f)$ 的描述给出 $\varphi(f)=p_{i,r}$。利用 \SourceCite{LLY} 中更详细的 $\varphi$ 描述，还可证明 $F_r\times F_{d-r}$ 是 $G$ 固定点集的一个分支，且 $\varphi^{-1}(p_{i,r})$ 中的所有 $G$ 不动分支都由此得到。

因此，用局部化计算式~\eqref{eq:geometric-euler-data-localization-integral} 右端积分时，只需要形如 $F_r\times F_{d-r}$ 的分支。先处理 $r\ne0,d$ 的情形。以下记号较为有用。令 $\mathcal L_{k,1}$ 为 $\overline M_{0,1}(\mathbb P^n,k)$ 上的线丛，其在 $f:(C,p)\to\mathbb P^n$ 处的纤维是余切空间 $T_p^*(C)$，如\hyperref[subsec:gravitational-correlator-definition]{第~\ref*{subsec:gravitational-correlator-definition} 小节}开头所定义。对任意分支 $F_k$，以 $N(F_k)$ 表示 $F_k$ 在 $\overline M_{0,1}(\mathbb P^n,k)$ 中的法丛。最后，令 $\pi_1:\overline M_{0,1}(\mathbb P^n,k)\to\overline M_{0,0}(\mathbb P^n,k)$ 为图~\eqref{eq:euler-data-universal-curve} 中定义丛 $\mathcal V_k$ 时使用的映射。将推论~\ref{cor:orbifold-localization-integration} 应用于式~\eqref{eq:geometric-euler-data-localization-integral} 右端积分，断言
\begin{equation}\label{eq:geometric-euler-data-product-localization}
\begin{aligned}
\widehat Q_d(\lambda_i+r\hbar)
&=\int_{(M_d)_G}\varphi^*(\phi_{p_{i,r}})\chi_d^{\mathcal V}\\
&=-\hbar^{-2}(\ell\lambda_i)^{-1}i_{p_{i,r}}^*(\phi_{p_{i,r}})
\prod_{j\ne i}(\lambda_i-\lambda_j)\\
&\quad\times\sum_{F_r}\int_{(F_r)_G}
\frac{\pi_1^*\operatorname{Euler}_G(\mathcal V_r)}
{a_r\operatorname{Euler}_G(N(F_r))(\hbar-c_1^G(\mathcal L_{r,1}))}\\
&\quad\times\sum_{F_{d-r}}\int_{(F_{d-r})_G}
\frac{\pi_1^*\operatorname{Euler}_G(\mathcal V_{d-r})}
{a_{d-r}\operatorname{Euler}_G(N(F_{d-r}))(-\hbar-c_1^G(\mathcal L_{d-r,1}))}.
\end{aligned}
\end{equation}
其中 $a_k$ 为与 $F_k$ 相关联的有限自同构群的阶。（在 \SourceCite{LLY} 中，$a_k$ 已包含在轨形积分的定义中，因此未在公式里明确出现。）下面概述这个计算，省略许多细节。

\SourcePage{364}
固定一般稳定映射 $f\in F_r\times F_{d-r}$ 的自同构，只能来自 $f|_{C_1}\in F_r$ 和 $f|_{C_2}\in F_{d-r}$ 的自同构，所以局部化公式所需的自同构群的阶为 $a_ra_{d-r}$。这解释了式~\eqref{eq:geometric-euler-data-product-localization} 分母中的因子 $a_r$、$a_{d-r}$。

还需要 $F_r\times F_{d-r}$ 在 $M_d$ 中的法丛 $N(F_r\times F_{d-r})$ 的等变 Euler 类。这个计算是论证中技术性最强的部分，因为有若干项贡献于 $\operatorname{Euler}_G(N(F_r\times F_{d-r}))$。有些地方只作简略说明，更完整的证明见 \SourceCite{LLY}。为研究此法丛，令 $M^i(k)\subset\overline M_{0,1}(\mathbb P^n,k)$ 为将标记点映到 $q_i\in\mathbb P^n$ 的稳定映射集。注意，此嵌入的法丛是 $T_{q_i}\mathbb P^n$ 经由 $e_1$ 的拉回。考虑以下图：
\[
\begin{tikzcd}[column sep=1.5em,row sep=1.2em]
F_r\times F_{d-r}\arrow[r]&M^i(r)\times M^i(d-r)\arrow[r]\arrow[d]&
\overline M_{0,1}(\mathbb P^n,r)\times\overline M_{0,1}(\mathbb P^n,d-r)\\
&\overline M(r,d-r)\arrow[d,"\widetilde\psi"]&\\
&M_d&
\end{tikzcd}
\]
其中 $\widetilde\psi$ 来自式~\eqref{eq:euler-data-graph-gluing-map}。式~\eqref{eq:euler-data-gluing-fiber-product} 的定义表明，$e_1^{-1}(q_i)=M^i(r)\times M^i(d-r)\subset\overline M(r,d-r)$，其中 $e_1$ 是图~\eqref{eq:euler-data-gluing-map} 的赋值映射。因此，此嵌入的法丛也是 $T_{q_i}\mathbb P^n$ 的拉回。

计算 $F_r\times F_{d-r}\subset M_d$ 的法丛时，上图意味着，在 $F_r\times F_{d-r}$ 的等变 $K$ 群中，有
\[
N(F_r\times F_{d-r})=N(F_r)+N(F_{d-r})-2T_{q_i}\mathbb P^n
+T_{q_i}\mathbb P^n+N(\widetilde\psi).
\]
为清楚起见，此式省去了若干拉回。取等变 Euler 类，得到（仍省去拉回）
\[
\operatorname{Euler}_G(N(F_r\times F_{d-r}))
=\frac{\operatorname{Euler}_G(N(F_r))\operatorname{Euler}_G(N(F_{d-r}))}
{\prod_{j\ne i}(\lambda_i-\lambda_j)}\operatorname{Euler}_G(N(\widetilde\psi)),
\]
因为由式~\eqref{eq:projective-fixed-euler}，$\operatorname{Euler}_G(T_{q_i}\mathbb P^n)=\prod_{j\ne i}(\lambda_i-\lambda_j)$。这出现在局部化定理的分母中，因而又解释了式~\eqref{eq:geometric-euler-data-product-localization} 中的三个因子。

还须考虑映射 $\widetilde\psi$ 的法丛。先考虑引理~\ref{lem:node-smoothing-tangent-tensor} 研究过的\IndexLocation{indexloc:0940:364}结点形变。对于 $p_1=C_0\cap C_1$，引理要求考察 $T_{p_1}C_0\otimes T_{p_1}C_1$。由于 $C_0\simeq\mathbb P^1$ 为等变同构，而 $\mathcal L_{r,1}$ 由 $C_1$ 的余切丛定义，权重为 $\hbar$ 和 $-c_1^G(\mathcal L_{r,1})$。这解释了式~\eqref{eq:geometric-euler-data-product-localization} 分母中的 $\hbar-c_1^G(\mathcal L_{r,1})$；类似地，形变 $p_2=C_0\cap C_2$ 导出分母中的 $-\hbar-c_1^G(\mathcal L_{d-r,1})$。\footnote{这与引理~\ref{lem:graph-space-descendant-factorization} 有关，并解释了式~\eqref{eq:giventals-flat-section-symbolic} 等公式中的 $\hbar-c$。}最后研究 $f|_{C_0}$ 的形变。知道 $C_0=\mathbb P^1$，而 $\overline M(r,d-r)$ 和 $\widetilde\psi$ 的定义表明，此类映射分别将 $0$、$\infty$ 映到 $(0,q)$、$(\infty,q)$，其中 $q\in\mathbb P^n$。允许 $(0,q)$、$(\infty,q)$ 的第一坐标变化，就可形变这个映射。这些法向形变可识别为 $\mathbb P^1$ 的无穷小自同构模去保持 $0$、$\infty$ 的那些自同构。用 $\mathbb P^1$ 的标准坐标 $(z_0,z_1)$，它们由向量场 $z_0\partial/\partial z_1$、$z_1\partial/\partial z_0$ 给出，产生因子 $-\hbar^2$。由于此因子出现在局部化定理的分母中，便得到式~\eqref{eq:geometric-euler-data-product-localization} 中的因子 $-\hbar^{-2}$。

应用局部化公式还需要最后一个要素：确定 $\varphi^*(\phi_{p_{i,r}})\chi_d^{\mathcal V}=\varphi^*(\phi_{p_{i,r}})\pi^*(\operatorname{Euler}_G(\mathcal V_d))$ 在 $F_r\times F_{d-r}$ 上的限制。由于 $\varphi(F_r\times F_{d-r})=p_{i,r}$，$\varphi^*(\phi_{p_{i,r}})$ 的限制为 $i_{p_{i,r}}^*(\phi_{p_{i,r}})$，这就是式~\eqref{eq:geometric-euler-data-product-localization} 的一个因子。至于
\SourcePage{365}
$\pi^*(\operatorname{Euler}_G(\mathcal V_d))$ 的限制，粘合引理（引理~\ref{lem:euler-data-gluing}）意味着，在 $F_r\times F_{d-r}$ 上有恒等式
\[
\operatorname{Euler}_G(\mathcal V)\operatorname{Euler}_G(\mathcal V_d)
=\operatorname{Euler}_G(\mathcal V_r)\operatorname{Euler}_G(\mathcal V_{d-r}),
\]
为方便起见，省去了拉回。由于 $\operatorname{Euler}_G(\mathcal V)=\ell p$，可知 $\chi_d^{\mathcal V}$ 在 $F_r\times F_{d-r}$ 上的限制为
\[
(\ell p)^{-1}\pi_1^*\operatorname{Euler}_G(\mathcal V_r)
\pi_1^*\operatorname{Euler}_G(\mathcal V_{d-r}).
\]
这些因子出现在局部化公式的分子中，从而完成对式~\eqref{eq:geometric-euler-data-product-localization} 的解释。

对 $r=0$ 作类似分析，得到公式
\begin{equation}\label{eq:geometric-euler-data-endpoint-localization}
\widehat Q_d(\lambda_i)=-\hbar^{-1}i_{p_{i,0}}^*(\phi_{p_{i,0}})
\sum_{F_d}\int_{(F_d)_G}
\frac{\pi_1^*\operatorname{Euler}_G(\mathcal V_d)}
{a_d\operatorname{Euler}_G(N(F_d))(-\hbar-c_1^G(\mathcal L_{d,1}))}.
\end{equation}
现在可证明定理。在上式中以 $r$ 替换 $d$，并应用式~\eqref{eq:euler-data-involution}，得到
\[
\overline{\widehat Q_r(\lambda_i)}
=\hbar^{-1}\overline{i_{p_{i,0}}^*(\phi_{p_{i,0}})}
\sum_{F_r}\int_{(F_r)_G}
\frac{\pi_1^*\operatorname{Euler}_G(\mathcal V_r)}
{a_r\operatorname{Euler}_G(N(F_r))(\hbar-c_1^G(\mathcal L_{r,1}))},
\]
再在式~\eqref{eq:geometric-euler-data-endpoint-localization} 中以 $d-r$ 替换 $d$，得到
\[
\widehat Q_{d-r}(\lambda_i)
=-\hbar^{-1}i_{p_{i,0}}^*(\phi_{p_{i,0}})
\sum_{F_{d-r}}\int_{(F_{d-r})_G}
\frac{\pi_1^*\operatorname{Euler}_G(\mathcal V_{d-r})}
{a_{d-r}\operatorname{Euler}_G(N(F_{d-r}))(-\hbar-c_1^G(\mathcal L_{d-r,1}))}.
\]
将两式相乘，与式~\eqref{eq:geometric-euler-data-product-localization} 乘以 $\ell\lambda_i$ 后比较，所需等式~\eqref{eq:geometric-euler-data-gluing-target} 便归结为恒等式
\[
i_{p_{i,r}}^*(\phi_{p_{i,r}})\prod_{j\ne i}(\lambda_i-\lambda_j)
=\overline{i_{p_{i,0}}^*(\phi_{p_{i,0}})}\,i_{p_{i,0}}^*(\phi_{p_{i,0}}).
\]
回顾式~\eqref{eq:linear-model-fixed-class-restriction}，$i_{p_{i,r}}^*(\phi_{p_{i,r}})$ 正是式~\eqref{eq:linear-model-normal-euler} 所计算的 $\operatorname{Euler}_G(N_{i,r})$。因此，上式归结为
\[
\begin{aligned}
&\prod_{(j,s)\ne(i,r)}\bigl(\lambda_i-\lambda_j+(r-s)\hbar\bigr)
\prod_{j\ne i}(\lambda_i-\lambda_j)\\
&\qquad=\prod_{(j,s)\ne(i,0)}\bigl(\lambda_i-\lambda_j-s\hbar\bigr)
\prod_{(j,s)\ne(i,0)}\bigl(\lambda_i-\lambda_j+s\hbar\bigr).
\end{aligned}
\]
其中左端的 $s$ 从 $0$ 到 $d$；右端第一个乘积的 $s$ 从 $0$ 到 $d-r$，第二个乘积的 $s$ 从 $0$ 到 $r$；所有乘积中 $j$ 均从 $0$ 到 $n$。此恒等式容易验证，从而完成 $r\ne0,d$ 时的证明。

由于 $r=0$ 时 Euler 恒等式显然成立（见定义~\ref{def:omega-euler-data} 后的讨论），只须再考虑 $r=d$。此时可为 $\widehat Q_d(\lambda_i+d\hbar)$ 导出类似于式~\eqref{eq:geometric-euler-data-endpoint-localization} 中 $\widehat Q_d(\lambda_i)$ 的公式，比较这些公式便得到所需恒等式。细节见 \SourceCite{LLY}。
\end{proof}

\subsection{相关联的 Euler 数据}\label{subsec:linked-euler-data}
现在转向一个重要概念。令 $\mathcal S$ 为序列 $P=\{P_d\}_{d=1}^{\infty}$ 组成的集合，其中对所有 $d$，有 $P_d\in\mathcal R_GH_G^*(N_d)$；注意 $i_{p_{i,0}}^*$ 将 $\mathcal R_GH_G^*(N_d)$ 映到 $\mathcal R_G$。
\begin{definition}\label{def:linked-euler-data}\index{Euler 数据!关联关系|hyperpage}
若对所有 $i\ne j$、$d>0$，$i_{p_{i,0}}^*(P_d-Q_d)\in\mathcal R_G$ 在 $\hbar=(\lambda_i-\lambda_j)/d$ 处有定义且为零，则称序列 $P,Q\in\mathcal S$ \emph{相关联}。
\end{definition}
\SourcePage{366}
为使此定义精确，需要解释 $\hbar=(\lambda_i-\lambda_j)/d$ 的含义。定义子群 $G'\subset G$ 为
\[
G'=\{(t,t_0,\ldots,t_n):t^d=t_i/t_j\}.
\]
包含 $G'\hookrightarrow G$ 诱导映射 $BG'\to BG$，从而在上同调上诱导拉回 $\operatorname{res}_{ij}:H^*(BG)\to H^*(BG')$。由\hyperref[subsec:equivariant-cohomology]{第~\ref*{subsec:equivariant-cohomology} 小节}的讨论，容易看出 $\operatorname{res}_{ij}$ 来自特征标的限制。这意味着 $\operatorname{res}_{ij}(\hbar)=\operatorname{res}_{ij}((\lambda_i-\lambda_j)/d)$；有时简写为 $\hbar=(\lambda_i-\lambda_j)/d$，并理解为这些符号位于 $G'$ 等变上同调中。类似地，若空间 $X$ 上有 $G$ 作用，则有映射 $\operatorname{res}_{ij}:H_G^*(X)\to H_{G'}^*(X)$。于是得到复合
\begin{equation}\label{eq:linked-data-restriction-first}
H_G^*(N_d)\xrightarrow{i_{p_{i,0}}^*}H_G^*(p_{i,0})
\xrightarrow{\operatorname{res}_{ij}}H_{G'}^*(p_{i,0}).
\end{equation}
定义~\ref{def:linked-euler-data} 的条件意味着，在所有相关分母有定义的局部化上，对所有 $i\ne j$、$d$，有 $\operatorname{res}_{ij}\circ i_{p_{i,0}}^*(P_d-Q_d)=0$。为后续使用，注意上述复合等于
\begin{equation}\label{eq:linked-data-restriction-second}
H_G^*(N_d)\xrightarrow{\operatorname{res}_{ij}}H_{G'}^*(N_d)
\xrightarrow{i_{p_{i,0}}^*}H_{G'}^*(p_{i,0}),
\end{equation}
所以，相关联也意味着，在这样的局部化上，对所有 $i\ne j$、$d$，有 $i_{p_{i,0}}^*\circ\operatorname{res}_{ij}(P_d-Q_d)=0$。

子群 $G'$ 值得研究，因为它固定了直线 $\overline{q_iq_j}\subset\mathbb P^n$ 的多重覆盖所组成的空间。考虑 $N_d$ 中参数化给出的直线
\[
L_{i,j}=\{(0,\ldots,0,aw_0^d,0,\ldots,0,bw_1^d,0,\ldots,0)\},
\]
其中 $aw_0^d$ 在第 $j$ 个位置，$bw_1^d$ 在第 $i$ 个位置。注意，$p_{i,0}\in L_{i,j}$ 对应于 $a=0$。$L_{i,j}$ 的一般点，是直线 $\overline{q_iq_j}\subset\mathbb P^n$ 的一个次数 $d$ 的多重覆盖在 $\varphi$ 下的像。子群 $G'$ 逐点固定 $L_{i,j}$，所以 $H_{G'}^*(L_{i,j})=H^*(L_{i,j})\otimes H^*(BG')$。因此，若 $u\in L_{i,j}$ 为任意点，则\IndexLocation{indexloc:0249:366}限制映射 $H_{G'}^*(L_{i,j})\to H_{G'}^*(u)$ 不依赖于 $u$。于是图
\[
\begin{tikzcd}[column sep=large,row sep=large]
H_G^*(N_d)\arrow[r,"\operatorname{res}_{ij}"]&H_{G'}^*(N_d)\arrow[r]\arrow[dr,"i_u^*"']&H_{G'}^*(L_{i,j})\arrow[d]\\
&&H_{G'}^*(u)
\end{tikzcd}
\]
中的映射也不依赖于 $u$。由于式~\eqref{eq:linked-data-restriction-second} 是 $u=p_{i,0}$ 时的特殊情形，可知在其后讨论的“限制”映射 $i_{p_{i,0}}^*\circ\operatorname{res}_{ij}$ 中，可用任意 $u\in L_{i,j}$ 替换 $p_{i,0}$。因此，$P$ 与 $Q$ 相关联，当且仅当它们在 $\mathbb P^n$ 的坐标直线的多重覆盖上相等。

现在陈述一个关键结果。
\begin{theorem}\label{thm:linked-geometric-hypergeometric-data}
在 $\hbar=(\lambda_i-\lambda_j)/d$、$i\ne j$ 时，有
\[
i_{p_{i,0}}^*(\widehat Q_d)
=\prod_{m=0}^{\ell d}\left(\ell\lambda_i-m\frac{\lambda_i-\lambda_j}{d}\right).
\]
特别地，Euler 数据 $\widehat P$ 与 $\widehat Q$ 相关联。
\end{theorem}
\SourcePage{367}
\begin{proof}
先回顾 $\widehat P_d=\prod_{m=0}^{\ell d}(\ell\kappa-m\hbar)$。由式~\eqref{eq:linked-data-restriction-first}，可用复合 $\operatorname{res}_{ij}\circ i_{p_{i,0}}^*$ 理解 $\hbar=(\lambda_i-\lambda_j)/d$。于是 $\operatorname{res}_{ij}\circ i_{p_{i,0}}^*(\hbar)=(\lambda_i-\lambda_j)/d$，而第~\ref{subsec:stable-moduli-linear-models} 小节给出 $i_{p_{i,0}}^*(\kappa)=\lambda_i$。因此
\[
\operatorname{res}_{ij}\circ i_{p_{i,0}}^*(\widehat P_d)
=\prod_{m=0}^{\ell d}\left(\ell\lambda_i-m\frac{\lambda_i-\lambda_j}{d}\right).
\]
所以，计算出 $\operatorname{res}_{ij}\circ i_{p_{i,0}}^*(\widehat Q_d)$ 后，即可得到定理的最后一个断言。

上面已证明，$\widehat Q_d$ 的限制也可由任意 $u\in L_{i,j}$ 对应的复合 $i_u^*\circ\operatorname{res}_{ij}(\widehat Q_d)$ 表示。特别地，可选
\[
u=(0,\ldots,0,w_0^d,0,\ldots,0,w_1^d,0,\ldots,0)\in L_{i,j}\subset N_d,
\]
其中 $w_0^d$、$w_1^d$ 分别位于第 $j$、第 $i$ 个位置。

证明的策略是：由于 $\widehat Q_d=\varphi_!(\chi_d^{\mathcal V})$，可在 $M_d$ 上进行计算。对于上述 $u$，考虑映射 $f:\mathbb P^1\to\mathbb P^1\times\mathbb P^n$，定义为
\[
f(w_0,w_1)=\bigl((w_1,w_0),
(0,\ldots,0,w_0^d,0,\ldots,0,w_1^d,0,\ldots,0)\bigr).
\]
这是 $M_d$ 的一个元素，显然 $\varphi(f)=u$。关键事实是，映射 $\varphi:M_d\to N_d$ 在这些点的邻域中为局部同构。考虑交换图
\[
\begin{tikzcd}[column sep=large,row sep=large]
H_G^*(M_d)\arrow[r,"\operatorname{res}_{ij}"]\arrow[d,"\varphi_!"']&H_{G'}^*(M_d)\arrow[r,"i_f^*"]\arrow[d,"\varphi_!"']&H_{G'}^*(f)\arrow[d,"\varphi_!"']\\
H_G^*(N_d)\arrow[r,"\operatorname{res}_{ij}"']&H_{G'}^*(N_d)\arrow[r,"i_u^*"']&H_{G'}^*(u)
\end{tikzcd}
\]
右端的竖直映射为同构，所以只须计算 $i_f^*\circ\operatorname{res}_{ij}(\chi_d^{\mathcal V})$。

回顾 $\chi_d^{\mathcal V}$ 为 $\pi^*\mathcal V_d$ 的等变 Euler 类。因此，在 $f$ 上须计算 $G'$ 对 $H^0(\mathbb P^1,f^*\mathcal O_{\mathbb P^n}(\ell))$ 作用的权重。一组截面基为 $w_0^mw_1^{\ell d-m}$（$0\le m\le\ell d$）。由于 $f$ 的次数为 $d$，$w_0$ 的权重为 $\lambda_j/d$，$w_1$ 的权重为 $\lambda_i/d$。这类似于\hyperref[chap:localization]{第九章}式~\eqref{eq:multiple-cover-cech-basis} 之后的讨论。于是截面 $w_0^mw_1^{\ell d-m}$ 的权重为
\[
m\frac{\lambda_j}d+(\ell d-m)\frac{\lambda_i}d
=\ell\lambda_i-m\frac{\lambda_i-\lambda_j}d.
\]
将这些权重相乘，就得到 $\hbar=(\lambda_i-\lambda_j)/d$ 时 $i_{p_{i,0}}^*(\widehat Q_d)$ 的所需公式。
\end{proof}

还需要一个唯一性结果。
\begin{theorem}\label{thm:linked-euler-data-uniqueness}
设 $P,Q$ 为相关联的 $\Omega$--Euler 数据，且对所有 $0\le i\le n$、$d\ge1$，有
\[
\deg_\hbar i_{p_{i,0}}^*(P_d-Q_d)\le(n+1)d-2.
\]
则 $P=Q$。
\end{theorem}
\begin{proof}
只简述证明。完整论证需要\IndexLocation{indexloc:1207:367}互易引理 \SourceCite[引理 2.4]{LLY}，该引理研究作为 Euler 数据所带来的结果。证明从以下观察开始：$\{\kappa^s\}=\{1,\kappa,\kappa^2,\ldots\}$ 生成
\SourcePage{368}
$H_G^*(N_d)$ 这个 $H^*(BG)$ 模。由于交配对非退化，只须证明，对所有 $s\ge0$，有
\begin{equation}\label{eq:euler-data-uniqueness-integral}
\int_{(N_d)_G}\kappa^s(P_d-Q_d)=0.
\end{equation}
以 $L_s$ 表示此式左端的积分。Euler 数据的定义意味着 $P_d$、$Q_d$ 关于 $\hbar$ 为多项式，故 $L_s$ 也是 $\hbar$ 的多项式。这一点稍后很重要。

对 $d$ 归纳，假定对 $0\le j\le d-1$，有 $P_j=Q_j$（回顾 $P_0=Q_0=\Omega$）。用局部化（推论~\ref{cor:localization-integration}）计算 $L_s$，得到对不动点 $p_{i,r}$ 的求和。由于 $P$ 为 Euler 数据，当 $0<r<d$ 时，可用 $P_1,\ldots,P_{d-1}$ 表示 $i_{p_{i,r}}^*(P_d)$；对 $Q$ 也一样。由归纳假设，仅 $r=0,d$ 的项可能贡献于式~\eqref{eq:euler-data-uniqueness-integral} 的积分。利用 \SourceCite[引理 2.4]{LLY}，可证明 $L_s$ 的局部化公式归结为
\begin{equation}\label{eq:euler-data-uniqueness-localized-sum}
L_s=\sum_{i=0}^n\left(\frac{\lambda_i^sA_i(\hbar)}{\hbar^d}
+\frac{(\lambda_i+d\hbar)^sA_i(-\hbar)}{(-\hbar)^d}\right),
\end{equation}
其中
\begin{equation}\label{eq:euler-data-uniqueness-a-polynomial}
A_i(\hbar)=\frac{(-1)^d}{d!\prod_{j\ne i}(\lambda_i-\lambda_j)}
\frac{i_{p_{i,0}}^*(P_d-Q_d)}
{\prod_{j\ne i}\prod_{s=1}^d(\lambda_i-(\lambda_j+s\hbar))}.
\end{equation}
分母部分由以下公式解释：
\[
\operatorname{Euler}_G(N_{p_{i,0}})=i_{p_{i,0}}^*(\phi_{p_{i,0}})
=\prod_{(j,s)\ne(i,0)}(\lambda_i-(\lambda_j+s\hbar)).
\]
式~\eqref{eq:euler-data-uniqueness-localized-sum} 和~\eqref{eq:euler-data-uniqueness-a-polynomial} 的分子、分母均含 $\hbar$。但 $P,Q$ 相关联，所以对所有 $j\ne i$，$i_{p_{i,0}}^*(P_d-Q_d)$ 在 $\hbar=(\lambda_i-\lambda_j)/d$ 处为零。分母中的因子 $\lambda_i-(\lambda_j+d\hbar)$ 也在此处为零，故与分子中的一个因子约去。此外，利用 \SourceCite[引理 2.4]{LLY} 和归纳假设，还可证明式~\eqref{eq:euler-data-uniqueness-a-polynomial} 分母中其他含 $\hbar$ 的因子也都约去。因此，$A_i$ 是 $\hbar$ 的多项式。于是式~\eqref{eq:euler-data-uniqueness-a-polynomial} 及定理中的次数界给出
\[
\deg_\hbar A_i\le(n+1)d-2-nd=d-2.
\]
由于对所有 $s$，$L_s$ 均为 $\hbar$ 的多项式，式~\eqref{eq:euler-data-uniqueness-localized-sum} 意味着 $\hbar^d$ 必须整除
\begin{equation}\label{eq:euler-data-uniqueness-divisibility}
\sum_{i=0}^n\bigl(\lambda_i^sA_i(\hbar)
+(-1)^d(\lambda_i+d\hbar)^sA_i(-\hbar)\bigr).
\end{equation}
现在证明所有 $A_i=0$。否则，令 $m$ 为各 $A_i$ 中非零项的最小次数，并写成
\[
A_i=a_i\hbar^m+b_i\hbar^{m+1}+\cdots,
\qquad\text{至少有一个 }a_i\ne0.
\]
$A_i$ 的次数界意味着 $m+1<d$，故式~\eqref{eq:euler-data-uniqueness-divisibility} 中 $\hbar^m$、$\hbar^{m+1}$ 的系数必须为零。若 $d+m$ 为偶数，$\hbar^m$ 的系数消失给出 $\sum_{i=0}^n2\lambda_i^sa_i=0$，
\SourcePage{369}
由此得到所有 $a_i=0$，矛盾。若 $d+m$ 为奇数，则类似地使用 $\hbar^{m+1}$ 的系数。
\end{proof}

\subsection{超几何函数与镜像变换}\label{subsec:hypergeometric-mirror-transformations}
先引入一些取等变上同调值的形式函数。令 $\mathcal S_0$ 表示序列 $B=\{B_d\}_{d=1}^{\infty}$ 所组成的集合，其中对所有 $d$，有 $B_d\in\mathcal R_GH_G^*(\mathbb P^n)$。

现在为每个 $B\in\mathcal S_0$ 关联一个形式函数，它是超几何函数的推广。
\begin{definition}\label{def:euler-data-hypergeometric-function}\index{Euler 数据!超几何函数|hyperpage}
对 $B=\{B_d\}\in\mathcal S_0$ 和可逆的 $\Omega\in\mathcal R_TH_G^*(\mathbb P^n)$，定义
\[
HG[B](t)=e^{-pt/\hbar}\left(\Omega+\sum_{d=1}^{\infty}
\frac{B_de^{dt}}{\prod_{k=0}^n\prod_{m=1}^d(p-\lambda_k-m\hbar)}\right),
\]
其中 $p\in H_G^*(\mathbb P^n)$ 为 $\mathbb P^n$ 的等变超平面类。
\end{definition}
若非等变极限 $\lambda_i\to0$ 存在，则得到形式表达式
\[
\lim_{\lambda_i\to0}HG[B](t)
=e^{-Ht/\hbar}\left(\overline\Omega+\sum_{d=1}^{\infty}
\frac{\overline B_de^{dt}}{\prod_{k=0}^n\prod_{m=1}^d(H-m\hbar)}\right),
\]
其中 $H\in H^2(\mathbb P^n)$ 为通常的超平面类，$\overline B_d=\lim_{\lambda_i\to0}B_d$，$\overline\Omega=\lim_{\lambda_i\to0}\Omega$。这是取 $H_{\mathbb C^*}^*(\mathbb P^n)\otimes\mathbb C(\hbar)$ 中形式值的函数。由于 $\mathbb C^*$ 在 $\mathbb P^n$ 上平凡作用，得到取 $H^*(\mathbb P^n)\otimes\mathbb C(\hbar)$ 中形式值的函数，此时 $\hbar$ 仅为参数。

前面几章已见过这类形式函数，\hyperref[chap:quantum-differential-equations]{第十章}中 $X=\mathbb P^n$ 的 Givental $J$ 函数就是一个明显例子。此外，下面\hyperref[subsec:lly-quintic-mirror-proof]{第~\ref*{subsec:lly-quintic-mirror-proof} 小节}将说明，形式函数 $HG[B](t)$ 在五次三维簇镜像定理的证明中起关键作用。

接下来，要把 $HG[B](t)$ 与一直研究的 Euler 数据联系起来。给定可逆的 $\Omega\in\mathcal R_TH_G^*(\mathbb P^n)$，如定义~\ref{def:omega-euler-data}，以 $\mathcal A^\Omega$ 表示所有 $\Omega$--Euler 数据的集合。还有 $\mathcal A^\Omega\subset\mathcal S$，其中 $\mathcal S$ 为序列 $Q=\{Q_d\}_{d=1}^{\infty}$ 所组成的集合，对所有 $d\ge1$，有 $Q_d\in\mathcal R_GH_G^*(N_d)$。

$\mathcal S$ 与 $\mathcal S_0$ 有何关系并不明显：后者的整个序列都位于 $\mathcal R_GH_G^*(\mathbb P^n)$ 中，而前者各项则位于不同空间的等变上同调中。幸而，有自然方法从 $\mathcal S$ 转到 $\mathcal S_0$。对于每个 $d$，有 $G$ 等变映射
\begin{equation}\label{eq:euler-data-restriction-embedding}
I_d:\mathbb P^n\longrightarrow N_d,\qquad
(a_0,\ldots,a_n)\longmapsto(a_0w_1^d,\ldots,a_nw_1^d).
\end{equation}
注意，$I_d$ 将不动点 $q_i\in\mathbb P^n$ 映到不动点 $p_{i,0}\in N_d$。定义映射 $\mathcal I:\mathcal S\to\mathcal S_0$ 为 $\mathcal I(P)_d=I_d^*(P_d)$。

通过 $\mathcal I$，任意 $\Omega$--Euler 数据 $Q$ 给出 $\mathcal I(Q)\in\mathcal S_0$，进而确定形式超几何函数 $HG[\mathcal I(Q)](t)$。在一个熟悉的例子中具体计算。
\SourceExampleNumber{11.1.5}{1}
\begin{example}\label{ex:hypergeometric-euler-data-function}
对于 $\Omega=\ell p$，有式~\eqref{eq:hypergeometric-euler-data} 定义的 $\ell p$--Euler 数据 $\widehat P$。为计算 $\mathcal I(\widehat P)$，先注意 $I_d^*(\kappa)=p$，其中 $\kappa\in H_G^*(N_d)$、$p\in H_G^*(\mathbb P^n)$ 为等变超平面类。由于 $I_d:\mathbb P^n\to N_d$ 是射影空间的等变嵌入，这一点立即成立。

由 $\widehat P$ 的定义~\eqref{eq:hypergeometric-euler-data}，得到
\[
I_d^*(\widehat P_d)=I_d^*\left(\prod_{m=0}^{\ell d}(\ell\kappa-m\hbar)\right)
=\prod_{m=0}^{\ell d}(\ell p-m\hbar).
\]
\SourcePage{370}
由此得到形式函数
\begin{equation}\label{eq:hypergeometric-euler-data-function}
HG[\mathcal I(\widehat P)](t)
=e^{-pt/\hbar}\sum_{d=0}^{\infty}
\frac{\prod_{m=0}^{\ell d}(\ell p-m\hbar)}
{\prod_{m=1}^d\prod_{k=0}^n(p-\lambda_k-m\hbar)}e^{dt},
\end{equation}
除 $\hbar$ 的符号外，它与 \SourceCite{Givental2} 中为描述 $\mathbb P^n$ 的次数 $\ell$ 超曲面的量子上同调而研究的函数 $I_{\mathcal V}$ 一致。\hyperref[subsec:lly-quintic-mirror-proof]{第~\ref*{subsec:lly-quintic-mirror-proof} 小节}将说明，当 $\mathcal V=\mathcal O_{\mathbb P^4}(5)$ 时，非等变极限 $\lim_{\lambda_i\to0}HG[\mathcal I(\widehat P)](t)$ 给出\hyperref[subsec:hypergeometric-mirror-map]{第~\ref*{subsec:hypergeometric-mirror-map} 小节}研究的五次镜像 Picard--Fuchs 方程的解 $y_0,y_1,y_2,y_3$。
\end{example}

除 $\widehat P$ 外，另一组重要的 $\ell p$--Euler 数据是定理~\ref{thm:geometric-euler-data-gluing} 中的 $\widehat Q$。此时得到形式函数
\[
HG[\mathcal I(\widehat Q)](t)
=e^{-pt/\hbar}\left(\Omega+\sum_{d=1}^{\infty}
\frac{I_d^*(\widehat Q_d)e^{dt}}
{\prod_{k=0}^n\prod_{m=1}^d(p-\lambda_k-m\hbar)}\right).
\]
可如下计算 $I_d^*(\widehat Q_d)$。
\begin{lemma}\label{lem:geometric-euler-data-hypergeometric-bound}\index{Euler 数据!次数界|hyperpage}
令 $\phi^d=\prod_{k=0}^n\prod_{m=1}^d(p-\lambda_k-m\hbar)\in H_G^*(\mathbb P^n)$，表示 $HG[\mathcal I(\widehat Q)](t)$ 中 $e^{dt}$ 系数的分母。则：
\begin{enumerate}[label=(\roman*)]
\item 若 $\pi_1,e_1$ 如图~\eqref{eq:euler-data-universal-curve}，$\mathcal L_{d,1}$ 如定理~\ref{thm:geometric-euler-data-gluing} 的证明，则
\[
I_d^*(\widehat Q_d)=\phi^d e_{1!}\left(
\frac{\pi_1^*\operatorname{Euler}_G(\mathcal V_d)}
{\hbar(\hbar+c_1^G(\mathcal L_{d,1}))}\right).
\]
\item $\deg_\hbar I_d^*(\widehat Q_d)\le(n+1)d-2$。
\end{enumerate}
\end{lemma}
\begin{proof}
为证明第 (i) 部分，只须证明，对每个 $i=0,\ldots,n$，两端在 $q_i$ 处的局部化相等。对于左端，由 $I_d(q_i)=p_{i,0}$，有
\[
i_{q_i}^*I_d^*(\widehat Q_d)=i_{p_{i,0}}^*(\widehat Q_d)=\widehat Q_d(\lambda_i).
\]
结合式~\eqref{eq:geometric-euler-data-endpoint-localization}，得到
\[
i_{q_i}^*I_d^*(\widehat Q_d)=i_{p_{i,0}}^*(\phi_{p_{i,0}})
\sum_{F_d}\int_{(F_d)_G}
\frac{\pi_1^*\operatorname{Euler}_G(\mathcal V_d)}
{a_d\operatorname{Euler}_G(N(F_d))\hbar(\hbar+c_1^G(\mathcal L_{d,1}))}.
\]
另一方面，若如式~\eqref{eq:projective-fixed-dual-class} 令 $\phi_{q_i}=\prod_{j\ne i}(p-\lambda_j)$，则式~\eqref{eq:linear-model-normal-euler} 和 $\phi^d$ 的定义立即给出
\[
i_{p_{i,0}}^*(\phi_{p_{i,0}})=i_{q_i}^*(\phi_{q_i})i_{q_i}^*(\phi^d).
\]
因此，上式成为
\begin{equation}\label{eq:geometric-euler-data-restriction-localized}
\begin{aligned}
i_{q_i}^*I_d^*(\widehat Q_d)
&=i_{q_i}^*(\phi_{q_i})i_{q_i}^*(\phi^d)\\
&\quad\times\sum_{F_d}\int_{(F_d)_G}
\frac{\pi_1^*\operatorname{Euler}_G(\mathcal V_d)}
{a_d\operatorname{Euler}_G(N(F_d))\hbar(\hbar+c_1^G(\mathcal L_{d,1}))}.
\end{aligned}
\end{equation}
转向右端，令
\[
\alpha^d=\frac{\pi_1^*\operatorname{Euler}_G(\mathcal V_d)}
{\hbar(\hbar+c_1^G(\mathcal L_{d,1}))},
\]
\SourcePage{371}
于是须计算 $\phi^d e_{1!}(\alpha^d)$ 在 $q_i$ 处的局部化。将式~\eqref{eq:geometric-euler-data-localization-integral} 适配到 $\mathbb P^n$，得到
\[
i_{q_i}^*(\phi^d e_{1!}(\alpha^d))
=\int_{(\mathbb P^n)_G}\phi_{q_i}\phi^d e_{1!}(\alpha^d)
=\int_{\overline M_{0,1}(\mathbb P^n,d)_G}e_1^*(\phi_{q_i}\phi^d)\alpha^d.
\]
用局部化计算右端积分时，因子 $\phi_{q_i}$ 意味着只须考虑在 $e_1$ 下映到 $q_i$ 的固定点集分支。这些恰好是定理~\ref{thm:geometric-euler-data-gluing} 证明中使用的分支 $\{F_d\}$；将局部化公式应用于此积分，立即得到式~\eqref{eq:geometric-euler-data-restriction-localized} 的右端。这就证明了第 (i) 部分。

对于第 (ii) 部分，由于 $\widehat Q$ 为 Euler 数据，$\widehat Q_d$、进而 $I_d^*(\widehat Q_d)$ 均为 $\hbar$ 的多项式。不过，由 $\phi^d$ 的定义，有 $\deg_\hbar\phi^d=(n+1)d$。于是，$I_d^*(\widehat Q_d)$ 的所需次数界由第 (i) 部分得到。
\end{proof}
注意，引理第二部分所证明的 $\deg_\hbar I_d^*(\widehat Q_d)$ 的界，恰好就是\IndexLocation{indexloc:0270:371}唯一性定理（定理~\ref{thm:linked-euler-data-uniqueness}）所需的次数界。

最终目标是在五次三维簇的情形下，联系 $HG[\mathcal I(\widehat P)](t)$ 与 $HG[\mathcal I(\widehat Q)](t)$。但这两个函数并不相等，Euler 数据 $\widehat P$、$\widehat Q$ 也不相等。因此，需要一种将一组 Euler 数据变成另一组的方法。这导出以下定义。
\begin{definition}\label{def:euler-data-mirror-transformation}\index{Euler 数据!镜像变换|hyperpage}
可逆映射 $\mu:\mathcal A^\Omega\to\mathcal A^\Omega$ 若对所有 $P\in\mathcal A^\Omega$，$\mu(P)$ 与 $P$ 相关联，则称它为\emph{镜像变换}。
\end{definition}
马上将看到，第~\ref{subsec:quintic-mirror-theorem-statement} 小节讨论的镜像映射，可描述为此定义意义下的镜像变换。

目标是构造一类特殊镜像变换 $\mu$，使 $HG[\mathcal I(P)](t)$ 与 $HG[\mathcal I(\mu(P))](t)$ 容易联系。先考虑超几何函数 $HG[B](t)$ 的一些初等变换。
\begin{lemma}\label{lem:hypergeometric-formal-transformation}
给定任意 $B\in\mathcal S_0$，给定任意 $f,g\in e^t\mathcal R_T[[e^t]]$。则存在唯一 $\widetilde B\in\mathcal S_0$，使
\[
e^{f/\hbar}HG[B](t+g)=HG[\widetilde B](t).
\]
\end{lemma}
\begin{proof}
这是直接的形式幂级数计算，见 \SourceCite{LLY}。
\end{proof}
利用此引理，可描述我们关心的镜像变换类。
\begin{proposition}\label{prop:hypergeometric-mirror-transformation-lift}\index{Euler 数据!Lagrange 提升|hyperpage}
令 $f=a\hbar+b$，其中 $a,b,g\in e^t\mathcal R_T[[e^t]]$。则存在镜像变换 $\mu:\mathcal A^\Omega\to\mathcal A^\Omega$，使对所有 Euler 数据 $P\in\mathcal A^\Omega$，有
\[
e^{f/\hbar}HG[\mathcal I(P)](t+g)=HG[\mathcal I(\mu(P))](t).
\]
\end{proposition}
\begin{proof}
给定 $f,g$，引理~\ref{lem:hypergeometric-formal-transformation} 中的映射 $B\mapsto\widetilde B$ 定义变换 $\mu_0:\mathcal S_0\to\mathcal S_0$。用此记号，前一引理意味着，对 $B\in\mathcal S_0$，有
\[
e^{f/\hbar}HG[B](t+g)=HG[\mu_0(B)](t).
\]
为证明命题，只须证明 $\mu_0$ 可提升为满足 $\mathcal I(\mu(P))=\mu_0(\mathcal I(P))$ 的镜像变换 $\mu$。

\SourcePage{372}
回顾 $\mathcal I:\mathcal S\to\mathcal S_0$ 由 $\mathcal I(P)_d=I_d^*(P_d)$ 定义，$I_d$ 如式~\eqref{eq:euler-data-restriction-embedding}。为定义 $\mu$，须构造反方向的映射，做法如下。命题~\ref{prop:atiyah-bott-localization} 对 $N_d$ 所给的同构可写成
\[
\mathcal R_GH_G^*(N_d)\xrightarrow{\ \bigoplus i_{p_{i,r}}^*\ }
(\mathcal R_G)^{(N_d)^G},
\]
其中 $(N_d)^G$ 如往常一样表示 $G$ 固定点集。因此，为指定元素 $\omega\in\mathcal R_GH_G^*(N_d)$，只须指定所有 $i,r$ 对应的限制 $i_{p_{i,r}}^*(\omega)$。据此，用以下条件唯一定义 $\mathcal L:\mathcal S_0\to\mathcal S$：
\begin{equation}\label{eq:euler-data-lagrange-lift}
i_{p_{i,r}}^*(\mathcal L(B)_d)
=(i_{q_i}^*(\Omega))^{-1}\overline{i_{q_i}^*(B_r)}\,i_{q_i}^*(B_{d-r}),
\qquad0\le i\le n,\quad0\le r\le d,
\end{equation}
对任意 $B\in\mathcal S_0$ 均如此。这些等式中，$B_0=\Omega$。

映射 $\mathcal L$ 有一些良好性质。首先，$\mathcal I\circ\mathcal L=\operatorname{id}_{\mathcal S_0}$。为证明这一点，注意 $I_d(q_i)=p_{i,0}$；式~\eqref{eq:euler-data-lagrange-lift} 在 $r=0$ 时给出
\[
\begin{aligned}
i_{q_i}^*(\mathcal I(\mathcal L(B))_d)
&=i_{q_i}^*\bigl(I_d^*(\mathcal L(B)_d)\bigr)\\
&=i_{p_{i,0}}^*(\mathcal L(B)_d)\\
&=(i_{q_i}^*(\Omega))^{-1}\overline{i_{q_i}^*(B_0)}\,i_{q_i}^*(B_d)
=i_{q_i}^*(B_d),
\end{aligned}
\]
最后一个等号来自 $B_0=\Omega\in\mathcal R_TH_T^*(\mathbb P^n)$。由 $\mathcal R_TH_T^*(\mathbb P^n)$ 的局部化（并保留形式参数 $\hbar$），得到 $\mathcal I(\mathcal L(B))=B$。$\mathcal L$ 的第二个良好性质是：对任意 $B\in\mathcal S_0$，$\mathcal L(B)$ 满足定义~\ref{def:omega-euler-data} 中的 Euler 恒等式。因此，若所有 $d$ 都有 $\mathcal L(B)_d\in\mathcal R_TH_G^*(N_d)$，则 $\mathcal L(B)$ 是 Euler 数据。最后，由定义也立即可知，若 $Q\in\mathcal A^\Omega\subset\mathcal S$，则 $Q=\mathcal L\circ\mathcal I(Q)$。

既然有了 $\mathcal L:\mathcal S_0\to\mathcal S$ 和 $\mathcal I:\mathcal S\to\mathcal S_0$，映射 $\mu_0:\mathcal S_0\to\mathcal S_0$ 可延伸为 $\mu=\mathcal L\circ\mu_0\circ\mathcal I:\mathcal S\to\mathcal S$。用 \SourceCite{LLY} 的术语，$\mu$ 是 $\mu_0$ 的 \emph{Lagrange 提升}。如此定义 $\mu$ 后，$\mathcal I\circ\mu=\mu_0\circ\mathcal I$ 立即来自 $\mathcal I\circ\mathcal L=\operatorname{id}_{\mathcal S_0}$。因此，还须证明 $\mu$ 是镜像变换。给定 $P\in\mathcal A^\Omega$，$P$ 与 $\mu(P)$ 相关联这一条件，即 $i_{p_{i,0}}^*(P_d-\mu(P)_d)$ 在 $\hbar=(\lambda_i-\lambda_j)/d$ 处为零，立即来自证明引理~\ref{lem:hypergeometric-formal-transformation} 时所作的显式幂级数计算。最后，须证明 $\mu(P)$ 是 $\Omega$--Euler 数据。前一段表明，$\mu(P)$ 自动满足 Euler 条件，因此最后一步是证明 $\mu(P)_d\in\mathcal R_TH_G^*(N_d)$。证明见 \SourceCite[引理 2.15]{LLY}。
\end{proof}

现在几乎已有证明\IndexLocation{indexloc:0864:372}镜像定理所需的全部准备。不过，在开始证明之前，应解释\hyperref[subsec:lly-euler-data]{第~\ref*{subsec:lly-euler-data}}、\hyperref[subsec:linked-euler-data]{\ref*{subsec:linked-euler-data}}、\hyperref[subsec:hypergeometric-mirror-transformations]{\ref*{subsec:hypergeometric-mirror-transformations} 小节}与 \SourceCite{LLY} 的关系。一直使用的丛 $\mathcal O_{\mathbb P^n}(\ell)$ 是\emph{凸丛}\index{凸丛|hyperpage}的特殊情形；凸丛是指一个 $T$ 等变丛 $\mathcal V$，其 $\operatorname{Euler}_T(\mathcal V)$ 在 $\mathcal R_TH_T^*(\mathbb P^n)$ 中可逆，且对每个无标记亏格零稳定映射 $f:C\to\mathbb P^n$，有 $H^1(C,f^*\mathcal V)=0$。凸丛的其他例子有 $\mathcal V=\bigoplus_i\mathcal O_{\mathbb P^n}(\ell_i)$，其中 $\ell_i>0$。类似地，\SourceCite{LLY} 将满足相应条件、但要求对每个无标记亏格零稳定映射 $f$ 有 $H^0(C,f^*\mathcal V)=0$ 的丛 $\mathcal V$ 称为\emph{凹丛}\index{凹丛|hyperpage}。最后，凸丛与凹丛的直和称为\emph{凸凹丛}\index{凸凹丛|hyperpage}。

给定 $\mathbb P^n$ 上的凸凹丛 $\mathcal V$，如下定义 $\overline M_{0,0}(\mathbb P^n,d)$ 上的丛 $\mathcal V_d$。若 $\mathcal V$ 凸，则如图~\eqref{eq:euler-data-universal-curve} 后的讨论一样，取 $\mathcal V_d=\pi_{1*}e_1^*(\mathcal V)$。若 $\mathcal V$ 凹，则取 $\mathcal V_d=R^1\pi_{1*}e_1^*(\mathcal V)$。若 $\mathcal V$ 凸凹，则 $\mathcal V_d$ 为将这两种构造分别应用于凸部分、凹部分后所得丛的直和。

最后，设 $\mathcal V$ 为形如 $\bigoplus_i\mathcal O_{\mathbb P^n}(\ell_i)$（$\ell_i\ne0$）的凸凹丛。用 $\mathcal V_d$，仍如式~\eqref{eq:geometric-euler-data} 一样定义 $\widehat Q_d$。此外，还可
\SourcePage{373}
推广式~\eqref{eq:hypergeometric-euler-data}，对 $d>0$ 定义
\begin{equation}\label{eq:concavex-hypergeometric-euler-data}
\widehat P_d=
\prod_{\ell_i>0}\prod_{m=0}^{\ell_i d}(\ell_i\kappa-m\hbar)
\times\prod_{\ell_i<0}\prod_{m=1}^{-\ell_i d-1}(\ell_i\kappa+m\hbar).
\end{equation}
关键在于，对于 $\mathcal V=\mathcal O_{\mathbb P^n}(\ell)$ 已证明的 $\widehat Q$、$\widehat P$ 的结果，在 $\mathcal V=\bigoplus_i\mathcal O_{\mathbb P^n}(\ell_i)$ 时仍成立；此时取 $\Omega=\operatorname{Euler}_T(\mathcal V^+)/\operatorname{Euler}_T(\mathcal V^-)$，其中 $\mathcal V^+$、$\mathcal V^-$ 为凸、凹部分，。细节见 \SourceCite{LLY}。
% 第八批接续：原 PDF 物理页373中部。
\subsection{临界丛}\label{subsec:critical-bundles}\index{临界丛|hyperpage}
\hypertarget{source:373-batch8}{}
在证明五次三维簇 $V\subset\mathbb P^4$ 的镜像定理之前，需要解释 $V$ 的 Gromov--Witten 不变量 $K_d=\langle I_{0,0,d}\rangle$ 与此前讨论的联系。定理~\ref{thm:critical-quintic-hypergeometric-potential} 将表明，对丛 $\mathcal V=\mathcal O_{\mathbb P^4}(5)$，$K_d$ 自然出现在超几何函数 $HG[\mathcal I(\widehat Q)](t)$ 中。还将讨论凸凹丛的一类子类，即\emph{临界丛}；它们满足该定理的一个修改版本。

回顾\hyperref[subsec:quintic-mirror-theorem-statement]{第~\ref*{subsec:quintic-mirror-theorem-statement} 小节}，$V$ 的 Gromov--Witten 不变量 $K_d$ 由
\[
K_d=\int_{\overline M_{0,0}(\mathbb P^4,d)}\operatorname{Euler}(\mathcal V_d)
\]
给出。用 $K_d$ 可将五次三维簇的 Gromov--Witten 势 $\Phi$ 写成
\[
\Phi=\frac{5t^3}{6}+\sum_{d=1}^{\infty}K_de^{dt}.
\]
于是有如下结果。
\begin{theorem}\label{thm:critical-quintic-hypergeometric-potential}
令 $\mathcal V=\mathcal O_{\mathbb P^4}(5)$，并如式~\eqref{eq:geometric-euler-data} 定义 $\widehat Q_d$。则在非等变极限 $\lambda_i\to0$ 下，有
\[
\lim_{\lambda_i\to0}HG[\mathcal I(\widehat Q)](t)
=5H\left(1-t\frac H\hbar+\frac{\Phi'}5\frac{H^2}{\hbar^2}
-\frac{t\Phi'-2\Phi}5\frac{H^3}{\hbar^3}\right),
\]
其中 $H$ 为 $H^*(\mathbb P^4)$ 中的超平面类。
\end{theorem}
\begin{proof}
定义~\ref{def:euler-data-hypergeometric-function} 的非等变极限给出
\[
\lim_{\lambda_i\to0}HG[\mathcal I(\widehat Q)](t)
=e^{-Ht/\hbar}\left(5H+\sum_{d=1}^{\infty}
\frac{\lim_{\lambda_i\to0}I_d^*(\widehat Q_d)}{\prod_{m=1}^d(H-m\hbar)^5}e^{dt}\right).
\]
现在取此表达式中 $e^{dt}$ 的系数，并用 $H^*(\mathbb P^4)$ 的基 $1,H,\ldots,H^4$ 表示：
\[
e^{-Ht/\hbar}\frac{\lim_{\lambda_i\to0}I_d^*(\widehat Q_d)}{\prod_{m=1}^d(H-m\hbar)^5}
=A_0+A_1H+A_2H^2+A_3H^3+A_4H^4.
\]
将此等式乘以 $H^k$ 并在 $\mathbb P^4$ 上积分，即可求出 $A_i$。这给出
\[
A_{4-k}=\int_{\mathbb P^4}H^ke^{-Ht/\hbar}
\frac{\lim_{\lambda_i\to0}I_d^*(\widehat Q_d)}{\prod_{m=1}^d(H-m\hbar)^5}.
\]
\SourcePage{374}
以下公式如何推出本定理，留给读者验证：
\begin{equation}\label{eq:critical-quintic-integral-coefficients}
\int_{\mathbb P^4}H^ke^{-Ht/\hbar}
\frac{\lim_{\lambda_i\to0}I_d^*(\widehat Q_d)}{\prod_{m=1}^d(H-m\hbar)^5}
=\begin{cases}
\hbar^{-3}(2-dt)K_d,&k=0,\\
\hbar^{-2}dK_d,&k=1,\\
0,&k>1.
\end{cases}
\end{equation}
式~\eqref{eq:critical-quintic-integral-coefficients} 的 $k=0$ 情形在 \SourceCite{LLY} 中明确给出，而 $k\ge1$ 的情形则隐含于其中。

为证明式~\eqref{eq:critical-quintic-integral-coefficients}，使用引理~\ref{lem:geometric-euler-data-hypergeometric-bound} 给出的 $I_d^*(\widehat Q_d)$ 公式。令 $\lambda_i\to0$，取其非等变极限，得到
\[
\lim_{\lambda_i\to0}I_d^*(\widehat Q_d)
=\prod_{m=1}^d(H-m\hbar)^5\,
 e_{1!}\left(\frac{\pi_1^*\operatorname{Euler}(\mathcal V_d)}
 {\hbar(\hbar+c_1(\mathcal L_{d,1}))}\right).
\]
将其代入式~\eqref{eq:critical-quintic-integral-coefficients} 左端的积分，积分简化为
\begin{equation}\label{eq:critical-quintic-pushforward-integral}
\begin{aligned}
&\int_{\mathbb P^4}H^ke^{-Ht/\hbar}
 e_{1!}\left(\frac{\pi_1^*\operatorname{Euler}(\mathcal V_d)}
 {\hbar(\hbar+c_1(\mathcal L_{d,1}))}\right)\\
&\quad=\int_{\overline M_{0,1}(\mathbb P^4,d)}
 e_1^*H^k e^{-e_1^*Ht/\hbar}
 \frac{\pi_1^*\operatorname{Euler}(\mathcal V_d)}{\hbar(\hbar+c_1(\mathcal L_{d,1}))}\\
&\quad=\int_{\overline M_{0,0}(\mathbb P^4,d)}\operatorname{Euler}(\mathcal V_d)\,
 \pi_{1!}\left(e_1^*H^k\frac{e^{-e_1^*Ht/\hbar}}
 {\hbar(\hbar+c_1(\mathcal L_{d,1}))}\right).
\end{aligned}
\end{equation}
因为 $\operatorname{Euler}(\mathcal V_d)$ 的次数为最高次，这个积分恰为 $K_d$ 乘以被积式中第二个因子的零次分量。又因为 $\pi_1$ 的相对维数为 $1$，该分量就是
\[
\begin{cases}
\pi_{1!}\bigl(\hbar^{-3}(-e_1^*Ht-c_1(\mathcal L_{d,1}))\bigr),&k=0,\\
\pi_{1!}(\hbar^{-2}e_1^*H),&k=1,\\
0,&k>1
\end{cases}
\]
的零次值。在 $\pi_1$ 的一条纤维上积分后，它化为
\[
\begin{cases}
\hbar^{-3}(-dt+2),&k=0,\\
\hbar^{-2}d,&k=1,\\
0,&k>1.
\end{cases}
\]
这就得到式~\eqref{eq:critical-quintic-integral-coefficients}，从而证明了定理。
\end{proof}

事实上，这个定理在更一般的情形下也成立。考察上述证明可知，关键步骤之一是式~\eqref{eq:critical-quintic-pushforward-integral} 后的论证，其中用到了 $\operatorname{Euler}(\mathcal V_d)$ 在 $\overline M_{0,0}(\mathbb P^4,d)$ 上为最高次类这一事实。这导出如下定义。
\begin{definition}\label{def:critical-bundle}
若 $\mathbb P^n$ 上的凸凹丛 $\mathcal V$ 所诱导的丛 $\mathcal V_d$ 在 $\overline M_{0,0}(\mathbb P^n,d)$ 上，对所有 $d>0$ 都满足
\[
\operatorname{rank}\mathcal V_d=(n+1)d+n-3=\dim\overline M_{0,0}(\mathbb P^n,d),
\]
则称 $\mathcal V$ 为\emph{临界丛}。
\end{definition}
\SourcePage{375}
当 $\mathcal V$ 为临界丛时，$\operatorname{Euler}(\mathcal V_d)$ 是 $\overline M_{0,0}(\mathbb P^n,d)$ 上的最高次类，因而可以定义
\[
K_d=\int_{\overline M_{0,0}(\mathbb P^n,d)}\operatorname{Euler}(\mathcal V_d),
\]
推广五次三维簇情形的构造。在此情形下，由于 $\mathcal V$ 凸凹，引理~\ref{lem:geometric-euler-data-hypergeometric-bound} 仍适用；由于 $\mathcal V$ 临界，定理~\ref{thm:critical-quintic-hypergeometric-potential} 的论证也适用。因此，已证明
\begin{equation}\label{eq:critical-bundle-integral-coefficients}
\int_{\mathbb P^n}H^ke^{-Ht/\hbar}
\frac{\lim_{\lambda_i\to0}I_d^*(\widehat Q_d)}{\prod_{m=1}^d(H-m\hbar)^{n+1}}
=\begin{cases}
\hbar^{-3}(2-dt)K_d,&k=0,\\
\hbar^{-2}dK_d,&k=1,\\
0,&k>1.
\end{cases}
\end{equation}
细节见 \SourceCite{LLY}。下面将看到，这有一些有趣的应用。

\hyperref[subsec:lly-euler-data]{第~\ref*{subsec:lly-euler-data} 小节}考虑了形如 $\mathcal V=\mathcal O_{\mathbb P^n}(\ell)$ 的丛。读者容易验证，此类型中唯一的临界丛是 $\mathcal V=\mathcal O_{\mathbb P^4}(5)$。不过，对于 $\mathcal V=\bigoplus_i\mathcal O_{\mathbb P^n}(\ell_i)$，还有其他值得研究的临界丛。以下例子与我们的讨论尤其相关。
\SourceExampleNumber{11.1.6}{1}
\begin{example}\label{ex:critical-multiple-cover-bundle}
考虑 $\mathbb P^1$ 上的凹丛
\[
\mathcal V=\mathcal O_{\mathbb P^1}(-1)\oplus\mathcal O_{\mathbb P^1}(-1).
\]
这里，$\mathcal V_d$ 定义为
\[
\mathcal V_d=R^1\pi_{1*}e_1^*\bigl(\mathcal O_{\mathbb P^1}(-1)\oplus\mathcal O_{\mathbb P^1}(-1)\bigr).
\]
对代表 $\overline M_{0,0}(\mathbb P^1,d)$ 中一点的映射 $f:C\to\mathbb P^1$，容易验证
\[
\begin{aligned}
\operatorname{rank}\mathcal V_d
&=\dim H^1\bigl(C,f^*\mathcal O_{\mathbb P^1}(-1)\oplus f^*\mathcal O_{\mathbb P^1}(-1)\bigr)\\
&=2d-2=\dim\overline M_{0,0}(\mathbb P^1,d).
\end{aligned}
\]
因此 $\mathcal V$ 为临界丛。于是式~\eqref{eq:critical-bundle-integral-coefficients} 表明，$HG[\mathcal I(\widehat Q)](t)$ 的非等变极限涉及
\[
K_d=\int_{\overline M_{0,0}(\mathbb P^1,d)}
\operatorname{Euler}\!\left(R^1\pi_{1*}e_1^*
\bigl(\mathcal O_{\mathbb P^1}(-1)\oplus\mathcal O_{\mathbb P^1}(-1)\bigr)\right).
\]
这恰是定理~\ref{thm:multiple-cover-localization} 计算的积分；该定理计算了 Calabi--Yau 三维簇中刚性嵌入的光滑有理曲线的\IndexLocation{indexloc:0916:375}多重覆盖贡献。即将应用 \SourceCite{LLY} 的方法，给出该定理的第二个证明。
\end{example}
还应指出，\SourceCite{LLY} 给出了 $\mathbb P^n$ 上所有由线丛直和构成的临界丛的完整清单。
\subsection{定理~\ref{thm:lly-quintic-mirror} 的证明}\label{subsec:lly-quintic-mirror-proof}
现在可以证明\IndexLocation{indexloc:0867:375}五次三维簇的镜像定理了。在本节余下部分，均假设 $\mathcal V=\mathcal O_{\mathbb P^4}(5)$，并使用\hyperref[subsec:lly-euler-data]{第~\ref*{subsec:lly-euler-data} 小节}定义的 $\widehat P$、$\widehat Q$。证明的关键是构造两个\IndexLocation{indexloc:0268:375}镜像变换 $\mu$、$\nu$，使 $\mu(\widehat P)=\nu(\widehat Q)$。这将导出超几何函数 $HG[\mathcal I(\widehat P)](t)$ 与 $HG[\mathcal I(\widehat Q)](t)$ 之间的关系；取非等变极限 $\lambda_i\to0$，便得到所需结果。
\SourcePage{376}
为构造所需的镜像变换，首先更仔细地考察 $HG[\mathcal I(\widehat P)](t)$。取式~\eqref{eq:hypergeometric-euler-data-function} 在 $\lambda_i\to0$ 时的非等变极限，得到
\begin{equation}\label{eq:quintic-hypergeometric-period-expansion}
\begin{aligned}
\lim_{\lambda_i\to0}HG[\mathcal I(\widehat P)](t)
&=e^{-Ht/\hbar}\sum_{d=0}^{\infty}
\frac{\prod_{m=0}^{5d}(5H-m\hbar)}{\prod_{m=1}^d(H-m\hbar)^5}e^{dt}\\
&=5H\left(y_0-y_1\frac H\hbar+y_2\frac{H^2}{\hbar^2}-y_3\frac{H^3}{\hbar^3}\right).
\end{aligned}
\end{equation}
这些就是\hyperref[subsec:quintic-mirror-theorem-statement]{第~\ref*{subsec:quintic-mirror-theorem-statement} 小节}讨论的函数 $y_0,y_1,y_2,y_3$；它们出现在镜像映射 $\Psi(t)=y_1/y_0$ 中，也出现在定理~\ref{thm:lly-quintic-mirror} 的陈述中。

现在两次使用命题~\ref{prop:hypergeometric-mirror-transformation-lift}，构造镜像变换 $\mu$、$\nu$。对于 $\mu$，取
\begin{equation}\label{eq:quintic-first-mirror-transformation-parameters}
\begin{gathered}
f_1=-(\log y_0)\hbar+\frac{h_2}{y_0}\sum_{k=0}^4\lambda_k,\\
h_2=\sum_{d=1}^{\infty}\frac{(5d)!}{(d!)^5}
\left(\sum_{m=1}^d\frac1m\right)e^{dt},
\end{gathered}
\end{equation}
并取 $g_1=0$，随后命题~\ref{prop:hypergeometric-mirror-transformation-lift} 给出镜像变换 $\mu$，使
\begin{equation}\label{eq:quintic-first-mirror-transformation}
e^{f_1/\hbar}HG[\mathcal I(\widehat P)](t)
=HG[\mathcal I(\mu(\widehat P))](t).
\end{equation}
对于第二个镜像变换，取 $f_2=0$、$g_2=\Psi(t)-t\in e^t\mathcal R[[e^t]]$，其中 $\Psi(t)=y_1/y_0$ 是镜像映射。再次应用上述命题，得到镜像变换 $\nu$，使
\begin{equation}\label{eq:quintic-second-mirror-transformation}
HG[\mathcal I(\nu(\widehat Q))](t)
=HG[\mathcal I(\widehat Q)](\Psi(t)).
\end{equation}
镜像变换所得数据 $\mu(\widehat P)$ 与 $\nu(\widehat Q)$ 有如下关系。
\begin{theorem}\label{thm:quintic-transformed-euler-data-equality}
$\mu(\widehat P)=\nu(\widehat Q)$。
\end{theorem}
\begin{proof}
遵循 \SourceCite{LLY} 的证明。由定理~\ref{thm:linked-geometric-hypergeometric-data}，$\widehat P$、$\widehat Q$ 是相关联的 Euler 数据。此外，因为 $\mu$、$\nu$ 是镜像变换，$\widehat P,\mu(\widehat P)$ 及 $\widehat Q,\nu(\widehat Q)$ 也分别相关联。由传递性，$\mu(\widehat P)$、$\nu(\widehat Q)$ 是相关联的 Euler 数据。因此，只要证明次数界
\[
\deg_{\hbar}i_{p_{i,0}}^*\bigl(\mu(\widehat P)_d-\nu(\widehat Q)_d\bigr)\le5d-2,
\]
本定理就由定理~\ref{thm:linked-euler-data-uniqueness} 得到。由于 $i_{p_{i,0}}=I_d\circ i_{q_i}$，只须证明
\[
\frac{I_d^*(\mu(\widehat P)_d-\nu(\widehat Q)_d)}
{\prod_{k=0}^4\prod_{m=1}^d(p-\lambda_k-m\hbar)}
=0\pmod{\hbar^{-2}}.
\]
不过，由定义~\ref{def:euler-data-hypergeometric-function}，这等价于
\[
HG[\mathcal I(\nu(\widehat Q))](t)
=HG[\mathcal I(\mu(\widehat P))](t)\pmod{\hbar^{-2}}.
\]
再用式~\eqref{eq:quintic-first-mirror-transformation}、\eqref{eq:quintic-second-mirror-transformation}，问题归结为证明
\begin{equation}\label{eq:quintic-transformed-hypergeometric-congruence}
e^{f_1/\hbar}HG[\mathcal I(\widehat P)](t)
=HG[\mathcal I(\widehat Q)](\Psi(t))\pmod{\hbar^{-2}}.
\end{equation}
为此，首先按 $\hbar^{-1}$ 的幂展开 $HG[\mathcal I(\widehat P)](t)$，得到
\[
HG[\mathcal I(\widehat P)](t)
=5p\left(y_0-\hbar^{-1}\left(py_1+h_2\sum_{k=0}^4\lambda_k\right)+O(\hbar^{-2})\right).
\]
于是，由式~\eqref{eq:quintic-first-mirror-transformation-parameters} 及一些代数运算可得
\[
e^{f_1/\hbar}HG[\mathcal I(\widehat P)](t)
=5p-\hbar^{-1}5p^2\frac{y_1}{y_0}+O(\hbar^{-2}).
\]
\SourcePage{377}
再看式~\eqref{eq:quintic-transformed-hypergeometric-congruence} 的右端。首先，由引理~\ref{lem:geometric-euler-data-hypergeometric-bound} 的第 (ii) 部分，有 $\deg_{\hbar}I_d^*(\widehat Q_d)\le5d-2$。于是定义~\ref{def:euler-data-hypergeometric-function} 给出
\[
\begin{aligned}
HG[\mathcal I(\widehat Q)](t)
&=e^{-pt/\hbar}\bigl(5p+O(\hbar^{-2})\bigr)\\
&=5p-\hbar^{-1}5p^2t+O(\hbar^{-2}).
\end{aligned}
\]
因为 $\Psi=y_1/y_0$，有
\[
HG[\mathcal I(\widehat Q)](\Psi(t))
=5p-\hbar^{-1}5p^2\frac{y_1}{y_0}+O(\hbar^{-2}).
\]
比较上述 $e^{f_1/\hbar}HG[\mathcal I(\widehat P)](t)$ 与 $HG[\mathcal I(\widehat Q)](\Psi(t))$ 的表达式，就立即得到所需的式~\eqref{eq:quintic-transformed-hypergeometric-congruence}。定理得证。
\end{proof}
注意，以上证明给出的 $HG[\mathcal I(\widehat P)](t)$ 与 $HG[\mathcal I(\widehat Q)](t)$ 的展开，说明了选择镜像变换 $\mu$、$\nu$ 的依据：基本想法是确保同余式~\eqref{eq:quintic-transformed-hypergeometric-congruence} 成立。
\begin{corollary}\label{cor:quintic-hypergeometric-mirror-equality}
在非等变极限 $\lambda_i\to0$ 下，有
\[
\lim_{\lambda_i\to0}HG[\mathcal I(\widehat Q)](\Psi(t))
=\frac1{y_0}\lim_{\lambda_i\to0}HG[\mathcal I(\widehat P)](t).
\]
\end{corollary}
\begin{proof}
将定理~\ref{thm:quintic-transformed-euler-data-equality} 与式~\eqref{eq:quintic-first-mirror-transformation}、\eqref{eq:quintic-second-mirror-transformation} 合并，并取非等变极限，得到
\[
\lim_{\lambda_i\to0}HG[\mathcal I(\widehat Q)](\Psi(t))
=\lim_{\lambda_i\to0}e^{f_1/\hbar}HG[\mathcal I(\widehat P)](t).
\]
不过，由式~\eqref{eq:quintic-first-mirror-transformation-parameters}，知道
$f_1=-(\log y_0)\hbar+(h_2/y_0)\sum_{k=0}^4\lambda_k$。这给出 $\lim_{\lambda_i\to0}e^{f_1/\hbar}=1/y_0$，从而得到推论。
\end{proof}

\begin{proof}[定理~\ref{thm:lly-quintic-mirror} 的证明]
现在容易证明定理~\ref{thm:lly-quintic-mirror}。由定理~\ref{thm:critical-quintic-hypergeometric-potential}，有
\[
\begin{aligned}
&\lim_{\lambda_i\to0}HG[\mathcal I(\widehat Q)](\Psi(t))\\
&\quad=5H\left(1-\Psi(t)\frac H\hbar+\frac{\Phi'(\Psi(t))}{5}\frac{H^2}{\hbar^2}
-\frac{\Psi(t)\Phi'(\Psi(t))-2\Phi(\Psi(t))}{5}\frac{H^3}{\hbar^3}\right),
\end{aligned}
\]
又由式~\eqref{eq:quintic-hypergeometric-period-expansion}，有
\[
\lim_{\lambda_i\to0}HG[\mathcal I(\widehat P)](t)
=5H\left(y_0-y_1\frac H\hbar+y_2\frac{H^2}{\hbar^2}-y_3\frac{H^3}{\hbar^3}\right).
\]
把这些等式与推论~\ref{cor:quintic-hypergeometric-mirror-equality} 合并，得到
\[
\begin{aligned}
&5H\left(1-\Psi(t)\frac H\hbar+\frac{\Phi'(\Psi(t))}{5}\frac{H^2}{\hbar^2}
-\frac{\Psi(t)\Phi'(\Psi(t))-2\Phi(\Psi(t))}{5}\frac{H^3}{\hbar^3}\right)\\
&\qquad=5H\left(1-\frac{y_1}{y_0}\frac H\hbar
+\frac{y_2}{y_0}\frac{H^2}{\hbar^2}-\frac{y_3}{y_0}\frac{H^3}{\hbar^3}\right).
\end{aligned}
\]
比较系数，简单计算即得
\[
\Phi(\Psi(t))=\frac52\left(\frac{y_1}{y_0}\frac{y_2}{y_0}-\frac{y_3}{y_0}\right).
\]
这完成了五次三维簇\IndexLocation{indexloc:0864:377}\IndexLocation{indexloc:0867:377}镜像定理的证明。
\end{proof}
\SourcePage{378}
刚才的证明十分漂亮；若再认识到，本节先前的工作还能轻而易举地给出其他重要结果，就更能体会它的价值。\SourceCite{LLY} 给出了若干这样的例子，这里只介绍一个与我们的目标有关的例子。
\SourceExampleNumber{11.1.7}{1}
\begin{example}\label{ex:lly-multiple-cover-proof}
如例题~\ref{ex:critical-multiple-cover-bundle}，考虑 $\mathcal V=\mathcal O_{\mathbb P^1}(-1)\oplus\mathcal O_{\mathbb P^1}(-1)$。由式~\eqref{eq:critical-bundle-integral-coefficients}，知道
\begin{equation}\label{eq:lly-multiple-cover-integral}
\int_{\mathbb P^1}e^{-Ht/\hbar}
\frac{\lim_{\lambda_i\to0}I_d^*(\widehat Q_d)}{\prod_{m=1}^d(H-m\hbar)^2}
=\hbar^{-3}(2-dt)K_d,
\end{equation}
其中
\[
K_d=\int_{\overline M_{0,0}(\mathbb P^1,d)}
\operatorname{Euler}\!\left(R^1\pi_{1*}e_1^*
\bigl(\mathcal O_{\mathbb P^1}(-1)\oplus\mathcal O_{\mathbb P^1}(-1)\bigr)\right).
\]
可以如下计算 $K_d$。已知 $\widehat P$、$\widehat Q$ 是相关联的 Euler 数据，而引理~\ref{lem:geometric-euler-data-hypergeometric-bound} 给出 $\deg_{\hbar}I_d^*(\widehat Q_d)\le2d-2$。不过，在此情形下，式~\eqref{eq:concavex-hypergeometric-euler-data} 表明
\[
\widehat P_d=\prod_{m=1}^{d-1}(-\kappa+m\hbar)^2
=\prod_{m=1}^{d-1}(\kappa-m\hbar)^2.
\]
因此 $I_d^*(\widehat P_d)=\prod_{m=1}^{d-1}(p-m\hbar)^2$，其关于 $\hbar$ 的次数也不超过 $2d-2$。所以，此情形满足定理~\ref{thm:linked-euler-data-uniqueness} 的次数条件，从而 $\widehat P=\widehat Q$。于是
\[
\begin{aligned}
&\int_{\mathbb P^1}e^{-Ht/\hbar}
\frac{\lim_{\lambda_i\to0}I_d^*(\widehat Q_d)}{\prod_{m=1}^d(H-m\hbar)^2}\\
&\quad=\int_{\mathbb P^1}e^{-Ht/\hbar}
\frac{\lim_{\lambda_i\to0}I_d^*(\widehat P_d)}{\prod_{m=1}^d(H-m\hbar)^2}\\
&\quad=\int_{\mathbb P^1}e^{-Ht/\hbar}
\frac{\prod_{m=1}^{d-1}(H-m\hbar)^2}{\prod_{m=1}^d(H-m\hbar)^2}\\
&\quad=\int_{\mathbb P^1}e^{-Ht/\hbar}\frac1{(H-d\hbar)^2}
=\hbar^{-3}(2-dt)d^{-3},
\end{aligned}
\]
最后一个等号利用了 $H^*(\mathbb P^1)$ 中的 $H^2=0$。与式~\eqref{eq:lly-multiple-cover-integral} 比较，即得 $K_d=d^{-3}$；这完成了定理~\ref{thm:multiple-cover-localization} 的第二个证明。
\end{example}

这个例子表明，有必要拓宽我们对“镜像定理”的认识：导出五次三维簇镜像对称的同样思想与技巧，也会同时给出其他有趣结果。论文 \SourceCite{LLY} 是朝理解镜像对称的更广阔背景迈出的重要一步。

下一节将看到，与 \SourceCite{LLY} 相比，Givental 的镜像定理方法具有相似而又不同的结构。他考虑两个取上同调值的形式函数 $I_{\mathcal V}$（式~\eqref{eq:givental-complete-intersection-i-function}）和 $J_{\mathcal V}$（式~\eqref{eq:givental-general-j-function}），然后作变量变换，使一个唯一性定理适用。参见定理~\ref{thm:givental-projective-mirror} 及 \SourceCite{Givental2}。例题~\ref{ex:givental-quintic-mirror-potential} 将以五次三维簇说明此方法；\hyperref[subsec:givental-general-bundle-proof]{第~\ref*{subsec:givental-general-bundle-proof} 小节}末附近还会进一步比较这两种方法。
\section{Givental 的方法}\label{sec:givental-mirror-approach}
\SourcePage{379}
本节解释 Givental 对射影空间中完全交的镜像定理的表述（定理~\ref{thm:givental-projective-mirror}）。其方法涉及构造取形式上同调值的函数 $I_{\mathcal V}$、$J_{\mathcal V}$，其中 $J_{\mathcal V}$ 是\hyperref[sec:quantum-cohomology-relations]{第~\ref*{sec:quantum-cohomology-relations} 节}讨论的函数 $J$ 的一个变体。镜像定理则断言，作变量变换后，$I_{\mathcal V}=J_{\mathcal V}$。

\SourceCite{BDPP,Pandharipande3} 更详细地解释了 \SourceCite{Givental2} 的方法，也可作为本节的有用参考文献。五次簇的情形见 \SourceCite{Givental6}。这里将综合 \SourceCite{Givental2,Givental4,Kim2,Pandharipande3} 中的思想，概述 Givental 关于 $\mathbb P^n$ 中 nef 完全交的定理的一个略有不同的证明；读者可参考这些文献。\SourceCite{Givental2} 引入的一些所需思想，已在\hyperref[chap:localization]{第九章}和\hyperref[chap:quantum-differential-equations]{第十章}解释过。还将讨论环面完全交，并给出一些说明性的例子。
\subsection{射影空间中 nef 完全交的镜像定理}\label{subsec:givental-nef-projective-mirror}
考虑 $\mathbb P^n$ 中由 $\ell$ 个次数分别为 $a_1,\ldots,a_\ell$ 的超曲面构成的光滑完全交 $X$。令 $\mathcal V=\bigoplus_{i=1}^{\ell}\mathcal O_{\mathbb P^n}(a_i)$，于是 $X$ 是 $\mathcal V$ 的一个整体截面的零点集。还假设 $X$ 的反典范类为 nef；在此情形下，这等价于 $\sum_{i=1}^{\ell}a_i\le n+1$。

首先要严谨陈述定理~\ref{thm:givental-projective-mirror}，它联系了下文定义的两个取上同调值的形式函数 $I_{\mathcal V}$、$J_{\mathcal V}$。其中 $I_{\mathcal V}$ 较为初等。遵循 \SourceCite{Givental2,Kim2}，定义
\begin{equation}\label{eq:givental-complete-intersection-i-function}\index{I函数@$I$ 函数|hyperpage}
\begin{aligned}
I_{\mathcal V}&=I_{\mathcal V}(t_0,t_1,\hbar^{-1})\\
&=e^{(t_0+t_1H)/\hbar}\operatorname{Euler}(\mathcal V)
\sum_{d=0}^{\infty}e^{dt_1}
\frac{\prod_{i=1}^{\ell}\prod_{m=1}^{a_id}(a_iH+m\hbar)}
{\prod_{m=1}^d(H+m\hbar)^{n+1}},
\end{aligned}
\end{equation}
其中 $H\in H^2(\mathbb P^n)$ 是超平面类。记号 $I_{\mathcal V}(t_0,t_1,\hbar^{-1})$ 表示，有时把 $I_{\mathcal V}$ 看作按 $\hbar^{-1}$ 的幂展开的形式级数。若 $X=\mathbb P^n$，即 $\mathcal V=0$，则把 $I_{\mathcal V}$ 记为 $I_{\mathbb P^n}$。

另一方面，形式函数 $J_{\mathcal V}$ 的定义较为复杂。\hyperref[chap:quantum-differential-equations]{第十章}研究了 Givental 函数 $J=J_X$。若 $\dim X\ge3$，则 $H^2(X)$ 由 $i^*H$ 生成，于是引理~\ref{lem:giventals-j-two-formulas} 给出 $J_X$ 的两个公式：
\begin{equation}\label{eq:complete-intersection-ordinary-j-function}
\begin{aligned}
J_X&=e^{(t_0+t_1i^*H)/\hbar}
\left(1+\sum_{d=1}^{\infty}\sum_{a=0}^m e^{dt_1}
\left\langle\frac{T_a}{\hbar-c},1\right\rangle_{0,d}T^a\right)\\
&=e^{(t_0+t_1i^*H)/\hbar}
\left(1+\sum_{d=1}^{\infty}e^{dt_1}PD^{-1}e_{1*}
\left(\frac1{\hbar-c}\cap[\overline M_{0,2}(X,d)]^{\mathrm{virt}}\right)\right).
\end{aligned}
\end{equation}
这里，$T_a$ 是 $H^*(X)$ 的上同调基，$T^a$ 是其对偶基，$PD$ 是 Poincar\'e 对偶，$e_1:\overline M_{0,2}(X,d)\to X$ 是赋值映射，而
\begin{equation}\label{eq:complete-intersection-symbolic-descendants}
\left\langle\frac{T_a}{\hbar-c},1\right\rangle_{0,d}
=\sum_{k=0}^{\infty}\hbar^{-(k+1)}\langle\tau_kT_a,1\rangle_{0,d}
\end{equation}
是\hyperref[chap:quantum-differential-equations]{第十章}采用的引力关联数的符号记法。此前把 Givental 函数简记为 $J$，但现在须区分 $X$ 与 $\mathbb P^n$，故改记为 $J_X$。函数 $J_X$ 在量子微分方程理论中起核心作用（见定理~\ref{thm:quantum-differential-relations}）。
\SourcePage{380}
注意，不能直接比较 $I_{\mathcal V}$ 与 $J_X$，因为前者取值于 $H^*(\mathbb P^n)$，后者取值于 $H^*(X)$。为此，将引入一个取值于 $H^*(\mathbb P^n)$ 的 $J$ 函数变体 $J_{\mathcal V}$。然后，Givental 的定理将说明，当 $X$ 的反典范类为 nef 时，$I_{\mathcal V}$ 与 $J_{\mathcal V}$ 如何联系。

解释此定理的安排如下。先讨论 $\mathcal V=0$ 的情形；此时 Givental 的定理就是简单的等式 $I_{\mathbb P^n}=J_{\mathbb P^n}$（见下文命题~\ref{prop:projective-space-i-equals-j}）。然后定义 $\mathbb P^n$ 上一般丛所对应的 $J_{\mathcal V}$，并陈述 Givental 的镜像定理（定理~\ref{thm:givental-projective-mirror}）。证明将推迟至\hyperref[subsec:givental-projective-space-proof]{第~\ref*{subsec:givental-projective-space-proof}}、\hyperref[subsec:givental-general-bundle-proof]{\ref*{subsec:givental-general-bundle-proof} 小节}。还将解释 $J_{\mathcal V}$ 与 $J_X$ 的联系，并说明定理~\ref{thm:givental-projective-mirror} 蕴含五次三维簇的镜像对称。

此时可以联系 \SourceCite{LLY} 的方法。Givental 的函数 $I_{\mathcal V}$ 在 $t_0=0$ 时，除去 $\hbar$ 的符号变化，正是以 \SourceCite{LLY} 记号表示的\IndexLocation{indexloc:0945:380}非等变极限 $\lim_{\lambda_i\to0}HG[\mathcal I(\widehat P)](t_1)$。在将本节计算与 \SourceCite{LLY} 的相应计算比较时，不再明确重复指出这一符号差别。此外，一般情形下，定理~\ref{thm:givental-projective-mirror} 的内容是：变量变换后，$I_{\mathcal V}$ 与 $J_{\mathcal V}$ 相同。在 \SourceCite{LLY} 的方法中，对应陈述则是：作同样变量变换后，$HG[\mathcal I(\widehat P)](t_1)$ 与 $HG[\mathcal I(\widehat Q)](t_1)$ 相同。

现在可以陈述 Givental 版本镜像定理在 $X=\mathbb P^n$，即射影空间本身这一特殊情形下的结果。用式~\eqref{eq:givental-complete-intersection-i-function} 中 $\ell=0$ 定义 $I_{\mathbb P^n}$，并如式~\eqref{eq:complete-intersection-ordinary-j-function} 定义 $J_{\mathbb P^n}$。
\begin{proposition}\label{prop:projective-space-i-equals-j}
对于射影空间 $\mathbb P^n$，有 $I_{\mathbb P^n}=J_{\mathbb P^n}$。
\end{proposition}
将在\hyperref[subsec:givental-projective-space-proof]{第~\ref*{subsec:givental-projective-space-proof} 小节}证明此命题。不过，\hyperref[chap:quantum-differential-equations]{第十章}已验证了 $n=1$ 的情形，下面先回顾这一结果。
\SourceExampleNumber{11.2.1}{1}
\begin{example}\label{ex:projective-line-i-j-equality}
当 $X=\mathbb P^1$ 时，由\hyperref[subsec:gravitational-correlator-properties]{第~\ref*{subsec:gravitational-correlator-properties} 小节}的次数公理，式~\eqref{eq:complete-intersection-ordinary-j-function} 的第一行简化为
\[
J_{\mathbb P^1}=e^{(t_0+t_1H)/\hbar}
\left(1+\sum_{d=1}^{\infty}\left(
\frac{e^{dt_1}}{\hbar^{2d}}\langle\tau_{2d-1}H,1\rangle_{0,d}1
+\frac{e^{dt_1}}{\hbar^{2d+1}}\langle\tau_{2d},1\rangle_{0,d}H\right)\right).
\]
此外，式~\eqref{eq:givental-complete-intersection-i-function} 在 $n=1$、$\ell=0$ 时为
\[
I_{\mathbb P^1}=e^{(t_0+t_1H)/\hbar}\sum_{d=0}^{\infty}
\frac{e^{dt_1}}{\bigl((H+\hbar)(H+2\hbar)\cdots(H+d\hbar)\bigr)^2}.
\]
例题~\ref{ex:projective-line-gravitational} 计算了 $\langle\tau_{2d-1}H,1\rangle_{0,d}$ 与 $\langle\tau_{2d},1\rangle_{0,d}$；例题~\ref{ex:projective-line-j-function} 使用这些公式证明了 $I_{\mathbb P^1}=J_{\mathbb P^1}$。

从命题~\ref{prop:projective-space-i-equals-j} 的角度看，也可反过来进行这一过程。一旦知道 $I_{\mathbb P^1}=J_{\mathbb P^1}$，按 $H$ 的幂展开 $I_{\mathbb P^1}$（注意，$H^*(\mathbb P^1)$ 中 $H^2=0$），就容易得到 $\langle\tau_{2d-1}H,1\rangle_{0,d}$ 与 $\langle\tau_{2d},1\rangle_{0,d}$ 的公式。因此，可把命题~\ref{prop:projective-space-i-equals-j} 看作对 $\mathbb P^1$ 的\IndexLocation{indexloc:0463:380}\IndexLocation{indexloc:0469:380}引力关联数的计算。
\end{example}

接下来，定义上述 $\mathcal V=\bigoplus_{i=1}^{\ell}\mathcal O_{\mathbb P^n}(a_i)$ 对应的 $J_{\mathcal V}$。先引入一些记号。给定任意 $d>0$，当 $f:C\to\mathbb P^n$ 在 $\overline M_{0,k}(\mathbb P^n,d)$ 中变化时，截面空间 $H^0(C,f^*\mathcal V)$ 构成该模空间上的向量丛 $\mathcal V_{d,k}$。更准确地，令 $\pi_{k+1}:\overline M_{0,k+1}(\mathbb P^n,d)\to\overline M_{0,k}(\mathbb P^n,d)$ 为忘掉第 $k+1$ 个标记点的映射，令 $e_{k+1}:\overline M_{0,k+1}(\mathbb P^n,d)\to\mathbb P^n$ 为该标记点处的\IndexLocation{indexloc:0276:380}赋值映射。于是
\[
\mathcal V_{d,k}=\pi_{k+1*}e_{k+1}^*\mathcal V.
\]
\SourcePage{381}
特别地，$\mathcal V_{d,0}=\pi_{1*}e_1^*\mathcal V$，恰是本章先前记作 $\mathcal V_d$ 的丛。还可证明，当 $k>0$ 时，$\mathcal V_{d,k}=\pi_k^*(\mathcal V_{d,k-1})$。

现在固定 $1\le i\le k$。对 $\overline M_{0,k}(\mathbb P^n,d)$ 中每个 $f:C\to\mathbb P^n$，考虑 $H^0(C,f^*\mathcal V)$ 中在第 $i$ 个标记点处消失的截面。它们构成 $\overline M_{0,k}(\mathbb P^n,d)$ 上丛 $\mathcal V'_{d,k,i}$ 的纤维，且有正合列
\begin{equation}\label{eq:marked-bundle-evaluation-exact-sequence}
0\longrightarrow\mathcal V'_{d,k,i}\longrightarrow\mathcal V_{d,k}
\longrightarrow e_i^*\mathcal V\longrightarrow0.
\end{equation}
这给出
\begin{equation}\label{eq:marked-bundle-euler-factorization}
\operatorname{Euler}(\mathcal V_{d,k})
=\operatorname{Euler}(\mathcal V'_{d,k,i})\cup e_i^*\operatorname{Euler}(\mathcal V).
\end{equation}
特别地，得到 $\overline M_{0,2}(\mathbb P^n,d)$ 上的丛 $\mathcal V'_{d,2,1}$。利用此丛，以以下公式定义 $J_{\mathcal V}$：
\begin{equation}\label{eq:givental-projective-bundle-j-function}
\begin{aligned}
J_{\mathcal V}&=e^{(t_0+t_1H)/\hbar}\operatorname{Euler}(\mathcal V)\times\\
&\quad\left(1+\sum_{d=1}^{\infty}\sum_{b=0}^{\infty}
\frac{e^{dt_1}}{\hbar^{b+1}}
 e_{1!}\bigl(c_1(\mathcal L_1)^b\cup\operatorname{Euler}(\mathcal V'_{d,2,1})\bigr)\right).
\end{aligned}
\end{equation}
用上述式~\eqref{eq:complete-intersection-ordinary-j-function} 的\IndexLocation{indexloc:0411:381}符号记法，可写成
\begin{equation}\label{eq:givental-projective-bundle-j-symbolic}
J_{\mathcal V}=e^{(t_0+t_1H)/\hbar}\operatorname{Euler}(\mathcal V)
\left(1+\sum_{d=1}^{\infty}e^{dt_1}e_{1!}
\left(\frac{\operatorname{Euler}(\mathcal V'_{d,2,1})}{\hbar-c}\right)\right).
\end{equation}
注意 $J_{\mathcal V}$ 的两个重要性质。首先，当 $\mathcal V=0$ 时，对应完全交为 $X=\mathbb P^n$。因为 $\mathbb P^n$ 为凸，$\overline M_{0,2}(\mathbb P^n,d)$ 的\IndexLocation{indexloc:1406:381}虚拟基本类就是通常基本类，因而式~\eqref{eq:complete-intersection-ordinary-j-function} 的第二个公式简化为
\[
J_{\mathbb P^n}=e^{(t_0+t_1H)/\hbar}
\left(1+\sum_{d=1}^{\infty}e^{dt_1}e_{1!}\left(\frac1{\hbar-c}\right)\right).
\]
零丛的 Euler 类为 $1$，故这与式~\eqref{eq:givental-projective-bundle-j-symbolic} 在 $\mathcal V=0$ 时一致。因此，当 $\mathcal V=0$ 时，$J_{\mathcal V}=J_{\mathbb P^n}$。

再看第二个性质。考虑式~\eqref{eq:givental-projective-bundle-j-function} 的括号表达式；其 $\hbar^{-1}$ 的系数包含因子
\[
\operatorname{Euler}(\mathcal V)\cup e_{1!}\bigl(\operatorname{Euler}(\mathcal V'_{d,2,1})\bigr)
=e_{1!}\bigl(\operatorname{Euler}(\mathcal V_{d,2})\bigr)=0.
\]
第一个等号使用了式~\eqref{eq:marked-bundle-euler-factorization}。第二个等号源于 $e_1$ 可经由忘掉第二个标记点的映射 $\pi_2:\overline M_{0,2}(\mathbb P^n,d)\to\overline M_{0,1}(\mathbb P^n,d)$ 分解。上文已说明 $\mathcal V_{d,2}=\pi_2^*\mathcal V_{d,1}$，故
\[
\pi_{2!}\bigl(\operatorname{Euler}(\mathcal V_{d,2})\bigr)
=\pi_{2!}\pi_2^*\bigl(\operatorname{Euler}(\mathcal V_{d,1})\bigr)=0,
\]
因为 $\pi_2$ 的纤维具有正维数。这给出
\begin{equation}\label{eq:givental-bundle-j-normalization}
J_{\mathcal V}=e^{(t_0+t_1H)/\hbar}\operatorname{Euler}(\mathcal V)
\bigl(1+o(\hbar^{-1})\bigr),
\end{equation}
类似于 Givental $J$ 函数的性质~\eqref{eq:j-function-asymptotic-normalization}。

现在可以陈述 Givental 关于 $\mathbb P^n$ 中\IndexLocation{indexloc:0856:381}完全交的镜像定理。
\begin{theorem}\label{thm:givental-projective-mirror}\index{镜像定理!射影完全交|hyperpage}
设 $X\subset\mathbb P^n$ 为 $\ell$ 个次数为 $a_i$ 的超曲面的完全交，且 $\sum_{i=1}^{\ell}a_i\le n+1$。若 $\mathcal V=\bigoplus_{i=1}^{\ell}\mathcal O_{\mathbb P^n}(a_i)$，则取形式上同调值的函数 $I_{\mathcal V}$、$J_{\mathcal V}$ 在以下加权齐次的三角变量变换后相同：
\[
t_0\longmapsto t_0+f(e^{t_1})\hbar+h(e^{t_1}),\qquad
 t_1\longmapsto t_1+g(e^{t_1}),
\]
\SourcePage{382}
其中 $f,g,h$ 是加权齐次幂级数，满足 $\deg f=\deg g=0$、$\deg h=1$，且 $\deg e^{t_1}=n+1-\sum_i a_i$。
\end{theorem}

如先前所述，条件 $\sum_{i=1}^{\ell}a_i\le n+1$ 表示 $X$ 的反典范类为 nef。特别地，$\sum_{i=1}^{\ell}a_i$ 的最大可能值为 $n+1$；椭圆曲线、K3 曲面或维数至少为 $3$ 的 Calabi--Yau 流形都出现这种情形。还应注意，该定理包含命题~\ref{prop:projective-space-i-equals-j} 作为特例（$\mathcal V=0$）。将在\hyperref[subsec:givental-general-bundle-proof]{第~\ref*{subsec:givental-general-bundle-proof} 小节}概述定理~\ref{thm:givental-projective-mirror} 的证明。

另一个有用的说明是，定理~\ref{thm:givental-projective-mirror} 中的幂级数 $f,g,h$ 可由 $I_{\mathcal V}$ 的系数唯一确定。为此，对 $J_{\mathcal V}$ 作变量代换 $t_0\mapsto t_0+f(e^{t_1})\hbar+h(e^{t_1})$、$t_1\mapsto t_1+g(e^{t_1})$。略去 $I_{\mathcal V}$ 与 $J_{\mathcal V}$ 中的因子 $e^{t_0/\hbar}$，则由式~\eqref{eq:givental-bundle-j-normalization}，代换后的 $J_{\mathcal V}$ 为
\[
J_{\mathcal V}=e^f\operatorname{Euler}(\mathcal V)
\left(1+h\hbar^{-1}1+(t_1+g)\hbar^{-1}H+o(\hbar^{-1})\right).
\]
令它与 $I_{\mathcal V}$ 相等，即得
\begin{equation}\label{eq:givental-mirror-variable-coefficients}
\begin{aligned}
e^f&\longleftrightarrow\widetilde I_{\mathcal V}\text{ 中 }1\text{ 的系数},\\
e^fh&\longleftrightarrow\widetilde I_{\mathcal V}\text{ 中 }\hbar^{-1}1\text{ 的系数},\\
e^f(t_1+g)&\longleftrightarrow\widetilde I_{\mathcal V}\text{ 中 }\hbar^{-1}H\text{ 的系数},
\end{aligned}
\end{equation}
其中 $\widetilde I_{\mathcal V}$ 是式~\eqref{eq:givental-complete-intersection-i-function} 去掉因子 $\operatorname{Euler}(\mathcal V)$ 后的函数（故 $I_{\mathcal V}=\operatorname{Euler}(\mathcal V)\widetilde I_{\mathcal V}$）。如所断言，$I_{\mathcal V}$ 确定了镜像定理中的坐标变换。特别地，如果在原有变量 $t_0,t_1$ 下就已有
\[
I_{\mathcal V}=e^{(t_0+t_1H)/\hbar}\operatorname{Euler}(\mathcal V)
\bigl(1+o(\hbar^{-1})\bigr),
\]
则必有 $f=g=h=0$，从而 $I_{\mathcal V}=J_{\mathcal V}$。以下是一个这样的例子。
\SourceExampleNumber{11.2.1}{2}
\begin{example}\label{ex:fano-hypersurface-i-j-equality}
设 $X\subset\mathbb P^n$ 是次数 $\ell<n$ 的超曲面。关于 $\ell$ 的假设容易给出 $I_{\mathcal V}=e^{(t_0+t_1H)/\hbar}\operatorname{Euler}(\mathcal V)(1+o(\hbar^{-1}))$，故由上述讨论，在定理~\ref{thm:givental-projective-mirror} 中有 $f=g=h=0$。因此
\[
J_{\mathcal V}=I_{\mathcal V}=\ell H e^{(t_0+t_1H)/\hbar}
\left(1+\sum_{d=1}^{\infty}e^{dt_1}
\frac{(\ell H+\hbar)(\ell H+2\hbar)\cdots(\ell H+\ell d\hbar)}
{\bigl((H+\hbar)(H+2\hbar)\cdots(H+d\hbar)\bigr)^{n+1}}\right).
\]
很快将看到，这对 $X$ 的\IndexLocation{indexloc:1129:382}小量子上同调环有一些有趣的推论。
\end{example}

为充分理解定理~\ref{thm:givental-projective-mirror}，须先解释 $J_X$ 与 $J_{\mathcal V}$ 的关系。为此，使用包含映射 $i:X\hookrightarrow\mathbb P^n$ 诱导的 Gysin 映射 $i_!:H^*(X)\to H^{*+2\ell}(\mathbb P^n)$。有以下结果。
\begin{proposition}\label{prop:complete-intersection-j-pushforward}
设 $i:X\hookrightarrow\mathbb P^n$ 为 $\mathcal V=\bigoplus_{i=1}^{\ell}\mathcal O_{\mathbb P^n}(a_i)$ 所确定的完全交，且 $\dim X\ge3$。则 $J_{\mathcal V}=i_!(J_X)$。
\end{proposition}
\begin{proof}
$\dim X\ge3$ 的假设给出 $H^2(\mathbb P^n)\simeq H^2(X)$，因而 $J_{\mathcal V}$ 与 $J_X$ 都是 $t_0,t_1$ 的形式函数。还应注意，$i:X\hookrightarrow\mathbb P^n$ 诱导包含映射 $j:\overline M_{0,2}(X,d)\hookrightarrow\overline M_{0,2}(\mathbb P^n,d)$，使下图交换：
\[
\begin{tikzcd}[column sep=large,row sep=large]
\overline M_{0,2}(X,d)\arrow[r,"j"]\arrow[d,"e_1"']&\overline M_{0,2}(\mathbb P^n,d)\arrow[d,"e_1"]\\
X\arrow[r,"i"]&\mathbb P^n
\end{tikzcd}
\]
\SourcePage{383}
这里 $e_1$ 为赋值映射。

关键事实是，可以计算 $\overline M_{0,2}(X,d)$ 的虚拟基本类。按定义，$X$ 为 $\mathcal V$ 的一个截面 $s$ 的零点集，而 $s$ 诱导了上述丛 $\mathcal V_{d,k}$ 的一个截面 $\widetilde s$。与例题~\ref{ex:quintic-virtual-class} 一样，$\widetilde s$ 的零点概形（按代数叠的意义）自然等同于 $\overline M_{0,2}(X,d)$。于是可对向量丛 $\mathcal V_{d,2}$ 使用法锥构造，得到 $[\overline M_{0,2}(X,d)]^{\mathrm{virt}}$。推广式~\eqref{eq:quintic-virtual-class-pushforward}，可得
\begin{equation}\label{eq:complete-intersection-virtual-pushforward}
j_*\bigl([\overline M_{0,2}(X,d)]^{\mathrm{virt}}\bigr)
=\operatorname{Euler}(\mathcal V_{d,2})\cap[\overline M_{0,2}(\mathbb P^n,d)].
\end{equation}
因此，$\overline M_{0,2}(X,d)$ 的虚拟基本类是 $\mathcal V_{d,2}$ 的 Euler 类的一个精细化。

现在证明命题。从式~\eqref{eq:complete-intersection-ordinary-j-function} 给出的 $J_X$ 的第二个公式出发，利用投影公式及上述交换图，容易看到
\[
\begin{aligned}
i_!(J_X)&=e^{(t_0+t_1H)/\hbar}\left(\operatorname{Euler}(\mathcal V)+\sum_{d=1}^{\infty}e^{dt_1}
 e_{1!}PD^{-1}j_*\left(\frac1{\hbar-c}\cap[\overline M_{0,2}(X,d)]^{\mathrm{virt}}\right)\right).
\end{aligned}
\]
在符号记法 $1/(\hbar-c)$ 中，$c$ 表示 $\overline M_{0,2}(X,d)$ 上丛 $\mathcal L_1$ 的 Chern 类 $c_1(\mathcal L_1)$。它是 $\overline M_{0,2}(\mathbb P^n,d)$ 上同样记为 $\mathcal L_1$ 的丛经 $j$ 拉回所得。因此，由 $j$ 的投影公式，
\[
\begin{aligned}
j_*\left(\frac1{\hbar-c}\cap[\overline M_{0,2}(X,d)]^{\mathrm{virt}}\right)
&=\frac1{\hbar-c}\cap j_*\bigl([\overline M_{0,2}(X,d)]^{\mathrm{virt}}\bigr)\\
&=\frac1{\hbar-c}\cap\bigl(\operatorname{Euler}(\mathcal V_{d,2})\cap[\overline M_{0,2}(\mathbb P^n,d)]\bigr)\\
&=\frac{\operatorname{Euler}(\mathcal V_{d,2})}{\hbar-c}\cap[\overline M_{0,2}(\mathbb P^n,d)].
\end{aligned}
\]
于是 $i_!(J_X)$ 的公式简化为
\[
i_!(J_X)=e^{(t_0+t_1H)/\hbar}
\left(\operatorname{Euler}(\mathcal V)+\sum_{d=1}^{\infty}e^{dt_1}
 e_{1!}\left(\frac{\operatorname{Euler}(\mathcal V_{d,2})}{\hbar-c}\right)\right).
\]
不过，式~\eqref{eq:marked-bundle-euler-factorization} 给出
\[
\operatorname{Euler}(\mathcal V_{d,2})
=\operatorname{Euler}(\mathcal V'_{d,2,1})\cup e_1^*\operatorname{Euler}(\mathcal V).
\]
将它代入刚得到的 $i_!(J_X)$ 公式，并用 $e_1$ 的投影公式，得到
\[
i_!(J_X)=e^{(t_0+t_1H)/\hbar}\operatorname{Euler}(\mathcal V)
\left(1+\sum_{d=1}^{\infty}e^{dt_1}e_{1!}
\left(\frac{\operatorname{Euler}(\mathcal V'_{d,2,1})}{\hbar-c}\right)\right),
\]
这恰是 $J_{\mathcal V}$ 的定义。命题得证。
\end{proof}

既然已经知道 $i_!(J_X)=J_{\mathcal V}$，接下来研究从 $J_X$ 到 $J_{\mathcal V}$ 是否会丢失信息。使用以下两点观察：
\begin{itemize}
\item $X$ 是丰沛超曲面的完全交，公式 $i_!(i^*\alpha)=\alpha\cup\operatorname{Euler}(\mathcal V)$ 表明，$i_!$ 在 $i^*$ 的像上是单射。注意，$i^*H^*(\mathbb P^n)\subset H^{\mathrm{even}}(X)$。
\item \hyperref[subsec:gravitational-correlator-properties]{第~\ref*{subsec:gravitational-correlator-properties} 小节}的\IndexLocation{indexloc:0461:383}次数公理表明，$J_X$ 只涉及 $X$ 的偶次上同调。
\end{itemize}
\SourcePage{384}
当 $\dim X\ge3$ 为奇数时，由 Lefschetz 定理，$H^{\mathrm{even}}(\mathbb P^n)\to H^{\mathrm{even}}(X)$ 为满射。第一点说明，$i_!$ 在 $H^{\mathrm{even}}(X)$ 上为单射；结合第二点，可知能从 $J_{\mathcal V}=i_!(J_X)$ 恢复 $J_X$。

另一方面，当 $\dim X\ge3$ 为偶数时，记 $\dim X=2m$。此时，$J_X$ 仍取值于 $H^{\mathrm{even}}(X)$，但 $i^*:H^{\mathrm{even}}(\mathbb P^n)\to H^{\mathrm{even}}(X)$ 未必为满射。问题在于，$H^{2m}(X)$ 有 Lefschetz 分解
\[
H^{2m}(X)=\mathbb C\cdot i^*(H^m)\oplus H_0^{2m}(X),
\]
其中 $H_0^{2m}(X)$ 是本原上同调。$i^*$ 的像完全不包含 $H_0^{2m}(X)$ 中的非零元素；若 $J_X$ 有本原部分的分量，它便不在这个像中。对 $i^*:H^*(\mathbb P^n)\to H^*(X)$ 的像，有不同记号。按\hyperref[subsec:toric-mirror-conjecture]{第~\ref*{subsec:toric-mirror-conjecture} 小节}的术语，因为 $\mathbb P^n$ 为环面簇，把此像记作 $H_{\mathrm{toric}}^*(X)$。\SourceCite{Givental3} 则记作 $H^*(\mathcal V)$，这里采用后一记号。利用上述 $H^{2m}(X)$ 的分解，得到
\[
H^{\mathrm{even}}(X)=H^*(\mathcal V)\oplus H_0^{2m}(X),
\]
因为 $X$ 是完全交。这诱导了分解
\begin{equation}\label{eq:complete-intersection-j-primitive-decomposition}
J_X=J_X^{\mathcal V}+J_X^0,
\end{equation}
其中 $J_X^{\mathcal V}$ 取值于 $H^*(\mathcal V)$，$J_X^0$ 是 $J_X$ 的“本原部分”。为联系它与 $i_!(J_X)$，回顾 $H_0^{2m}(X)$ 正是 $i_!:H^{2m}(X)\to H^{2m+2\ell}(\mathbb P^n)$ 的核。对式~\eqref{eq:complete-intersection-j-primitive-decomposition} 应用 $i_!$，即得 $J_{\mathcal V}=i_!(J_X^{\mathcal V})$。因此，对于偶数维完全交，$J_{\mathcal V}$ 反映的是 $J_X$ 中来自 $\mathbb P^n$ 的部分。

可以把上述观察统一起来：式~\eqref{eq:complete-intersection-j-primitive-decomposition} 的分解 $J_X=J_X^{\mathcal V}+J_X^0$ 对任意维数都有意义（所以，当 $\dim X$ 为奇数时，$J_X^0=0$）。上述讨论表明，对所有 $X$，均可从 $J_{\mathcal V}=i_!(J_X^{\mathcal V})$ 恢复 $J_X^{\mathcal V}$。因此，镜像定理告诉我们如何计算 $X$ 的一部分亏格零引力关联数，即涉及来自环境空间的上同调类的关联数。

还可给出一个完全在 $H^*(\mathcal V)$ 内表述的镜像定理版本。若写成 $I_{\mathcal V}=\operatorname{Euler}(\mathcal V)\widetilde I_{\mathcal V}$，则 $i^*(\widetilde I_{\mathcal V})$ 与 $J_X^{\mathcal V}$ 都是取值于 $H^*(\mathcal V)$ 的形式函数。又因为 $i_!$ 在 $H^*(\mathcal V)$ 上为单射，分别把 $i^*(\widetilde I_{\mathcal V})$ 与 $J_X^{\mathcal V}$ 映到 $I_{\mathcal V}$ 与 $J_{\mathcal V}$，故定理~\ref{thm:givental-projective-mirror} 等价于断言：作适当类型的坐标变换后，$i^*(\widetilde I_{\mathcal V})$ 与 $J_X^{\mathcal V}$ 相等。这是 Givental 在 \SourceCite{Givental3} 中陈述镜像定理的方式，不过他允许环境空间为环面簇。\hyperref[subsec:givental-toric-mirror]{第~\ref*{subsec:givental-toric-mirror} 小节}将讨论这个更一般的情形。还应指出，这里使用 $H^*(\mathbb P^n)$ 而非 $H^*(\mathcal V)$ 的方法，基于 \SourceCite{Kim2}。

为看看镜像定理在具体情形下的样子，考虑五次三维簇。
\SourceExampleNumber{11.2.1}{3}
\begin{example}\leavevmode\phantomsection\label{ex:givental-quintic-mirror-potential}%
五次三维簇 $V\subset\mathbb P^4$ 对应于 $\mathcal V=\mathcal O_{\mathbb P^4}(5)$，故 $\operatorname{Euler}(\mathcal V)=5H$。这里 $H$ 是 $H^*(\mathbb P^4)$ 中的超平面类。于是
\[
I_{\mathcal V}=e^{(t_0+t_1H)/\hbar}5H\sum_{d=0}^{\infty}e^{dt_1}
\frac{(5H+\hbar)\cdots(5H+5d\hbar)}
{\bigl((H+\hbar)\cdots(H+d\hbar)\bigr)^5}.
\]
\SourcePage{385}
与式~\eqref{eq:quintic-hypergeometric-period-expansion} 比较，公式相同，只是 $\hbar$ 的符号不同，并多了一个因子 $e^{t_0/\hbar}$。因此
\[
I_{\mathcal V}=e^{t_0/\hbar}5H
\left(y_0+y_1\frac H\hbar+y_2\frac{H^2}{\hbar^2}+y_3\frac{H^3}{\hbar^3}\right),
\]
其中 $y_0,y_1,y_2,y_3$ 是\hyperref[subsec:quintic-mirror-theorem-statement]{第~\ref*{subsec:quintic-mirror-theorem-statement} 小节}讨论的五次镜像的 Picard--Fuchs 方程的一组解基。

镜像定理告诉我们，作适当变量变换后，$I_{\mathcal V}$ 等于 $J_{\mathcal V}$；由式~\eqref{eq:givental-mirror-variable-coefficients}，可从 $I_{\mathcal V}$ 读出 $f,g,h$。因此，上述 $I_{\mathcal V}$ 的公式给出
\[
\begin{aligned}
e^f=y_0&\ \Longrightarrow\ f=\log y_0,\\
e^fh=0&\ \Longrightarrow\ h=0,\\
e^f(t_1+g)=y_1&\ \Longrightarrow\ g=\frac{y_1}{y_0}-t_1.
\end{aligned}
\]
于是 $t_0\mapsto t_0+f\hbar=t_0+(\log y_0)\hbar$，$t_1\mapsto t_1+g=y_1/y_0$。注意，后者恰是\hyperref[subsec:quintic-mirror-theorem-statement]{第~\ref*{subsec:quintic-mirror-theorem-statement} 小节}使用的镜像映射。现在写 $J_{\mathcal V}=J_{\mathcal V}(t_0,t_1)$、$I_{\mathcal V}=I_{\mathcal V}(t_0,t_1)$。定理~\ref{thm:givental-projective-mirror} 给出
\[
J_{\mathcal V}\bigl(t_0+(\log y_0)\hbar,y_1/y_0\bigr)=I_{\mathcal V}(t_0,t_1).
\]
令 $s=y_1/y_0$，并注意 $e^{(t_0+\log(y_0)\hbar)/\hbar}=y_0e^{t_0/\hbar}$，上式就简化为
\[
J_{\mathcal V}(t_0,s)=\frac1{y_0}I_{\mathcal V}(t_0,t_1).
\]
结合上述 $I_{\mathcal V}$ 的公式，得到
\begin{equation}\label{eq:givental-quintic-j-periods}
J_{\mathcal V}(t_0,s)=e^{t_0/\hbar}5H
\left(1+s\frac H\hbar+\frac{y_2}{y_0}\frac{H^2}{\hbar^2}
+\frac{y_3}{y_0}\frac{H^3}{\hbar^3}\right).
\end{equation}
不过，\hyperref[subsec:cy-threefold-j-function]{第~\ref*{subsec:cy-threefold-j-function} 小节}还给出了 $J_V$ 的一些漂亮公式。使用 $H^{\mathrm{even}}(V)$ 的上同调基
$T_0=1$、$T_1=i^*H$、$T^1=i^*H^2/5$、$T^0=i^*H^3/5$。注意 $T_0,T^0$ 互为对偶，$T_1,T^1$ 亦然。若 $\Phi(s)$ 为五次三维簇的 Gromov--Witten 势，则命题~\ref{prop:cy-threefold-j-potential} 给出
\[
J_V(t_0,s)=e^{t_0/\hbar}
\left(1+\hbar^{-1}sT_1+\hbar^{-2}\Phi'(s)T^1
+\hbar^{-3}\bigl(s\Phi'(s)-2\Phi(s)\bigr)T^0\right).
\]
应用 $i_!$ 并使用命题~\ref{prop:complete-intersection-j-pushforward}，得到
\begin{equation}\label{eq:givental-quintic-j-potential}
J_{\mathcal V}(t_0,s)=e^{t_0/\hbar}5H
\left(1+s\frac H\hbar+\frac{\Phi'(s)}5\frac{H^2}{\hbar^2}
+\frac{s\Phi'(s)-2\Phi(s)}5\frac{H^3}{\hbar^3}\right).
\end{equation}
现在令 $s=\Psi(t_1)$。将式~\eqref{eq:givental-quintic-j-potential} 与式~\eqref{eq:givental-quintic-j-periods} 中的 $J_{\mathcal V}(t_0,s)$ 公式比较，容易得到
\[
\Phi(\Psi(t_1))=\frac52\left(\frac{y_1}{y_0}\frac{y_2}{y_0}-\frac{y_3}{y_0}\right).
\]
这恰是\hyperref[subsec:lly-quintic-mirror-proof]{第~\ref*{subsec:lly-quintic-mirror-proof} 小节}讨论的五次三维簇镜像定理。如定理~\ref{thm:quintic-mirror-potential-equivalence} 所指出，它蕴含\hyperref[sec:mirror-symmetry-revisited]{第~\ref*{sec:mirror-symmetry-revisited} 节}陈述的 Hodge 理论版本的镜像对称。

还须完成\hyperref[chap:quintic]{第二章}开始的讨论。回顾\hyperref[subsec:givental-ij]{第~\ref*{subsec:givental-ij} 小节}，那里给出了上述 $I_{\mathcal V}$、$J_{\mathcal V}$ 的公式，并讨论定理~\ref{thm:givental-projective-mirror} 如何
\SourcePage{386}
联系镜像对称。回看该小节，一个关键公式是式~\eqref{eq:iv-second-derivative}，它指出，当 $t_0=0$、$\hbar=1$ 时，有
\[
(I_{\mathcal V}/y_0)''=(5+2875q+\cdots)(H^3+sH^4),
\]
其中 $'$ 表示 $d/ds$。随后证明了括号中的级数 $5+2875q+\cdots$，正是五次镜像的归一化 Yukawa 耦合 $Y$。因此
\begin{equation}\label{eq:givental-quintic-yukawa-identity}
(I_{\mathcal V}/y_0)''=Y(H^3+sH^4).
\end{equation}
用它说明，Givental 的镜像定理如何给出五次三维簇的“经典”镜像对称表述。如上，定理~\ref{thm:givental-projective-mirror} 给出 $I_{\mathcal V}/y_0(t_0,t_1)=J_{\mathcal V}(t_0,s)$。因此
\[
(I_{\mathcal V}/y_0)''=J_{\mathcal V}'',
\]
其中 $'=d/ds$。不过，命题~\ref{prop:cy-threefold-j-potential} 表明
\[
J_V''=\hbar^{-2}e^{(t_0+sH)/\hbar}H*H,
\]
其中 $*$ 是 $H^*(V)$ 中的小量子乘积。严格地说，应将 $H$ 写为 $i^*H$，但这样会使记号过于繁杂。回顾式~\eqref{eq:one-parameter-quantum-yukawa}，$H*H=\langle H,H,H\rangle C$，其中 $C\in H^4(V)$ 是 $V$ 中一条直线的类。因此
\[
J_V''=\hbar^{-2}e^{(t_0+sH)/\hbar}\langle H,H,H\rangle C.
\]
令 $t_0=0$、$\hbar=1$。由于 $i_!(C)=H^3\in H^6(\mathbb P^4)$，有
\[
J_{\mathcal V}''=i_!(J_V)''=e^{sH}\langle H,H,H\rangle H^3
=\langle H,H,H\rangle(H^3+sH^4)
\]
在 $H^*(\mathbb P^4)$ 中成立。与式~\eqref{eq:givental-quintic-yukawa-identity} 比较，立即得到 $Y=\langle H,H,H\rangle$，从而恢复\hyperref[chap:quintic]{第二章}讨论的镜像对称版本。还应注意，五次三维簇的镜像对称可看作计算 $V$ 的瞬子数 $n_d$ 的一种方法。\SourceCite{Givental2} 解释了从镜像定理提取 $n_d$ 的另一种方式。
\end{example}

最后，解释 $J_{\mathcal V}$ 与 $J_X$ 的另一种联系，作为镜像定理讨论的收尾。其思想是，在 $H^*(\mathbb P^n)$ 上定义一个修改的“量子乘积” \SourceCite{Givental2,Kim2}，使得可用它构造 $J_{\mathcal V}$，就如 $J$ 是由通常的小量子乘积、式~\eqref{eq:giventals-flat-section-symbolic} 与式~\eqref{eq:giventals-j-pairing-definition} 得到的一样。\SourceCite{Pandharipande3} 详细定义并研究了这一乘积，记为 $*_X$。对我们而言，$*_X$ 最值得注意的性质如下。
\begin{proposition}\label{prop:modified-quantum-product-restriction}
设 $X\subset\mathbb P^n$ 是完全交。则
\[
i^*:H^*(\mathbb P^n)\otimes\mathbb C[[q]]\longrightarrow H^*(X)\otimes\mathbb C[[q]]
\]
是环同态，其中 $H^*(X)$ 采用小量子乘积 $*$，$H^*(\mathbb P^n)$ 采用修改的量子乘积 $*_X$。
\end{proposition}
\begin{proof}
证明见 \SourceCite[命题 4]{Pandharipande3}；该文将此结果归于 T. Graber。
\end{proof}
此外，函数 $J_{\mathcal V}$ 给出环 $(H^*(\mathbb P^n)\otimes\mathbb C[[q]],*_X)$ 中的关系，正如\hyperref[subsec:giventals-j-function]{第~\ref*{subsec:giventals-j-function} 小节}中 $J_X$ 给出 $X$ 的\IndexLocation{indexloc:1129:386}小\IndexLocation{indexloc:1127:386}\IndexLocation{indexloc:1128:386}量子上同调中的关系一样。
\begin{proposition}\label{prop:modified-quantum-j-relations}
设 $P(\hbar\partial/\partial t_1,e^{t_1},\hbar)J_{\mathcal V}=0$，其中 $P$ 是关于 $\hbar\partial/\partial t_1$ 的多项式，系数为 $e^{t_1}$、$\hbar$ 的幂级数。则关系 $P(H,q,0)=0$ 在 $(H^*(\mathbb P^n)\otimes\mathbb C[[q]],*_X)$ 中成立，即关系 $P(H,q,0)=0$ 在 $(H^*(\mathbb P^n)\otimes\mathbb C[[q]],*_X)/\ker i^*$ 中成立。
\end{proposition}
\SourcePage{387}
\begin{proof}
证明见 \SourceCite[引理 3]{Pandharipande3}。
\end{proof}
将命题~\ref{prop:modified-quantum-product-restriction} 与命题~\ref{prop:modified-quantum-j-relations} 合并，得到如下有用的推论。
\begin{corollary}\label{cor:complete-intersection-j-quantum-relations}
设 $P(\hbar\partial/\partial t_1,e^{t_1},\hbar)J_{\mathcal V}=0$，其中 $P$ 为上述形式的幂级数。则关系 $P(i^*H,q,0)=0$ 在 $X$ 的小量子上同调环中成立。
\end{corollary}
\begin{proof}
由命题~\ref{prop:modified-quantum-j-relations}，首先得到 $P(H,q,0)=0$，其中乘积采用 $*$。再由命题~\ref{prop:modified-quantum-product-restriction}，有 $i^*(P(H,q,0))=P(i^*H,q,0)$，右边按 $X$ 的小量子上同调环计算。因此，右边也为零。
\end{proof}
以下例子说明如何结合定理~\ref{thm:givental-projective-mirror} 与推论~\ref{cor:complete-intersection-j-quantum-relations}。
{\SourceExampleNumber{11.2.1}{2}
\renewcommand{\theHexam}{11.2.1.2.continued}
\begin{example}[续]\label{ex:fano-hypersurface-quantum-relation}
设 $X\subset\mathbb P^n$ 为次数 $\ell<n$ 的超曲面。在此例的第一部分，已经解释了 $\mathcal V=\mathcal O_{\mathbb P^n}(\ell)$ 的镜像定理如何给出等式
\[
J_{\mathcal V}=I_{\mathcal V}=\ell H e^{(t_0+t_1H)/\hbar}
\left(1+\sum_{d>0}e^{dt_1}
\frac{(\ell H+\hbar)(\ell H+2\hbar)\cdots(\ell H+\ell d\hbar)}
{\bigl((H+\hbar)(H+2\hbar)\cdots(H+d\hbar)\bigr)^{n+1}}\right).
\]
简单计算表明，$I_{\mathcal V}$ 被以下算子消去：
\[
\left(\hbar\frac d{dt_1}\right)^n
-e^{t_1}\ell\left(\ell\hbar\frac d{dt_1}+\hbar\right)
\cdots\left(\ell\hbar\frac d{dt_1}+(\ell-1)\hbar\right).
\]
因为 $J_{\mathcal V}=I_{\mathcal V}$，此算子也消去 $J_{\mathcal V}$。于是推论~\ref{cor:complete-intersection-j-quantum-relations} 给出超曲面 $X$ 的小量子上同调环中的关系
\[
H^n-q\ell(\ell H)^{\ell-1}=0,
\]
其中简记 $H=i^*H$。这与 \SourceCite{Beauville1} 所得到的关系一致。
\end{example}}

\subsection{量子超平面截面原理}\label{subsec:quantum-hyperplane-section-principle}\index{量子超平面截面原理|hyperpage}
迄今讨论的是完全交 $X\subset\mathbb P^n$ 的镜像定理。现在希望通过完全交的\emph{量子超平面截面原理}，把它放在更广泛的背景中。此原理由 \SourceCite{Givental3} 提出，并由 \SourceCite{Kim2} 给出正式表述；后者解释了广义旗流形中完全交的量子超平面截面原理。这一原理与把线性 $\sigma$ 模型的\IndexLocation{indexloc:0183:387}\IndexLocation{indexloc:0357:387}关联函数限制到反典范超曲面的量子限制公式 \SourceCite{MP1} 精神相近，但适用范围更广。

我们将在一个尚未获证的一般情形中陈述此原理。设 $Y$ 为光滑代数簇，$X\subset Y$ 为完全交，即向量丛 $\mathcal V=\bigoplus_{i=1}^{\ell}\mathcal L_i$ 的一个截面的零点集，其中 $\mathcal L_i$ 是 $Y$ 上的凸线丛。这里“凸”的含义与\hyperref[subsec:hypergeometric-mirror-transformations]{第~\ref*{subsec:hypergeometric-mirror-transformations} 小节}相同。因此，对所有无标记的亏格零\IndexLocation{indexloc:1280:387}稳定映射 $f:C\to Y$，均有 $H^1(C,f^*\mathcal L_i)=0$。

给定 $\beta\in H_2(Y,\mathbb Z)$，将\hyperref[subsec:givental-nef-projective-mirror]{第~\ref*{subsec:givental-nef-projective-mirror} 小节}的一些构造推广如下。特别地，令 $\mathcal V_{\beta,k}$ 为 $\overline M_{0,k}(Y,\beta)$ 上的向量丛，其在 $f$ 处的纤维为截面空间 $H^0(C,f^*\mathcal V)$。设 $\pi_k:\overline M_{0,k}(Y,\beta)\to\overline M_{0,k-1}(Y,\beta)$ 是忘掉第 $k$ 个标记点的映射，则 $\pi_k^*\mathcal V_{\beta,k-1}=\mathcal V_{\beta,k}$。对于 $1\le i\le k$，$H^0(C,f^*\mathcal V)$ 中在第 $i$ 个标记点消失的截面构成 $\overline M_{0,k}(Y,\beta)$ 上一个丛 $\mathcal V'_{\beta,k,i}$ 的纤维。这个丛有如下正合列
\[
0\longrightarrow\mathcal V'_{\beta,k,i}\longrightarrow\mathcal V_{\beta,k}\longrightarrow e_i^*\mathcal V\longrightarrow0.
\]
\SourcePage{388}
为简便起见，假设 $Y$ 为凸簇。现在定义取值于 $H^*(Y)$ 的形式函数 $J_{\mathcal V}$ 如下：
\begin{equation}\label{eq:givental-general-j-function}
J_{\mathcal V}=e^{(t_0+\delta)/\hbar}\operatorname{Euler}(\mathcal V)
\left(1+\sum_{\beta\ne0}q^\beta e_{1!}\left(\frac{\operatorname{Euler}(\mathcal V'_{\beta,2,1})}{\hbar-c}\right)\right).
\end{equation}
这里 $e_1$ 如常为赋值映射，$1/(\hbar-c)$ 是式~\eqref{eq:givental-projective-bundle-j-symbolic} 中所用的符号记法。采用约定 $\delta=\sum_{i=1}^r t_iT_i$，其中 $T_i$ 遍历 $H^2(Y)$ 的一组基，并令 $q^\beta=e^{\int_\beta\delta}$。当 $Y$ 为凸簇 $\mathbb P^n$ 时，这便化为式~\eqref{eq:givental-projective-bundle-j-symbolic}。此外，当 $Y$ 为一般的凸簇且 $\mathcal V=0$ 时，不难证明 $J_{\mathcal V}=J_Y$。

令 $\mathcal L_j(\beta)=\int_\beta c_1(\mathcal L_j)$，并定义另一个上同调值形式函数
\[
\begin{aligned}
\overline I_{\mathcal V}={}&e^{(t_0+\delta)/\hbar}\operatorname{Euler}(\mathcal V)\times\\
&\left(1+\sum_{0\ne\beta\in M(Y)}q^\beta
\prod_{j=1}^{\ell}\frac{\prod_{m=-\infty}^{\mathcal L_j(\beta)}(c_1(\mathcal L_j)+m\hbar)}{\prod_{m=-\infty}^0(c_1(\mathcal L_j)+m\hbar)}
 e_{1!}\left(\frac1{\hbar-c}\right)\right).
\end{aligned}
\]
如\hyperref[subsec:nef-gkz]{第~\ref*{subsec:nef-gkz} 小节}所述，乘积中除有限个因子外，其余全部约去。

现在陈述\IndexLocation{indexloc:1159:388}\IndexLocation{indexloc:1160:388}\IndexLocation{indexloc:1161:388}量子超平面截面原理。假设 $Y$ 的 Kähler 锥是由类 $T_i$ 生成的单纯锥。如常引入相应的变量 $t_i$ 与 $q_i=e^{t_i}$，用以下等式定义 $q_i$ 的次数：
\[
c_1(Y)-c_1(\mathcal V)=\sum_{i=1}^r(\deg q_i)T_i.
\]
注意，若 $Y=\mathbb P^n$、$\mathcal V=\bigoplus_{i=1}^{\ell}\mathcal O(a_i)$，则 $\deg q=n+1-\sum_{i=1}^{\ell}a_i$。
\begin{conjecture}\label{conj:quantum-hyperplane-section-principle}
如上，设 $X\subset Y$ 为与 $\mathcal V$ 相应的完全交。假设 $Y$ 为凸簇且所有 $i$ 均有 $\deg q_i\ge0$。则形式向量值函数 $\overline I_{\mathcal V}$ 与 $J_{\mathcal V}$ 在唯一的加权齐次三角变量变换后相等：
\[
t_0\longmapsto t_0+f_0(q)\hbar+h(q),\qquad t_i\longmapsto t_i+f_i(q),
\]
其中 $h,f_0,\ldots,f_r$ 是无常数项的加权齐次幂级数，且 $\deg f_0=\deg f_i=0$、$\deg h=1$。
\end{conjecture}
先考察此猜想在 $Y=\mathbb P^n$ 时的含义。从命题~\ref{prop:projective-space-i-equals-j} 开始，该命题断言 $I_{\mathbb P^n}=J_{\mathbb P^n}$。由式~\eqref{eq:givental-complete-intersection-i-function}，
\begin{equation}\label{eq:projective-space-i-hypergeometric}
I_{\mathbb P^n}=e^{(t_0+t_1H)/\hbar}\sum_{d=0}^{\infty}
\frac{q^d}{\prod_{m=1}^d(H+m\hbar)^{n+1}},
\end{equation}
其中 $q=e^{t_1}$。在式~\eqref{eq:givental-projective-bundle-j-symbolic} 后的说明中，已指出 $J_{\mathbb P^n}=J_{\mathcal V}$，这里 $\mathcal V=0$。于是式~\eqref{eq:givental-projective-bundle-j-symbolic} 给出
\begin{equation}\label{eq:projective-space-j-descendants}
J_{\mathbb P^n}=e^{(t_0+t_1H)/\hbar}\left(1+\sum_{d=1}^{\infty}e^{dt_1}e_{1!}\left(\frac1{\hbar-c}\right)\right).
\end{equation}
这是因为零秩平凡丛的 Euler 类为 $1$。比较两个表达式，得到
\begin{equation}\label{eq:projective-space-descendant-pushforward}
e_{1!}\left(\frac1{\hbar-c}\right)=\frac1{\prod_{m=1}^d(H+m\hbar)^{n+1}}.
\end{equation}
\SourcePage{389}
将它代入上述 $\overline I_{\mathcal V}$ 的公式，并注意此时 $\mathcal L_j(\beta)$ 就是 $a_jd$，便容易得出 $\overline I_{\mathcal V}$ 与式~\eqref{eq:givental-complete-intersection-i-function} 定义的 $I_{\mathcal V}$ 相同。

这表明：对于 $\mathbb P^n$，一旦证明命题~\ref{prop:projective-space-i-equals-j}，猜想~\ref{conj:quantum-hyperplane-section-principle} 陈述的量子超平面截面原理便等价于定理~\ref{thm:givental-projective-mirror} 的镜像定理。这是猜想~\ref{conj:quantum-hyperplane-section-principle} 发挥作用的典型方式：将式~\eqref{eq:projective-space-descendant-pushforward} 与量子超平面截面原理结合，就能得知 $\mathbb P^n$ 中完全交的镜像猜想应当是什么。更一般的情形也可作同样推理，因此，一旦知道凸簇 $Y$ 对应的式~\eqref{eq:projective-space-descendant-pushforward} 的右边，就立即知道 $Y$ 中完全交的相应镜像定理应当是什么。

作为猜想~\ref{conj:quantum-hyperplane-section-principle} 的证据，上面的讨论表明，证明定理~\ref{thm:givental-projective-mirror} 就能验证 $Y=\mathbb P^n$ 的情形。\SourceCite{Kim2} 给出了广义旗流形中完全交的证明；\SourceCite{BCFKvS1} 将其应用于 Grassmann 簇及完全旗流形中的 Calabi--Yau 完全交三维簇，计算有理曲线的虚数目。一般而言，量子超平面截面原理似乎对多种环境空间都给出合理结果。除 \SourceCite{Kim1,BCFKvS1} 研究的完全旗流形外，量子超平面截面原理还对某些二次曲面网簇产生了可以验证的结果（见 \SourceCite{Tjotta}）。

\subsection{命题~\ref{prop:projective-space-i-equals-j} 的证明}\label{subsec:givental-projective-space-proof}
现在概述命题~\ref{prop:projective-space-i-equals-j} 的证明。它断言
\[
I_{\mathbb P^n}=J_{\mathbb P^n},
\]
其中 $I_{\mathbb P^n}$ 与 $J_{\mathbb P^n}$ 分别由式~\eqref{eq:projective-space-i-hypergeometric} 和~\eqref{eq:projective-space-j-descendants} 给出。这里遵循 \SourceCite{Kim2} 的约定，其方法受到 \SourceCite{Givental4} 的启发。证明的关键步骤是引理~\ref{lem:projective-space-j-recursion} 给出的递推。\hyperref[subsec:givental-general-bundle-proof]{第~\ref*{subsec:givental-general-bundle-proof} 小节}将把这里的方法与一些重要思想结合，证明镜像定理（定理~\ref{thm:givental-projective-mirror}）。

策略是构造 $I_{\mathbb P^n}$ 和 $J_{\mathbb P^n}$ 的等变版本，再用类似于\hyperref[chap:localization]{第九章}和\hyperref[sec:quintic-mirror-theorem]{第~\ref*{sec:quintic-mirror-theorem} 节}的局部化技巧证明二者相等。采用例题~\ref{ex:projective-equivariant-localization} 所述 $T=(\mathbb C^*)^{n+1}$ 在 $\mathbb P^n$ 上的标准作用。它诱导 $\overline M_{0,2}(\mathbb P^n,d)$ 上的作用，使赋值映射 $e_1:\overline M_{0,2}(\mathbb P^n,d)\to\mathbb P^n$ 为等变映射。

先定义式~\eqref{eq:projective-space-j-descendants} 形式的 $J_{\mathbb P^n}$ 的等变版本：
\begin{equation}\label{eq:projective-space-equivariant-j}
J_T=e^{(t_0+t_1p)/\hbar}\left(1+\sum_{d=1}^{\infty}q^d e_{1!}\left(\frac1{\hbar-c_T}\right)\right),
\end{equation}
其中 $c_T$ 是等变 Chern 类 $c_1^T(\mathcal L_1)$ 的符号。如常设 $q=e^{t_1}$，$p$ 是 $\mathbb P^n$ 的\IndexLocation{indexloc:0250:389}等变超平面类。$J_T$ 取值于 $H_T^*(\mathbb P^n)$，取\IndexLocation{indexloc:0945:389}非等变极限 $\lambda_i\to0$、$p\to H$、$c_T\to c_1(\mathcal L_1)$，便能恢复 $J_{\mathbb P^n}$。更确切地，$J_{\mathbb P^n}=i_{\mathbb P^n}^*(J_T)$，其中 $i_{\mathbb P^n}$ 如\hyperref[sec:localization-theorem]{第~\ref*{sec:localization-theorem} 节}中定义。

去掉式~\eqref{eq:projective-space-equivariant-j} 的指数因子会更方便。令
\[
S=1+\sum_{d=1}^{\infty}q^d e_{1!}\left(\frac1{\hbar-c_T}\right),
\]
于是 $J_T=e^{(t_0+t_1p)/\hbar}S$。
\SourcePage{390}
接着把 $S$ 分解为较易研究的部分。令 $Z_i=\int_{(\mathbb P^n)_T}S\cup\phi_i$，其中 $\phi_i=\prod_{k\ne i}(p-\lambda_k)$ 是例题~\ref{ex:projective-equivariant-localization} 中 $H_T^*(\mathbb P^n)\otimes_{H^*(BT)}\mathcal R_T$ 的一组基。关键在于，$Z_i$ 足以确定 $S$，从而确定 $J_T$，最终也确定 $J_{\mathbb P^n}$：因为可以由 $S=\sum_{i=0}^nZ_i\cup\phi^i$ 恢复 $S$，其中 $\phi^i$ 是与 $\phi_i$ 对偶的基。

投影公式给出
\[
\int_{(\mathbb P^n)_T}e_{1!}\left(\frac1{\hbar-c_T}\right)\cup\phi_i
=\int_{\overline M_{0,2}(\mathbb P^n,d)_T}\frac{e_1^*(\phi_i)}{\hbar-c_T},
\]
所以可将 $Z_i$ 写为
\begin{equation}\label{eq:projective-space-j-fixed-components}
Z_i=1+\sum_{d=1}^{\infty}q^d\int_{\overline M_{0,2}(\mathbb P^n,d)_T}\frac{e_1^*(\phi_i)}{\hbar-c_T}.
\end{equation}
这里有一个简化。$\dim\overline M_{0,2}(\mathbb P^n,d)=(n+1)d+n-1$，因此其上按复次数计小于 $(n+1)d+n-1$ 的上同调类，积分均为零。另外，$\phi_i$ 的复次数为 $n$。将 $1/(\hbar-c_T)$ 按 $c_T/\hbar$ 的幂展开为几何级数，再代入式~\eqref{eq:projective-space-j-fixed-components}，可知含 $c_T^k$ 的项只有在 $k\ge(n+1)d-1$ 时才可能有贡献。因此，截去几何级数的低次项不会改变结果。出于稍后将明白的理由，我们截去的项略少于所能截去的项，从 $c_T^d$ 开始展开；这是可行的，因为 $d\le(n+1)d-1$。利用恒等式
\begin{equation}\label{eq:projective-space-descendant-truncation}
\sum_{k=d}^{\infty}\frac{c_T^k}{\hbar^{k+1}}=\left(\frac{c_T}{\hbar}\right)^d\frac1{\hbar-c_T},
\end{equation}
得到代替式~\eqref{eq:projective-space-j-fixed-components} 的另一形式：
\begin{equation}\label{eq:projective-space-j-truncated-components}
Z_i=Z_i(q,\hbar)=1+\sum_{d=1}^{\infty}\left(\frac q\hbar\right)^d
\int_{\overline M_{0,2}(\mathbb P^n,d)_T}\frac{e_1^*(\phi_i)c_T^d}{\hbar-c_T}.
\end{equation}
采用推论~\ref{cor:orbifold-localization-integration} 的形式，在 $\overline M_{0,2}(\mathbb P^n,d)$ 上局部化，计算 $Z_i$ 中出现的积分。每个积分
\[
\int_{\overline M_{0,2}(\mathbb P^n,d)_T}\frac{e_1^*(\phi_i)c_T^d}{\hbar-c_T}
\]
都会成为固定点集各分支上的积分之和；由\hyperref[subsec:kontsevich-localization]{第~\ref*{subsec:kontsevich-localization} 小节}，这些分支用某些带标记图 $\Gamma$ 描述。

设图 $\Gamma$ 对应于固定点分支 $\overline M_\Gamma\subset\overline M_{0,2}(\mathbb P^n,d)$。如\hyperref[subsec:kontsevich-localization]{第~\ref*{subsec:kontsevich-localization} 小节}，有包含映射 $i_\Gamma:\overline M_\Gamma\to\overline M_{0,2}(\mathbb P^n,d)$。首要任务是确定哪些分支在局部化后对上述积分有贡献。回顾 $T$ 在 $\mathbb P^n$ 上作用的固定点为 $q_j$，$0\le j\le n$。于是，除非 $e_1(\overline M_\Gamma)=q_i$，否则 $i_\Gamma^*(e_1^*(\phi_i))$ 为零。这是因为，若 $e_1(\overline M_\Gamma)=q_j$、$j\ne i$，则
\[
i_\Gamma^*(e_1^*(\phi_i))=i_{q_j}^*(\phi_i)=0,
\]
最后的等号由式~\eqref{eq:projective-dual-class-restriction} 得到。因此，局部化时只须对满足 $e_1(\overline M_\Gamma)=q_i$ 的 $\Gamma$ 求和。（注意，这与定理~\ref{thm:geometric-euler-data-gluing} 证明开头的做法相似。）
\SourcePage{391}
因为等变超平面类 $p$ 在 $q_i$ 上限制为 $\lambda_i$，式~\eqref{eq:projective-dual-class-restriction} 给出 $i_\Gamma^*(e_1^*\phi_i)=i_{q_i}^*(\phi_i)=\prod_{k\ne i}(\lambda_i-\lambda_k)$。也用符号 $c_\Gamma$ 表示 $c_T$ 在固定点分支上的限制。更准确地，定义
\[
\frac1{\hbar-c_\Gamma}=\sum_{k=0}^{\infty}\hbar^{-(k+1)}c_1^T(i_\Gamma^*\mathcal L_1)^k.
\]
最后，令 $a_\Gamma$ 为具有图 $\Gamma$ 的一般稳定映射的自同构群的阶。于是推论~\ref{cor:orbifold-localization-integration} 给出
\begin{equation}\label{eq:projective-space-j-localization}
Z_i(q,\hbar)=1+\sum_{d=1}^{\infty}\left(\frac q\hbar\right)^d\sum_\Gamma^{(i)}
\int_{(\overline M_\Gamma)_T}
\frac{c_\Gamma^d\prod_{k\ne i}(\lambda_i-\lambda_k)}{a_\Gamma(\hbar-c_\Gamma)\operatorname{Euler}_T(N_\Gamma)},
\end{equation}
其中 $\sum_\Gamma^{(i)}$ 表示仅对满足 $e_1(\overline M_\Gamma)=q_i$ 的 $\Gamma$ 求和。

还可以进一步减少需要考虑的分支。若 $(f,C,p_1,p_2)\in\overline M_\Gamma$，则 $e_1(\overline M_\Gamma)=q_i$ 意味着 $e_1(f)=f(p_1)=q_i$。令 $C'\subset C$ 为含 $p_1$ 的分支，区分以下两种情形：
\begin{itemize}
\item A 型：含 $p_1$ 的分支 $C'\subset C$ 被 $f$ 映到一条曲线上。
\item B 型：含 $p_1$ 的分支 $C'\subset C$ 被 $f$ 映到一个点上。
\end{itemize}
断言：B 型图 $\Gamma$ 对 $Z_i$ 没有贡献，因为此时 $c_\Gamma^d=0$。

为证明 B 型图的 $c_\Gamma^d=0$，须回顾\hyperref[subsec:kontsevich-localization]{第~\ref*{subsec:kontsevich-localization} 小节}对 $\overline M_\Gamma$ 的描述。图 $\Gamma$ 的顶点对应于 $f^{-1}(\{q_0,\ldots,q_n\})$ 的连通分支。因为 $\Gamma$ 为 B 型，$C'$ 包含在 $f^{-1}(q_i)$ 的某个连通分支 $C_v$ 内，$v$ 为相应顶点。如该小节，$n(v)$ 是 $C_v$ 中标记点的数目，加上将 $C_v$ 连到 $f$ 的次数 $d_e\ge1$ 的分支 $C_e$ 的结点数目。我们研究的是 $\overline M_{0,2}(\mathbb P^n,d)$，所以标记点至多有两个。另外，由图 $\Gamma$ 的定义，$\sum_ed_e=d$，所以这种结点至多有 $d$ 个。于是 $n(v)\le d+2$。

$n(v)$ 很重要，因为如该小节所解释，$\overline M_{0,n(v)}$ 是给出 $\overline M_\Gamma$ 的乘积中的一个因子。我们将研究 $c_\Gamma^d$ 在 $\overline M_{0,n(v)}$ 上的限制。因为 $n(v)\le d+2$，有 $\dim\overline M_{0,n(v)}=n(v)-3\le d-1$。此外，此时 $c_\Gamma$ 来自 $\overline M_{0,n(v)}$ 上相应丛 $\mathcal L_1$ 的等变 Chern 类 $c_1^T(\mathcal L_1)$。$T$ 在 $\overline M_{0,n(v)}$ 上的作用平凡，故 $H_T^*(\overline M_{0,n(v)})\simeq H^*(\overline M_{0,n(v)})\otimes H^*(BT)$。$T$ 在 $\mathcal L_1$ 上也平凡作用，因此在此同构下 $c_1^T(\mathcal L_1)=c_1(\mathcal L_1)\otimes1$。由于 $c_1(\mathcal L_1)$ 的实次数为 $2$，且 $\dim\overline M_{0,n(v)}\le d-1$，有 $c_1(\mathcal L_1)^d=0$。由此证明所断言的 $c_\Gamma^d=0$。

B 型图对积分没有贡献，所以用 $\sum_\Gamma^A$ 表示对 A 型图求和。仍假设 $e_1(\overline M_\Gamma)=q_i$，尽管记号中省略了下标 $i$。于是式~\eqref{eq:projective-space-j-localization} 成为
\begin{equation}\label{eq:projective-space-j-type-a-localization}
Z_i(q,\hbar)=1+\sum_{d=1}^{\infty}\sum_\Gamma^A\left(\frac q\hbar\right)^d
\int_{(\overline M_\Gamma)_T}\frac{c_\Gamma^d\prod_{k\ne i}(\lambda_i-\lambda_k)}{a_\Gamma(\hbar-c_\Gamma)\operatorname{Euler}_T(N_\Gamma)}.
\end{equation}
由此可证明以下引理。
\SourcePage{392}
\begin{lemma}[\IndexLocation{indexloc:0861:392}正则性引理]\label{lem:projective-space-j-regularity}
表达式 $Z_i$ 是环 $\mathbb Q(\lambda_0,\ldots,\lambda_n,\hbar)[[q]]$ 的元素。每个 $q^d$ 的系数都是 $\lambda_k$ 和 $\hbar$ 的有理函数，并且对于所有 $j\ne i$、$k\ge1$，在 $\hbar=(\lambda_i-\lambda_j)/k$ 处正则。
\end{lemma}
\begin{proof}
须考察式~\eqref{eq:projective-space-j-type-a-localization} 中 $q^d$ 的系数。为此，设 $\Gamma$ 为 A 型。分母中的 $\operatorname{Euler}_T(N_\Gamma)$ 属于 $H_T^*(\overline M_\Gamma)$，其局部化贡献关于 $\lambda_k$ 为有理函数，故只剩下分析 $c_\Gamma$。

回顾 $c_\Gamma=c_1^T(i_\Gamma^*\mathcal L_1)$。\hyperref[subsec:gravitational-correlator-definition]{第~\ref*{subsec:gravitational-correlator-definition} 小节}定义的 $\mathcal L_1$ 是 $\overline M_{0,2}(\mathbb P^n,d)$ 上的线丛，其在 $(f,C,p_1,p_2)$ 处的纤维为余切空间 $T_{p_1}^*C$。为计算 $c_1^T(i_\Gamma^*\mathcal L_1)$，须求 $(f,C,p_1,p_2)\in\overline M_\Gamma$ 时 $T$ 在 $T_{p_1}^*C$ 上作用的权。如前，令 $C'\subset C$ 为含 $p_1$ 的不可约分支。$\Gamma$ 为 A 型，因此 $C'$ 被映到直线 $\overline{q_iq_j}$ 上，次数为 $d'$。在 $q_i\in\overline{q_iq_j}\simeq\mathbb P^1$ 处，余切空间 $T_{q_i}^*\mathbb P^1$ 的权为 $\lambda_j-\lambda_i$；这可由例题~\ref{ex:projective-equivariant-localization} 的方法得出。因 $f|_{C'}$ 的次数为 $d'$，$T$ 在 $T_{p_1}^*C=T_{p_1}^*C'$ 上作用的权为 $(\lambda_j-\lambda_i)/d'$。所以 $c_\Gamma=(\lambda_j-\lambda_i)/d'$，从而
\[
\int_{(\overline M_\Gamma)_T}\frac{c_\Gamma^d\prod_{k\ne i}(\lambda_i-\lambda_k)}{a_\Gamma(\hbar-c_\Gamma)\operatorname{Euler}_T(N_\Gamma)}
=\int_{(\overline M_\Gamma)_T}\frac{((\lambda_j-\lambda_i)/d')^d\prod_{k\ne i}(\lambda_i-\lambda_k)}{a_\Gamma(\hbar-(\lambda_j-\lambda_i)/d')\operatorname{Euler}_T(N_\Gamma)}.
\]
由此容易验证所需结论。
\end{proof}
刚证明的引理表明，可以在 $Z_j$ 中代入 $\hbar=(\lambda_j-\lambda_i)/d$，得到有定义的表达式。这一点用于下面的递推引理。
\begin{lemma}[递推引理]\label{lem:projective-space-j-recursion}
对 $0\le i,j\le n$、$i\ne j$，令
\[
C_{i,j}(d)=\frac1{d!\displaystyle\prod_{\substack{1\le m\le d,\ k\ne i\\(k,m)\ne(j,d)}}(\lambda_i-\lambda_k+m(\lambda_j-\lambda_i)/d)}.
\]
则
\[
Z_i(q,\hbar)=1+\sum_{j\ne i}\sum_{d=1}^{\infty}\left(\frac q\hbar\right)^d
\frac{C_{i,j}(d)}{\lambda_i-\lambda_j+d\hbar}
Z_j\left(\frac q\hbar\frac{\lambda_j-\lambda_i}{d},\frac{\lambda_j-\lambda_i}{d}\right).
\]
\end{lemma}
\begin{proof}
从式~\eqref{eq:projective-space-j-type-a-localization} 推出此递推。如式~\eqref{eq:projective-space-j-fixed-components} 后的讨论，原本可把 $1/(\hbar-c_T)$ 的几何级数截到式~\eqref{eq:projective-space-descendant-truncation} 中从 $c_T^{(n+1)d-1}/\hbar^{(n+1)d}$ 开始。若这样做，会得到 \SourceCite{Givental2} 中的递推。这里的选择导出 \SourceCite{Kim2} 的递推关系，并避免了 \SourceCite{Givental2} 所用的 $q$ 的变量变换。

式~\eqref{eq:projective-space-j-type-a-localization} 中的图 $\Gamma$ 全为 A 型。对 $(f,C,p_1,p_2)\in\overline M_\Gamma$，有 $f(p_1)=q_i$。若 $C'\subset C$ 是含 $p_1$ 的分支，则存在 $j\ne i$，使 $f|_{C'}:C'\to\overline{q_iq_j}$ 为次数 $d'$ 的覆盖。若还有余下曲线，此时图 $\Gamma$ 的定义给出 $C=C'\cup C''$，其中 $C''$ 对应于图 $\Gamma''$，且 $p_2\in C''$。令 $p=C'\cap C''$，则 $q_j=f(p)$。把 $p$ 看作 $C''$ 上第一个标记点，则 $(f|_{C''},C'',p,p_2)$ 是
\SourcePage{393}
$\overline M_{0,2}(\mathbb P^n,d'')$ 中具有图 $\Gamma''$ 的二点稳定映射，其中 $d=d'+d''$。当 $C=C'$ 时，取 $d''=0$，相应贡献按常数项约定为 $1$。对于有余下曲线的图，断言
\begin{equation}\label{eq:projective-space-edge-removal-euler}
\begin{aligned}
a_\Gamma&=d'a_{\Gamma''},\\
c_\Gamma&=\frac{\lambda_j-\lambda_i}{d'},\\
\operatorname{Euler}_T(N_\Gamma)&=\operatorname{Euler}_T(N_{\Gamma''})
\left(\frac{\lambda_j-\lambda_i}{d'}-c_{\Gamma''}\right)
\prod_{m,k}^{(i)}\left(m\frac{\lambda_j-\lambda_i}{d'}+\lambda_i-\lambda_k\right),
\end{aligned}
\end{equation}
其中 $\prod_{m,k}^{(i)}$ 遍历 $0\le m\le d'-1$、$0\le k\le n$，但排除 $(m,k)=(0,i)$。

$a_\Gamma$ 的公式来自 $f|_{C'}$ 的自同构构成阶为 $d'$ 的循环群，正如式~\eqref{eq:fixed-graph-automorphisms} 及其后的讨论。$c_\Gamma$ 的公式已在引理~\ref{lem:projective-space-j-regularity} 的证明中得到。最后，对等变法丛 $N_\Gamma$ 作滤过，计算 $\operatorname{Euler}_T(N_\Gamma)$。因子 $\operatorname{Euler}_T(N_{\Gamma''})$ 来自 $f|_{C''}$ 的形变。因子
\[
\frac{\lambda_j-\lambda_i}{d'}-c_{\Gamma''}
\]
来自在结点 $p$ 处应用引理~\ref{lem:node-smoothing-tangent-tensor}，类似于定理~\ref{thm:geometric-euler-data-gluing} 的证明。最后的因子来自保持点 $p$ 不动的 $f|_{C'}$ 的形变。这些形变由向量场 $v_{m,k}=z_i^{d'-m}z_j^m\partial/\partial x_k$ 给出，其中 $0\le m<d'$。这里 $x_0,\ldots,x_n$ 是 $\mathbb P^n$ 上的坐标，$z_i,z_j$ 是 $C'$ 上的坐标，使 $f|_{C'}$ 由 $x_i=z_i^{d'}$、$x_j=z_j^{d'}$ 给出。$v_{m,k}$ 的权显然为
\[
(d'-m)\frac{\lambda_i}{d'}+m\frac{\lambda_j}{d'}-\lambda_k
=m\frac{\lambda_j-\lambda_i}{d'}+\lambda_i-\lambda_k.
\]
注意，向量场 $v_{d',k}$ 不保持 $p$ 不动，故没有贡献。另外，$v_{0,i}$ 除一个非零标量因子外，是 $C'$ 上保持 $p_1$ 和 $p$ 不动的向量场 $z_i\partial/\partial z_i$ 的像。它来自稳定映射的自同构，所以也须排除 $v_{0,i}$。把剩余各权相乘，便得到 $\operatorname{Euler}_T(N_\Gamma)$ 公式中的最后一个因子。

将式~\eqref{eq:projective-space-edge-removal-euler} 代入式~\eqref{eq:projective-space-j-type-a-localization}，计算得
\begin{equation}\label{eq:projective-space-j-recursion-factorization}
\begin{aligned}
Z_i=1+\sum_{j\ne i}\sum_{d'=1}^{\infty}
&\frac{(q/\hbar)^{d'}((\lambda_j-\lambda_i)/d')^{d'}\prod_{k\ne i}(\lambda_i-\lambda_k)}
{\left(\prod_{k\ne j}(\lambda_j-\lambda_k)\right)(\hbar-(\lambda_j-\lambda_i)/d')d'}\\
&\quad\times\frac{1}{\prod_{m,k}^{(i)}(m(\lambda_j-\lambda_i)/d'+\lambda_i-\lambda_k)}\\
&\quad\times\Biggl(1+\sum_{d''=1}^{\infty}\sum_{\Gamma''}^{(j)}\left(\frac q\hbar\right)^{d''}\\
&\qquad\quad\times\int_{(\overline M_{\Gamma''})_T}
\frac{((\lambda_j-\lambda_i)/d')^{d''}\prod_{k\ne j}(\lambda_j-\lambda_k)}
{\operatorname{Euler}_T(N_{\Gamma''})a_{\Gamma''}((\lambda_j-\lambda_i)/d'-c_{\Gamma''})}\Biggr).
\end{aligned}
\end{equation}
得到此式时，插入了相互约去的因子 $\prod_{k\ne j}(\lambda_j-\lambda_k)$，并使用 $d=d'+d''$。

现在从式~\eqref{eq:projective-space-j-recursion-factorization} 出发，注意以下各点，即可证明所需的 $Z_i$ 公式。此式最后一行正是
\[
Z_j\left(q(\lambda_j-\lambda_i)/(d'\hbar),(\lambda_j-\lambda_i)/d'\right),
\]
\SourcePage{394}
这可由式~\eqref{eq:projective-space-j-localization} 看出。至于第一行，其中的大分式含有因子
\[
\left(\frac q\hbar\right)^{d'}\frac1{\lambda_i-\lambda_j+d'\hbar}.
\]
所以，引理归结为证明剩余因子为 $C_{i,j}(d')$。也就是须证明
\begin{equation}\label{eq:projective-space-recursion-coefficient-cancellation}
C_{i,j}(d')=\frac{((\lambda_j-\lambda_i)/d')^{d'}\prod_{k\ne i}(\lambda_i-\lambda_k)}{\left(\prod_{k\ne j}(\lambda_j-\lambda_k)\right)\prod_{m,k}^{(i)}(m(\lambda_j-\lambda_i)/d'+\lambda_i-\lambda_k)}.
\end{equation}
为此，令 $A(m,k)=m(\lambda_j-\lambda_i)/d'+\lambda_i-\lambda_k$，注意三个涉及 $A(m,k)$ 的恒等式：
\[
\begin{aligned}
(d')!\left(\frac{\lambda_j-\lambda_i}{d'}\right)^{d'}&=\prod_{m=1}^{d'}A(m,i),\\
\lambda_i-\lambda_k&=A(0,k),\\
\lambda_j-\lambda_k&=A(d',k).
\end{aligned}
\]
于是可用 $A(m,k)$ 表示式~\eqref{eq:projective-space-recursion-coefficient-cancellation} 两边，并在分母中保留因子 $(d')!$。略去右边若干相互约去的 $A(m,k)$ 因子，就得到所需结果。
\end{proof}
初看时，引理~\ref{lem:projective-space-j-recursion} 的恒等式未必明显呈现为递推。将 $Z_i$ 写成
\[
Z_i(q,\hbar)=\sum_{d=0}^{\infty}Z_{i,d}(\hbar)q^d,
\]
其中 $Z_{i,0}(\hbar)=1$，递推便清楚了。对于 $d\ge1$，引理给出
\[
Z_{i,d}(\hbar)=\sum_{\substack{j\ne i,\ d=d'+d''\\d'\ge1}}
\frac{C_{i,j}(d')}{\lambda_i-\lambda_j+d'\hbar}
\left(\frac{\lambda_j-\lambda_i}{d'\hbar}\right)^{d''}\hbar^{-d'}
Z_{j,d''}\left(\frac{\lambda_j-\lambda_i}{d'}\right).
\]
因为 $d'\ge1$，$Z_{i,d}$ 由 $d''<d$ 的 $Z_{j,d''}$ 确定。再由引理给出 $Z_{i,0}=1$，可见 $Z_i$ 由引理中的恒等式唯一刻画。

有了\IndexLocation{indexloc:0860:394}递推引理，终于可以证明命题~\ref{prop:projective-space-i-equals-j}。
\begin{proof}[命题~\ref{prop:projective-space-i-equals-j} 的证明]
计划是构造 $I_{\mathbb P^n}$ 的等变版本 $I_T$，再证明它满足上述递推。回顾式~\eqref{eq:projective-space-i-hypergeometric}，
\[
I_{\mathbb P^n}=e^{(t_0+t_1H)/\hbar}\sum_d\frac{q^d}{((H+\hbar)\cdots(H+d\hbar))^{n+1}}.
\]
考虑等变表达式
\[
I_T=e^{(t_0+t_1p)/\hbar}\sum_{d=0}^{\infty}\frac{q^d}{\prod_{j=0}^n\prod_{m=1}^d(p-\lambda_j+m\hbar)},
\]
\SourcePage{395}
其非等变极限为 $I_{\mathbb P^n}$。目标是证明 $I_T=J_T$。先按对 $J_T$ 所作的同样方式分解 $I_T$。去掉指数因子，令
\[
S'=\sum_{d=0}^{\infty}\frac{q^d}{\prod_{j=0}^n\prod_{m=1}^d(p-\lambda_j+m\hbar)},
\]
再令 $Z'_i=\int_{(\mathbb P^n)_T}S'\cup\phi_i$。计算如下：把式~\eqref{eq:linear-model-restriction-integral} 的做法用于 $\mathbb P^n$，可知 $Z'_i=i_{q_i}^*(S')=S'(\lambda_i)$，从而
\[
Z'_i=1+\sum_{d=1}^{\infty}\frac{q^d}{\prod_{j=0}^n\prod_{m=1}^d(\lambda_i-\lambda_j+m\hbar)}.
\]
此式清楚表明 $Z'_i$ 属于 $\mathbb Q(\lambda_0,\ldots,\lambda_n,\hbar)[[q]]$，且满足引理~\ref{lem:projective-space-j-regularity} 的正则性。另外，直接计算可知 $Z'_i$ 满足引理~\ref{lem:projective-space-j-recursion} 的递推。由该引理证明后的讨论中指出的唯一性，所有 $i$ 均有 $Z_i=Z'_i$。为结束证明，回顾式~\eqref{eq:projective-space-equivariant-j} 后的讨论，
\[
J_T=e^{(t_0+t_1p)/\hbar}\sum_{i=0}^nZ_i\cup\phi^i.
\]
类似推理给出
\[
I_T=e^{(t_0+t_1p)/\hbar}\sum_{i=0}^nZ'_i\cup\phi^i,
\]
故 $I_T=J_T$。取非等变极限 $\lambda_i\to0$、$p\to H$，即得所需的 $I_{\mathbb P^n}=J_{\mathbb P^n}$。
\end{proof}

\subsection{定理~\ref{thm:givental-projective-mirror} 的证明}\label{subsec:givental-general-bundle-proof}
为证明\hyperref[subsec:givental-nef-projective-mirror]{第~\ref*{subsec:givental-nef-projective-mirror} 小节}陈述的镜像定理，将推广命题~\ref{prop:projective-space-i-equals-j} 证明中的各个步骤。不过，此时得到的递推（引理~\ref{lem:complete-intersection-j-recursion}）弱得多，所以还需要下面的引理~\ref{lem:complete-intersection-double-construction}、\ref{lem:complete-intersection-j-uniqueness} 和~\ref{lem:complete-intersection-p-transformations}。论述基于 \SourceCite{Givental2,Pandharipande3}，但也采用 \SourceCite{Kim2,Givental4} 的一些约定。这一讨论同时包含了射影空间中完全交的猜想~\ref{conj:quantum-hyperplane-section-principle} 的证明。

与命题~\ref{prop:projective-space-i-equals-j} 的证明相同，先给出 $J_{\mathcal V}$ 的等变版本。由式~\eqref{eq:givental-projective-bundle-j-symbolic} 中 $J_{\mathcal V}$ 的定义，令
\begin{equation}\label{eq:complete-intersection-equivariant-j}
J_T=e^{(t_0+t_1p)/\hbar}\operatorname{Euler}_T(\mathcal V)
\left(1+\sum_{d=1}^{\infty}q^d e_{1!}\left(\frac{\operatorname{Euler}_T(\mathcal V'_{d,2,1})}{\hbar-c_T}\right)\right).
\end{equation}
零秩平凡丛的 Euler 类为 $1$，所以当 $\mathcal V=0$ 时，此式化为式~\eqref{eq:projective-space-equivariant-j}。注意，$T$ 在 $\mathcal V=\bigoplus_{i=1}^{\ell}\mathcal O_{\mathbb P^n}(a_i)$ 上自然作用，故 $\operatorname{Euler}_T(\mathcal V'_{d,2,1})$ 有意义。如\hyperref[subsec:givental-projective-space-proof]{第~\ref*{subsec:givental-projective-space-proof} 小节}，去掉指数因子，令
\begin{equation}\label{eq:complete-intersection-equivariant-s}
S_{\mathcal V}=S_{\mathcal V}(q,\hbar)=\operatorname{Euler}_T(\mathcal V)
\left(1+\sum_{d=1}^{\infty}q^d e_{1!}\left(\frac{\operatorname{Euler}_T(\mathcal V'_{d,2,1})}{\hbar-c_T}\right)\right).
\end{equation}
也令
\[
\widetilde S_{\mathcal V}=1+\sum_{d>0}q^d e_{1!}\left(\frac{\operatorname{Euler}_T(\mathcal V'_{d,2,1})}{\hbar-c_T}\right),
\]
\SourcePage{396}
于是 $S_{\mathcal V}=\operatorname{Euler}_T(\mathcal V)\widetilde S_{\mathcal V}$。如例题~\ref{ex:projective-equivariant-localization} 所示，张量到 $\lambda_i$ 的有理函数域后，$p$ 可逆，故 $\operatorname{Euler}_T(\mathcal V)=(\prod_{i=1}^{\ell}a_i)p^\ell$ 也可逆。特别地，$S_{\mathcal V}$ 唯一确定 $\widetilde S_{\mathcal V}$，且
\[
\widetilde S_{\mathcal V}=(\operatorname{Euler}_T(\mathcal V))^{-1}S_{\mathcal V}.
\]
这里再次允许 $\lambda_i$ 的有理函数，使表达式有意义。本章余下部分将经常采用类似记号，用符号上的波浪号表示除以 $\operatorname{Euler}_T(\mathcal V)$。

最后，令
\[
Z_{i,\mathcal V}=\int_{(\mathbb P^n)_T}\widetilde S_{\mathcal V}\cup\phi_i=i_{q_i}^*(\widetilde S_{\mathcal V}).
\]
要研究 $Z_{i,\mathcal V}$ 是否与\hyperref[subsec:givental-projective-space-proof]{第~\ref*{subsec:givental-projective-space-proof} 小节}研究的函数 $Z_i$ 具有同样的性质。第一步很容易，因为可直接把式~\eqref{eq:projective-space-j-fixed-components} 推广为
\begin{equation}\label{eq:complete-intersection-j-fixed-components}
Z_{i,\mathcal V}=1+\sum_{d=1}^{\infty}q^d
\int_{\overline M_{0,2}(\mathbb P^n,d)_T}\frac{\operatorname{Euler}_T(\mathcal V'_{d,2,1})e_1^*(\phi_i)}{\hbar-c_T}.
\end{equation}
但由于 $\operatorname{Euler}_T(\mathcal V'_{d,2,1})$ 的出现，式~\eqref{eq:projective-space-descendant-truncation} 和~\eqref{eq:projective-space-j-truncated-components} 所作的截断一般不再适用。因此，不按式~\eqref{eq:projective-space-j-truncated-components} 到~\eqref{eq:projective-space-j-localization} 的方式局部化，而直接对式~\eqref{eq:complete-intersection-j-fixed-components} 中的积分应用局部化，容易得到
\begin{equation}\label{eq:complete-intersection-j-localization}
Z_{i,\mathcal V}=1+\sum_{d=1}^{\infty}q^d\sum_\Gamma^{(i)}
\int_{(\overline M_\Gamma)_T}\frac{i_\Gamma^*(\operatorname{Euler}_T(\mathcal V'_{d,2,1}))\prod_{k\ne i}(\lambda_i-\lambda_k)}{a_\Gamma(\hbar-c_\Gamma)\operatorname{Euler}_T(N_\Gamma)},
\end{equation}
其中 $\sum_\Gamma^{(i)}$ 与式~\eqref{eq:projective-space-j-localization} 中含义相同。

继续证明之前，暂停一下，对截断作一个有趣的观察。关键是，一般情形无法截断，这正是镜像定理需要变量变换的原因。例如，若 $\mathcal V=\bigoplus_{i=1}^{\ell}\mathcal O(a_i)$ 满足 $\sum_{i=1}^{\ell}a_i<n$，截断仍然可行。直接验证时用到 $\mathcal V'_{d,2,1}$ 的秩为 $d\sum_{i=1}^{\ell}a_i<dn$。所以 $d+\operatorname{rank}(\mathcal V'_{d,2,1})<d+nd$，从而 $d+\operatorname{rank}(\mathcal V'_{d,2,1})\le(n+1)d-1$，正好满足与式~\eqref{eq:projective-space-descendant-truncation} 前的论证类似的要求。此时，下面的证明显著简化，变得很像\hyperref[subsec:givental-projective-space-proof]{第~\ref*{subsec:givental-projective-space-proof} 小节}所作的证明。结果是，当 $\sum_{i=1}^{\ell}a_i<n$ 时，镜像定理断言 $I_{\mathcal V}=J_{\mathcal V}$。换言之，此时定理~\ref{thm:givental-projective-mirror} 中 $f_0=f_1=h=0$。例题~\ref{ex:fano-hypersurface-i-j-equality} 已给出这一情形的一个特例。一般地，当 $\sum_{i=1}^{\ell}a_i\le n+1$ 时，将看到无法截断如何使递推变复杂，并迫使我们改变变量。

回到镜像定理的证明，注意以下 \SourceCite{Pandharipande3} 的结果，它推广了引理~\ref{lem:projective-space-j-regularity}。
\begin{lemma}[\IndexLocation{indexloc:0861:396}正则性引理]\label{lem:complete-intersection-j-regularity}
表达式 $Z_{i,\mathcal V}$ 是环 $\mathbb Q(\lambda_0,\ldots,\lambda_n,\hbar)[[q]]$ 的元素。每个 $q^d$ 的系数都是 $\lambda_k$ 和 $\hbar$ 的有理函数，且对于所有 $j\ne i$、$k\ge1$，在 $\hbar=(\lambda_i-\lambda_j)/k$ 处正则。
\end{lemma}
\SourcePage{397}
\begin{proof}
对于 $\sum_\Gamma^{(i)}$ 中出现的图 $\Gamma$，须研究积分
\begin{equation}\label{eq:complete-intersection-regularity-local-integral}
\int_{(\overline M_\Gamma)_T}\frac{i_\Gamma^*(\operatorname{Euler}_T(\mathcal V'_{d,2,1}))\prod_{k\ne i}(\lambda_i-\lambda_k)}{a_\Gamma(\hbar-c_\Gamma)\operatorname{Euler}_T(N_\Gamma)}.
\end{equation}
\hyperref[subsec:givental-projective-space-proof]{第~\ref*{subsec:givental-projective-space-proof} 小节}把图 $\Gamma$ 分成 A 型和 B 型。对于 A 型图，引理~\ref{lem:projective-space-j-regularity} 的证明表明，存在 $d'\ge1$、$j\ne i$，使 $c_\Gamma=(\lambda_j-\lambda_i)/d'$。因此，这些分支的上述积分成为
\[
\int_{(\overline M_\Gamma)_T}\frac{i_\Gamma^*(\operatorname{Euler}_T(\mathcal V'_{d,2,1}))\prod_{k\ne i}(\lambda_i-\lambda_k)}{a_\Gamma(\hbar-(\lambda_j-\lambda_i)/d')\operatorname{Euler}_T(N_\Gamma)}.
\]
与引理~\ref{lem:projective-space-j-regularity} 相同，此表达式显然关于 $\lambda_k,\hbar$ 为有理函数，并对 $k\ge1$ 在 $\hbar=(\lambda_i-\lambda_j)/k$ 处正则。

再设 $\Gamma$ 为 B 型。在引理~\ref{lem:projective-space-j-regularity} 前的讨论中，已证明此时 $c_\Gamma^d=0$。于是，把式~\eqref{eq:complete-intersection-regularity-local-integral} 分母中的 $1/(\hbar-c_\Gamma)$ 展开为几何级数时，可忽略所有含 $c_\Gamma^k/\hbar^{k+1}$、$k\ge d$ 的项。因此，只须考虑 $k<d$ 的项。每项都是 $\lambda_i$ 的有理函数乘以 $\hbar^{-1}$ 的正次幂，所以关于 $\hbar$ 和 $\lambda_i$ 为有理函数，且满足所断言的正则性。
\end{proof}
如\hyperref[subsec:givental-projective-space-proof]{第~\ref*{subsec:givental-projective-space-proof} 小节}，此引理表明，可以在 $Z_{j,\mathcal V}$ 中代入 $\hbar=(\lambda_j-\lambda_i)/d$，得到有定义的表达式。
\begin{lemma}[递推引理]\label{lem:complete-intersection-j-recursion}
设 $\mathcal V=\bigoplus_{k=1}^{\ell}\mathcal O_{\mathbb P^n}(a_k)$。对于 $0\le i,j\le n$、$i\ne j$，令
\[
C_{i,j}^{\mathcal V}(d)=\frac{\prod_{k=1}^{\ell}\prod_{m=1}^{a_kd}(a_k\lambda_i+m(\lambda_j-\lambda_i)/d)}{d!\displaystyle\prod_{\substack{1\le m\le d,\ k\ne i\\(k,m)\ne(j,d)}}(\lambda_i-\lambda_k+m(\lambda_j-\lambda_i)/d)}.
\]
则差
\[
\begin{aligned}
R_i={}&Z_{i,\mathcal V}(q,\hbar)-\sum_{j\ne i}\sum_{d=1}^{\infty}\left(\frac q\hbar\right)^d
\frac{C_{i,j}^{\mathcal V}(d)}{\lambda_i-\lambda_j+d\hbar}\\
&\hspace{2em}\times Z_{j,\mathcal V}\left(\frac q\hbar\frac{\lambda_j-\lambda_i}{d},\frac{\lambda_j-\lambda_i}{d}\right)
\end{aligned}
\]
属于 $\mathcal R_T[\hbar^{-1}][[q]]$。
\end{lemma}
\begin{proof}
恒等式
\[
\frac1{\hbar-c_\Gamma}=\left(\frac{c_\Gamma}{\hbar}\right)^d\frac1{\hbar-c_\Gamma}
+\sum_{k=0}^{d-1}c_\Gamma^k\hbar^{-(k+1)}
\]
表明，式~\eqref{eq:complete-intersection-j-localization} 可改写为
\[
Z_{i,\mathcal V}=1+\sum_{d=1}^{\infty}\sum_\Gamma^{(i)}\left(\frac q\hbar\right)^d
\int_{(\overline M_\Gamma)_T}\frac{c_\Gamma^d i_\Gamma^*(\operatorname{Euler}_T(\mathcal V'_{d,2,1}))\prod_{k\ne i}(\lambda_i-\lambda_k)}{a_\Gamma(\hbar-c_\Gamma)\operatorname{Euler}_T(N_\Gamma)}
+\text{关于 $\hbar^{-1}$ 为多项式的项}.
\]
此外，引理~\ref{lem:complete-intersection-j-regularity} 的证明指出，B 型图有 $c_\Gamma^d=0$。故上式成为
\[
Z_{i,\mathcal V}=1+\sum_{d=1}^{\infty}\sum_\Gamma^A\left(\frac q\hbar\right)^d
\int_{(\overline M_\Gamma)_T}\frac{c_\Gamma^d i_\Gamma^*(\operatorname{Euler}_T(\mathcal V'_{d,2,1}))\prod_{k\ne i}(\lambda_i-\lambda_k)}{a_\Gamma(\hbar-c_\Gamma)\operatorname{Euler}_T(N_\Gamma)}
+\text{关于 $\hbar^{-1}$ 为多项式的项},
\]
\SourcePage{398}
其中如式~\eqref{eq:projective-space-j-type-a-localization}，$\sum_\Gamma^A$ 仅对满足 $e_1(\overline M_\Gamma)=q_i$ 的 A 型图 $\Gamma$ 求和。模去关于 $\hbar^{-1}$ 的多项式后，上式与式~\eqref{eq:projective-space-j-type-a-localization} 相同，只是分子多了因子 $i_\Gamma^*(\operatorname{Euler}_T(\mathcal V'_{d,2,1}))$。因此，按引理~\ref{lem:projective-space-j-recursion} 的论证，得到类似于式~\eqref{eq:projective-space-j-recursion-factorization} 的表达式，分子中额外出现若干因子。通过计算 $T$ 在丛 $\mathcal V'_{d,2,1}$ 上作用的权，可将这些因子与 $C_{i,j}^{\mathcal V}(d')$ 的分子等同；这里省略这一论证的细节。于是，$R_i$ 模去关于 $\hbar^{-1}$ 的多项式后为零，引理得证。
\end{proof}
将此引理与引理~\ref{lem:projective-space-j-recursion} 给出的 $\mathcal V=0$ 情形比较，可见该情形下 $R_i=1$。所以，对于零秩平凡丛，得到真正的递推：$Z_i$ 中 $q^d$ 的系数由各 $Z_j$ 中 $q^{d'}$、$d'<d$ 的系数决定。但当 $\mathcal V$ 任意时，只知道 $R_i$ 是 $\hbar^{-1}$ 的多项式，因而无法用 $\mathcal V=0$ 时的同样方式得到递推。为完成定理~\ref{thm:givental-projective-mirror} 的证明，还需要新内容。基本思想是，用下面引理~\ref{lem:complete-intersection-double-construction} 的“双重构造”补充引理~\ref{lem:complete-intersection-j-regularity} 和~\ref{lem:complete-intersection-j-recursion} 的正则性及递推条件，使这些性质足以保证唯一性。引理~\ref{lem:complete-intersection-j-uniqueness} 最终将实现这一点。
\begin{lemma}[\IndexLocation{indexloc:0850:398}双重构造引理]\label{lem:complete-intersection-double-construction}
设 $z$ 为变量，$\mathcal V=\bigoplus_{j=1}^{\ell}\mathcal O_{\mathbb P^n}(a_j)$。则表达式
\[
W(z,\hbar)=\sum_{i=0}^n\frac{E_i}{\prod_{j\ne i}(\lambda_i-\lambda_j)}e^{\lambda_i z}
Z_{i,\mathcal V}(qe^{z\hbar},\hbar)Z_{i,\mathcal V}(q,-\hbar),\qquad E_i=\prod_{j=1}^{\ell}(a_j\lambda_i),
\]
属于 $H^*(BT)[\hbar][[q,z]]$。
\end{lemma}
\begin{proof}
回顾\hyperref[subsec:stable-moduli-linear-models]{第~\ref*{subsec:stable-moduli-linear-models} 小节}讨论的映射 $\varphi:M_d\to N_d$ 及丛 $\mathcal V_d$，其中 $M_d=\overline M_{0,0}(\mathbb P^1\times\mathbb P^n,(1,d))$。这里采用此构造的一个变体：考虑 $\overline M_{0,2}(\mathbb P^1\times\mathbb P^n,(1,d))$，它具有\IndexLocation{indexloc:0276:398}赋值映射 $e_1,e_2$，取值于 $\mathbb P^1\times\mathbb P^n$。遵循 \SourceCite{Givental2}，定义
\[
L_d=e_1^{-1}(\{0\}\times\mathbb P^n)\cap e_2^{-1}(\{\infty\}\times\mathbb P^n).
\]
令 $i_{L_d}$ 为 $L_d$ 到 $\overline M_{0,2}(\mathbb P^1\times\mathbb P^n,(1,d))$ 的包含映射。映射 $\mu:L_d\to N_d$ 定义为 $\mu=\varphi\circ\pi_1\circ\pi_2\circ i_{L_d}$，其中 $\pi_k$ 是忘掉第 $k$ 个标记点的映射，$k=1,2$。\SourceCite{Givental2} 讨论了此映射，并将 $N_d$ 记为 $L'_d$。

还回顾由自然投影 $\mathbb P^1\times\mathbb P^n\to\mathbb P^n$ 诱导的映射 $\pi:M_d\to\overline M_{0,0}(\mathbb P^n,d)$。由它得到映射 $\widetilde\pi:L_d\to\overline M_{0,0}(\mathbb P^n,d)$，定义为 $\widetilde\pi=\pi\circ\pi_1\circ\pi_2\circ i_{L_d}$。对于 $d=0$，直接以 $\mathbb P^n$ 上的丛 $\mathcal V$ 定义相应零次项。令 $\widetilde\chi_d^{\mathcal V}=\operatorname{Euler}_G(\widetilde\pi^*\mathcal V_d)$。断言以下恒等式：
\begin{equation}\label{eq:complete-intersection-double-polynomiality-integral}
\sum_{d=0}^{\infty}q^d\int_{(L_d)_G}e^{\mu^*(\kappa)z}\widetilde\chi_d^{\mathcal V}
=\sum_{i=0}^n\frac{E_i}{\prod_{j\ne i}(\lambda_i-\lambda_j)}e^{\lambda_i z}
Z_{i,\mathcal V}(qe^{z\hbar},\hbar)Z_{i,\mathcal V}(q,-\hbar).
\end{equation}
一旦证明此式，就得到所需结论，因为 $\mu^*(\kappa)$ 与相应丛的 $G$ 等变 Euler 类属于等变上同调，因而关于 $\hbar$ 和 $\lambda_i$ 多项式地依赖。

将用推论~\ref{cor:orbifold-localization-integration} 的局部化公式证明此式。论证与定理~\ref{thm:geometric-euler-data-gluing} 的证明非常相似，因此仅概述计算，省略大部分细节。在 $\sum_ja_j=n+1$ 的情形，此式的完整证明见 \SourceCite[引理 7]{Pandharipande3}。
\SourcePage{399}
先考察式~\eqref{eq:complete-intersection-double-polynomiality-integral} 左边 $q^d$ 的系数
\begin{equation}\label{eq:complete-intersection-double-degree-integral}
\int_{(L_d)_G}e^{\mu^*(\kappa)z}\widetilde\chi_d^{\mathcal V}.
\end{equation}
第一步是确定 $L_d$ 的固定点分支。这些分支与定理~\ref{thm:geometric-euler-data-gluing} 证明中出现的固定点分支十分相似。

先固定 $0\le i\le n$，令 $q_i$ 为 $\mathbb P^n$ 中相应的 $T$ 固定点。对每个 $0<r\le d$，令 $F'_r$ 为 $\overline M_{0,2}(\mathbb P^n,r)^T$ 的一个分支，包含第一个标记点映到 $q_i$ 的稳定映射。若还有 $r<d$，选择 $\overline M_{0,2}(\mathbb P^n,d-r)^T$ 的一个分支 $F'_{d-r}$，其稳定映射的第一个标记点同样映到 $q_i$。将描述 $L_d^G$ 的一个同构于 $F'_r\times F'_{d-r}$ 的分支。略微滥用记号，仍用 $F'_r\times F'_{d-r}$ 表示它。一般映射 $f\in F'_r\times F'_{d-r}$ 的源曲线为 $C=C_0\cup C_1\cup C_2$，其中 $C_0=\mathbb P^1$，而 $f:C\to\mathbb P^1\times\mathbb P^n$ 给定为
\[
\begin{aligned}
f|_{C_0}(z)&=(z,q_i),\\
f|_{C_1}(z)&=(0,f_1(z)),\\
f|_{C_2}(z)&=(\infty,f_2(z)).
\end{aligned}
\]
曲线 $C_1$ 有两个标记点 $(p'_1,p_1)$，$f_1:(C_1,p'_1,p_1)\to\mathbb P^n$ 是 $F'_r$ 的元素。另外，$C_2$ 有两个标记点 $(p'_2,p_2)$，映射 $f_2:(C_2,p'_2,p_2)\to\mathbb P^n$ 属于 $F'_{d-r}$。最后，把 $0\in C_0$ 与 $p'_1\in C_1$ 粘合，把 $\infty\in C_0$ 与 $p'_2\in C_2$ 粘合。于是 $f:(C,p_1,p_2)\to\mathbb P^1\times\mathbb P^n$ 属于 $L_d$，所有这类映射构成的集合正是所要描述的分支 $F'_r\times F'_{d-r}$。$r=0$ 或 $r=d$ 时分支的描述更简单，留给读者。

接下来用局部化描述此分支对积分~\eqref{eq:complete-intersection-double-degree-integral} 的贡献。先限制被积式中的各项到 $F'_r\times F'_{d-r}$。注意，$\mu(F'_r\times F'_{d-r})$ 是 $N_d$ 的 $G$ 固定点 $p_{i,r}$，如定理~\ref{thm:geometric-euler-data-gluing} 的讨论。因此，由\hyperref[subsec:stable-moduli-linear-models]{第~\ref*{subsec:stable-moduli-linear-models} 小节}，$\mu^*(\kappa)$ 在 $F'_r\times F'_{d-r}$ 上限制为 $\lambda_i+r\hbar$。

再计算 $\widetilde\chi_d^{\mathcal V}$ 的限制。令 $\rho_1,\rho_2$ 为 $F'_r\times F'_{d-r}$ 到两个因子的投影，则有正合列
\[
\begin{aligned}
0\longrightarrow{}&\rho_1^*(\mathcal V'_{r,2,1}|_{F'_r})\oplus\rho_2^*(\mathcal V'_{d-r,2,1}|_{F'_{d-r}})\\
\longrightarrow{}&(\widetilde\pi^*\mathcal V_d)|_{F'_r\times F'_{d-r}}
\longrightarrow\mathcal V_{q_i}\otimes\mathcal O_{F'_r\times F'_{d-r}}\longrightarrow0.
\end{aligned}
\]
考察纤维就能说明此正合列的来源。在 $f\in F'_r\times F'_{d-r}$ 上，最后一个映射为限制映射
\[
H^0(C,f^*\mathcal V)\longrightarrow H^0(C_0,f|_{C_0}^*\mathcal V)=\mathcal V_{q_i}.
\]
其核由 $f_1^*\mathcal V$ 与 $f_2^*\mathcal V$ 的一对截面构成，二者分别在标记点 $p'_1$、$p'_2$ 消失。这些正是 $\rho_1^*(\mathcal V'_{r,2,1}|_{F'_r})$ 和 $\rho_2^*(\mathcal V'_{d-r,2,1}|_{F'_{d-r}})$ 的纤维。因此，正合列表明 $\widetilde\chi_d^{\mathcal V}$ 限制为
\[
E_i\operatorname{Euler}_G(\mathcal V'_{r,2,1}|_{F'_r})\operatorname{Euler}_G(\mathcal V'_{d-r,2,1}|_{F'_{d-r}}),
\]
因为 $\mathcal V_{q_i}=(\bigoplus_j\mathcal O_{\mathbb P^n}(a_j))_{q_i}$ 的各个权为 $a_j\lambda_i$，其 Euler 类为 $E_i=\prod_j(a_j\lambda_i)$。
\SourcePage{400}
最后，计算法丛的等变 Euler 类，结果为
\[
\begin{aligned}
\operatorname{Euler}_G(N(F'_r\times F'_{d-r}))={}&\prod_{j\ne i}(\lambda_i-\lambda_j)\\
&\times\frac{(\hbar-c_1^T(\mathcal L_{r,1}))\operatorname{Euler}_G(N(F'_r))}{\prod_{j\ne i}(\lambda_i-\lambda_j)}\\
&\times\frac{(-\hbar-c_1^T(\mathcal L_{d-r,1}))\operatorname{Euler}_G(N(F'_{d-r}))}{\prod_{j\ne i}(\lambda_i-\lambda_j)}.
\end{aligned}
\]
（暂不作显然的约分，其理由很快就会清楚。）此计算与定理~\ref{thm:geometric-euler-data-gluing} 中 $\operatorname{Euler}_G(N(F_r\times F_{d-r}))$ 的计算十分相似。唯一的区别是，后者有一个来自 $\mathbb P^1$ 上使 $0$ 与 $\infty$ 移动的自同构形变的因子 $-\hbar^2$。$\operatorname{Euler}_G(N(F'_r\times F'_{d-r}))$ 中没有此因子，因为按 $L_d$ 的定义，其中所有映射与投影到 $\mathbb P^1$ 因子的映射复合后，两个标记点的像始终为 $0$ 和 $\infty$。

汇总以上各项，可见 $F'_r\times F'_{d-r}$ 对积分~\eqref{eq:complete-intersection-double-degree-integral} 的贡献为
\[
\begin{aligned}
\frac{E_i}{\prod_{j\ne i}(\lambda_i-\lambda_j)}e^{(\lambda_i+r\hbar)z}
&\int_{(F'_r)_G}\frac{i_{F'_r}^*(\operatorname{Euler}_G(\mathcal V'_{r,2,1}))\prod_{j\ne i}(\lambda_i-\lambda_j)}{a_r(\hbar-c_{F'_r})\operatorname{Euler}_G(N(F'_r))}\\
\times{}&\int_{(F'_{d-r})_G}\frac{i_{F'_{d-r}}^*(\operatorname{Euler}_G(\mathcal V'_{d-r,2,1}))\prod_{j\ne i}(\lambda_i-\lambda_j)}{a_{d-r}(-\hbar-c_{F'_{d-r}})\operatorname{Euler}_G(N(F'_{d-r}))}.
\end{aligned}
\]
其中 $a_r,a_{d-r}$ 是相应有限自同构群的阶；$c_{F'_r},c_{F'_{d-r}}$ 分别为 $c_1^T(\mathcal L_{r,1})$ 和 $c_1^T(\mathcal L_{d-r,1})$ 在 $F'_r$、$F'_{d-r}$ 上的限制。

由推论~\ref{cor:orbifold-localization-integration}，积分~\eqref{eq:complete-intersection-double-degree-integral} 是上述贡献对所有 $F'_r$、$F'_{d-r}$ 求和，再对 $i=0,\ldots,n$ 求和的结果。随后乘以 $q^d$，再对所有 $d\ge0$ 求和，就得到式~\eqref{eq:complete-intersection-double-polynomiality-integral} 的左边。

为把它与该式右边联系起来，回顾式~\eqref{eq:complete-intersection-j-localization} 中 $Z_{i,\mathcal V}$ 的公式：
\[
Z_{i,\mathcal V}(q,\hbar)=1+\sum_{d=1}^{\infty}q^d\sum_\Gamma^{(i)}
\int_{(\overline M_\Gamma)_T}\frac{i_\Gamma^*(\operatorname{Euler}_T(\mathcal V'_{d,2,1}))\prod_{j\ne i}(\lambda_i-\lambda_j)}{a_\Gamma(\hbar-c_\Gamma)\operatorname{Euler}_T(N_\Gamma)}.
\]
如式~\eqref{eq:projective-space-j-localization} 所解释，$\sum_\Gamma^{(i)}$ 对 $\overline M_{0,2}(\mathbb P^n,d)$ 固定点集内、经 $e_1$ 映到 $q_i$ 的全部分支求和。将 $d$ 改记为 $r$，则上式可写成
\[
Z_{i,\mathcal V}(q,\hbar)=1+\sum_{r=1}^{\infty}q^r\sum_{F'_r}
\int_{(F'_r)_T}\frac{i_{F'_r}^*(\operatorname{Euler}_T(\mathcal V'_{r,2,1}))\prod_{j\ne i}(\lambda_i-\lambda_j)}{a_r(\hbar-c_{F'_r})\operatorname{Euler}_T(N(F'_r))}.
\]
此外，因 $\mathbb C^*$ 在 $\mathbb P^n$ 上平凡作用，可以将 $T$ 替换为 $G=\mathbb C^*\times T$。若用此 $Z_{i,\mathcal V}$ 公式展开式~\eqref{eq:complete-intersection-double-polynomiality-integral} 的右边，就容易看到结果恰是上一段所述的左边表达式。于是该式成立，引理也按前述论证得到证明。
\end{proof}
\SourcePage{401}
刚给出的 $Z_{i,\mathcal V}$ 公式表明，它自然可视为 $\hbar^{-1}$ 的形式幂级数。因此，引理~\ref{lem:complete-intersection-double-construction} 定义的 $W(z,\hbar)$ 竟是关于 $\hbar$ 的多项式，这颇为意外。将看到，这一点在唯一性引理（引理~\ref{lem:complete-intersection-j-uniqueness}）的证明中发挥重要作用。

接着将前三个引理所表述的 $Z_{i,\mathcal V}$ 的性质统一起来。考虑集合 $\{Z_i\}_{i=0}^n$，其中所有 $i$ 均有 $Z_i\in\mathbb Q(\lambda_0,\ldots,\lambda_n)[[\hbar^{-1},q]]$。遵循 \SourceCite{Pandharipande3}，作如下定义。
\begin{definition}\label{def:complete-intersection-class-p}\index{Givental 方法!$\mathcal P$ 类|hyperpage}
若集合 $\{Z_i\}_{i=0}^n$ 满足以下条件，就称它属于 $\mathcal P$ 类：
\begin{enumerate}[label=(\roman*)]
\item 所有 $i$ 均有 $Z_i(0,\hbar)=1$。
\item $Z_i$ 满足引理~\ref{lem:complete-intersection-j-regularity} 中 $Z_{i,\mathcal V}$ 所满足的正则性条件。
\item $Z_i$ 满足引理~\ref{lem:complete-intersection-j-recursion} 中 $Z_{i,\mathcal V}$ 所满足的递推条件。
\item 在引理~\ref{lem:complete-intersection-double-construction} 中用 $Z_i$ 代替 $Z_{i,\mathcal V}$ 定义 $W(z,\hbar)$，则 $W(z,\hbar)\in H^*(BT)[\hbar][[q,z]]$。
\end{enumerate}
\end{definition}
由此定义提及的三个引理，$\{Z_{i,\mathcal V}\}_{i=0}^n$ 属于 $\mathcal P$ 类。现在可陈述关键的唯一性引理。
\begin{lemma}[唯一性引理]\label{lem:complete-intersection-j-uniqueness}
设 $\{Z_i\}_{i=0}^n$ 与 $\{Z'_i\}_{i=0}^n$ 均属于 $\mathcal P$ 类。若所有 $0\le i\le n$ 的表达式按 $\hbar^{-1}$ 形式展开时，都有 $Z_i\equiv Z'_i\pmod{\hbar^{-2}}$，则所有 $i$ 均有 $Z_i=Z'_i$。
\end{lemma}
\begin{proof}
遵循 \SourceCite{Givental4,Kim2} 的证明。令 $Z_i=\sum_d Z_{i,d}q^d$，类似地写出 $Z'_i$。因两个集合均属于 $\mathcal P$ 类，有 $Z_{i,0}=Z'_{i,0}=1$。因此，可归纳假设：对于某个 $d_0>0$，所有 $0\le d<d_0$、$0\le i\le n$ 均有 $Z_{i,d}=Z'_{i,d}$。属于 $\mathcal P$ 类意味着引理~\ref{lem:complete-intersection-j-recursion} 的递推同时适用于两集合。结合归纳假设，对所有 $i$，$Z'_{i,d_0}-Z_{i,d_0}$ 可展开为 $\hbar^{-1}$ 的多项式。令
\[
D_i(\hbar)=Z'_{i,d_0}-Z_{i,d_0}=\hbar^{-2r_i}\bigl(A_i\hbar^{-1}+B_i+O(\hbar)\bigr),
\]
其中 $A_i,B_i\in\mathbb C(\lambda_0,\ldots,\lambda_n)$，$r_i\ge0$。由 $Z_i\equiv Z'_i\pmod{\hbar^{-2}}$，只要证明所有 $r_i$ 均可取为零，就可得 $Z'_{i,d_0}=Z_{i,d_0}$。

所以须证明：$r_i\ge1$ 时，最高阶的系数 $A_i,B_i$ 为零。为此，用引理~\ref{lem:complete-intersection-double-construction} 的公式，分别由 $\{Z_i\}$ 和 $\{Z'_i\}$ 构造 $W$、$W'$。归纳假设容易推出，$W'-W$ 中 $q^{d_0}$ 的系数为
\[
\sum_{i=0}^n\frac{E_i}{\prod_{j\ne i}(\lambda_i-\lambda_j)}e^{\lambda_i z}
\bigl(e^{d_0\hbar z}D_i(\hbar)+D_i(-\hbar)\bigr).
\]
将 $D_i$ 的上述公式代入并简化，可见它具有形式
\[
\sum_{i=0}^n\frac{E_i}{\prod_{j\ne i}(\lambda_i-\lambda_j)}e^{\lambda_i z}
\hbar^{-2r_i}\bigl(A_id_0z+2B_i+O(\hbar)\bigr).
\]
因为两个集合均属于 $\mathcal P$ 类，$W$、$W'$ 关于 $\hbar$ 都是多项式。取最大的 $r_i=r\ge1$，上述表达式中 $\hbar^{-2r}$ 的系数必须为零。对于互异的 $\lambda_i$，，而 $E_i/\prod_{j\ne i}(\lambda_i-\lambda_j)$ 在局部化后非零。因此，所有满足 $r_i=r$ 的 $A_i,B_i$ 均为零。逐次降低最高负幂，得到所有 $r_i\ge1$ 的系数消失，引理得证。
\end{proof}
证明\IndexLocation{indexloc:0847:401}\IndexLocation{indexloc:0850:401}\IndexLocation{indexloc:0860:401}\IndexLocation{indexloc:0861:401}\IndexLocation{indexloc:0863:401}镜像定理之前，还需两个准备步骤。首先，须给出现的变量赋予权重：赋予
\SourcePage{402}
$\lambda_i,p,\hbar$ 各权重 $1$，赋予 $q$ 权重 $n+1-\sum_i a_i$。下面“齐次”一词均相对于这些权重而言。以下引理表明，某些坐标变换不改变属于 $\mathcal P$ 类的性质。
\begin{lemma}\label{lem:complete-intersection-p-transformations}
考虑形式幂级数 $f(q)\in q\mathbb Q[[q]]$、$g(q),h(q)\in qH^*(BT)[[q]]$，均为齐次，且 $\deg f=\deg g=0$、$\deg h=1$。若 $\{Z_i\}_{i=0}^n$ 属于 $\mathcal P$ 类，则 $\{Z''_i\}_{i=0}^n$ 亦然，其中
\[
Z''_i(q,\hbar)=e^{f(q)+h(q)/\hbar+g(q)\lambda_i/\hbar}Z_i(qe^{g(q)},\hbar).
\]
\end{lemma}
\begin{proof}
当 $\sum_ja_j=n+1$ 时，详尽证明见 \SourceCite[引理 10]{Pandharipande3}。该论证容易推广到 $\sum_ja_j\le n+1$ 的情形，因此省略细节。
\end{proof}
终于可以概述 Givental 关于 $\mathbb P^n$ 中 nef 完全交的镜像定理的证明。
\begin{proof}[定理~\ref{thm:givental-projective-mirror} 的证明]
须证明 $I_{\mathcal V}$ 与 $J_{\mathcal V}$ 在适当的齐次变量变换后相等。与 $\mathcal V=0$ 的情形一样，策略是构造 $I_{\mathcal V}$ 与 $J_{\mathcal V}$ 的等变版本，随后证明定理的等变版本。

$J_{\mathcal V}$ 的等变版本是式~\eqref{eq:complete-intersection-equivariant-j} 中的 $J_T$。去掉指数因子得到式~\eqref{eq:complete-intersection-equivariant-s} 中的 $S_{\mathcal V}$，再去掉 $\operatorname{Euler}_T(\mathcal V)$ 得到 $\widetilde S_{\mathcal V}$。最后，由公式
\[
Z_{i,\mathcal V}=\int_{(\mathbb P^n)_T}\widetilde S_{\mathcal V}\cup\phi_i=i_{q_i}^*(\widetilde S_{\mathcal V})
\]
得到 $Z_{i,\mathcal V}$。知道 $J_T$ 由这些函数唯一确定，且由引理~\ref{lem:complete-intersection-j-regularity}、\ref{lem:complete-intersection-j-recursion} 和~\ref{lem:complete-intersection-double-construction}，集合 $\{Z_{i,\mathcal V}\}_{i=0}^n$ 属于 $\mathcal P$ 类。还注意，上述权重下 $\widetilde S_{\mathcal V}$ 是权重为零的齐次函数。这可由式~\eqref{eq:complete-intersection-equivariant-s} 及三个简单事实看出：丛 $\mathcal V'_{d,2,1}$ 的秩为 $\sum_{j=1}^{\ell}a_jd$；$e_1:\overline M_{0,2}(\mathbb P^n,d)\to\mathbb P^n$ 的纤维维数为 $(n+1)d-1$；$1/(\hbar-c_T)$ 展开中的一般项为 $c_T^k/\hbar^{k+1}$。细节留给读者。

转而考虑 $I_{\mathcal V}$ 的等变版本，定义
\begin{equation}\label{eq:complete-intersection-equivariant-i-series}
\widetilde S'_{\mathcal V}(q,\hbar)=1+\sum_{d=1}^{\infty}q^d
\frac{\prod_{j=1}^{\ell}\prod_{m=1}^{a_jd}(a_jp+m\hbar)}{\prod_{j=0}^n\prod_{m=1}^d(p-\lambda_j+m\hbar)},
\end{equation}
再令
\[
I_T=e^{(t_0+t_1p)/\hbar}\operatorname{Euler}_T(\mathcal V)\widetilde S'_{\mathcal V}(q,\hbar).
\]
当 $\lambda_i\to0$ 时，$I_T$ 的非等变极限显然为 $I_{\mathcal V}$。$\widetilde S'_{\mathcal V}$ 也是权重为零的齐次函数。最后定义
\[
Z'_{i,\mathcal V}=\int_{(\mathbb P^n)_T}\widetilde S'_{\mathcal V}\cup\phi_i=i_{q_i}^*(\widetilde S'_{\mathcal V}).
\]
可验证 $\{Z'_{i,\mathcal V}\}_{i=0}^n$ 属于 $\mathcal P$ 类。这一验证直接但略长，省略细节。

为比较 $I_T$ 与 $J_T$，先注意
\[
J_T=e^{(t_0+t_1p)/\hbar}\operatorname{Euler}_T(\mathcal V)\bigl(1+o(\hbar^{-1})\bigr).
\]
\SourcePage{403}
这与证明式~\eqref{eq:givental-bundle-j-normalization} 时的论证相同，该式给出渐近展开 $J_{\mathcal V}=e^{(t_0+t_1H)/\hbar}\operatorname{Euler}(\mathcal V)(1+o(\hbar^{-1}))$。

现在考虑变量变换 $t_0\mapsto t_0+f(q)\hbar+h(q)$、$t_1\mapsto t_1+g(q)$，令 $J''_T$ 为在 $J_T$ 中作此代换后的表达式。上述渐近展开给出
\[
J''_T=e^{(t_0+t_1p)/\hbar}\operatorname{Euler}_T(\mathcal V)
\bigl(e^{f(q)}+e^{f(q)}h(q)\hbar^{-1}+e^{f(q)}g(q)p\hbar^{-1}+o(\hbar^{-1})\bigr).
\]
以同样方式展开 $I_T$，可见存在唯一一组 $f,g,h$，使 $I_T\equiv J''_T\pmod{\hbar^{-2}}$；这类似于式~\eqref{eq:givental-mirror-variable-coefficients} 中所作的计算。注意，$f,g$ 的次数为零，$h$ 的次数为一，因为 $\widetilde S_{\mathcal V}$ 与 $\widetilde S'_{\mathcal V}$ 的次数均为零。

为证明定理，须证明 $I_T=J''_T$。首先注意，$t_1\mapsto t_1+g(q)$ 将 $q=e^{t_1}$ 变为 $qe^{g(q)}$。因此，在
\[
J_T=e^{(t_0+t_1p)/\hbar}\operatorname{Euler}_T(\mathcal V)\widetilde S_{\mathcal V}(q,\hbar)
\]
中作变量变换，得到
\[
\begin{aligned}
J''_T&=e^{(t_0+f(q)\hbar+h(q)+(t_1+g(q))p)/\hbar}\operatorname{Euler}_T(\mathcal V)\widetilde S_{\mathcal V}(qe^{g(q)},\hbar)\\
&=e^{(t_0+t_1p)/\hbar}\operatorname{Euler}_T(\mathcal V)e^{f(q)+h(q)/\hbar+g(q)p/\hbar}\widetilde S_{\mathcal V}(qe^{g(q)},\hbar)\\
&=e^{(t_0+t_1p)/\hbar}\operatorname{Euler}_T(\mathcal V)\widetilde S''_{\mathcal V}(q,\hbar),
\end{aligned}
\]
其中
\[
\widetilde S''_{\mathcal V}(q,\hbar)=e^{f(q)+h(q)/\hbar+g(q)p/\hbar}\widetilde S_{\mathcal V}(qe^{g(q)},\hbar).
\]
若令
\[
Z''_{i,\mathcal V}=\int_{(\mathbb P^n)_T}\widetilde S''_{\mathcal V}\cup\phi_i=i_{q_i}^*(\widetilde S''_{\mathcal V}),
\]
则 $I_T=J''_T$ 等价于所有 $i$ 均有 $Z'_{i,\mathcal V}=Z''_{i,\mathcal V}$。

再断言集合 $\{Z''_{i,\mathcal V}\}_{i=0}^n$ 属于 $\mathcal P$ 类。为此，注意上述 $\widetilde S''_{\mathcal V}$、$Z''_{i,\mathcal V}$ 的公式给出
\[
Z''_{i,\mathcal V}(q,\hbar)=e^{f(q)+h(q)/\hbar+g(q)\lambda_i/\hbar}Z_{i,\mathcal V}(qe^{g(q)},\hbar),
\]
因为 $i_{q_i}^*(p)=\lambda_i$。由引理~\ref{lem:complete-intersection-p-transformations}，$\{Z''_{i,\mathcal V}\}_{i=0}^n$ 属于 $\mathcal P$ 类，因为 $\{Z_{i,\mathcal V}\}_{i=0}^n$ 属于此类。

现在定理容易得证。已知 $\{Z'_{i,\mathcal V}\}$ 与 $\{Z''_{i,\mathcal V}\}$ 均属于 $\mathcal P$ 类，而 $I_T\equiv J''_T\pmod{\hbar^{-2}}$ 意味着所有 $i$ 均有 $Z'_{i,\mathcal V}\equiv Z''_{i,\mathcal V}\pmod{\hbar^{-2}}$。因此，唯一性引理（引理~\ref{lem:complete-intersection-j-uniqueness}）的假设成立，可得所有 $i$ 均有 $Z'_{i,\mathcal V}=Z''_{i,\mathcal V}$。如前所述，这意味着 $I_T=J''_T$。取非等变极限 $\lambda_i\to0$，表明所需的变量变换后 $I_{\mathcal V}$ 等于 $J_{\mathcal V}$。镜像定理的证明完成。
\end{proof}
对本节与\hyperref[sec:quintic-mirror-theorem]{第~\ref*{sec:quintic-mirror-theorem} 节}所讨论的镜像定理证明方法之间的关系，再作几点评论。引理~\ref{lem:complete-intersection-j-recursion} 的递推关系，与\hyperref[subsec:lly-euler-data]{第~\ref*{subsec:lly-euler-data} 小节}中引理~\ref{lem:euler-data-gluing} 的粘合恒等式起着类似作用。把递推关系与引理~\ref{lem:complete-intersection-double-construction} 的双重构造的多项式性质结合，对应于该小节讨论的 Euler 数据性质。此外，引理~\ref{lem:complete-intersection-j-uniqueness} 的唯一性对应于\hyperref[subsec:linked-euler-data]{第~\ref*{subsec:linked-euler-data} 小节}的定理~\ref{thm:linked-euler-data-uniqueness}。注意，该定理次数界 $5d-2$ 中的 $-2$，
\SourcePage{404}
对应于这里唯一性引理证明中两个系数（$1$、$\hbar^{-1}$ 的系数）的匹配。最后，两种方法都使用镜像变换，将两个不同的形式级数等同。\SourceCite{Givental5} 的一个长篇脚注更详细地比较了这两种方法，该文也给出了亏格一的镜像定理。

\subsection{环面完全交}\label{subsec:givental-toric-mirror}\index{镜像定理!环面完全交|hyperpage}
接下来讨论镜像定理到光滑环面簇中完全交的推广 \SourceCite{Givental4}。主要结果将在下面作为定理~\ref{thm:givental-toric-mirror} 重述。随后讨论环面镜像定理与本书一个主要主题的关系，即 Batyrev 镜像构造所给出的 Calabi--Yau 流形的镜像对称。最后给出几个有趣例子。

设 $X_\Sigma$ 为扇 $\Sigma$ 对应的光滑环面簇，$\mathcal L_1,\ldots,\mathcal L_\ell$ 是其上由全局截面生成的线丛。令 $\mathcal V=\bigoplus_{i=1}^{\ell}\mathcal L_i$，$X\subset X_\Sigma$ 为 $\mathcal V$ 的一般全局截面的零点集所定义的光滑完全交。在此情形，要定义上同调值函数 $I_{\mathcal V}$ 和 $J_{\mathcal V}$，推广\hyperref[subsec:givental-nef-projective-mirror]{第~\ref*{subsec:givental-nef-projective-mirror} 小节}的构造。

先约定一些记号。对每个 $\rho\in\Sigma(1)$，略微滥用记号，也用 $D_\rho$ 表示相应除子在 $H^2(X_\Sigma)$ 中的上同调类。遵循 \SourceCite{Givental4}，令 $\mathcal L_i(\beta)=\int_\beta c_1(\mathcal L_i)$、$D_\rho(\beta)=\int_\beta D_\rho$。还选取 $H^2(X_\Sigma,\mathbb Z)$ 的整基 $T_1,\ldots,T_r$，使各基元素位于 Kähler 锥的闭包内。如常令 $\delta=\sum_{i=1}^r t_iT_i$。

定义两个上同调值形式函数。先定义 $I_{\mathcal V}$：
\begin{equation}\label{eq:toric-complete-intersection-i-function}
\begin{aligned}
I_{\mathcal V}={}&e^{(t_0+\delta)/\hbar}\operatorname{Euler}(\mathcal V)\\
&\times\sum_{\beta\in M(X_\Sigma)}q^\beta
\frac{\prod_{i=1}^{\ell}\prod_{m=-\infty}^{\mathcal L_i(\beta)}(c_1(\mathcal L_i)+m\hbar)\prod_\rho\prod_{m=-\infty}^0(D_\rho+m\hbar)}{\prod_{i=1}^{\ell}\prod_{m=-\infty}^0(c_1(\mathcal L_i)+m\hbar)\prod_\rho\prod_{m=-\infty}^{D_\rho(\beta)}(D_\rho+m\hbar)}.
\end{aligned}
\end{equation}
其中 $q_i=e^{t_i}$、$q^\beta=\prod_{i=1}^r q_i^{\int_\beta T_i}$。注意，若 $\Sigma$ 是 $\mathbb P^n$ 的标准扇，就恢复式~\eqref{eq:givental-complete-intersection-i-function}。再定义 $J_{\mathcal V}$：
\begin{equation}\label{eq:toric-complete-intersection-j-function}
\begin{aligned}
J_{\mathcal V}={}&e^{(t_0+\delta)/\hbar}\operatorname{Euler}(\mathcal V)\\
&\times\left(1+\sum_{\beta\ne0}q^\beta PD^{-1}e_{1*}
\left(\frac{\operatorname{Euler}(\mathcal V'_{\beta,2,1})}{\hbar-c}\cap[\overline M_{0,2}(X_\Sigma,\beta)]^{\mathrm{virt}}\right)\right),
\end{aligned}
\end{equation}
其中 $[\overline M_{0,2}(X_\Sigma,\beta)]^{\mathrm{virt}}$ 是 $\overline M_{0,2}(X_\Sigma,\beta)$ 的虚拟基本类，$PD$ 是 Poincaré 对偶。注意，当 $X_\Sigma$ 为凸簇时，虚拟基本类就是通常的基本类，所以 $J_{\mathcal V}$ 的公式可简化。例如，当 $X_\Sigma$ 为凸簇 $\mathbb P^n$ 时，式~\eqref{eq:toric-complete-intersection-j-function} 化为式~\eqref{eq:givental-general-j-function}。

此时，变量 $q_i$ 有次数。如\hyperref[subsec:quantum-hyperplane-section-principle]{第~\ref*{subsec:quantum-hyperplane-section-principle} 小节}，用等式
\[
c_1(X_\Sigma)-c_1(\mathcal V)=\sum_{i=1}^r(\deg q_i)T_i
\]
定义 $\deg q_i$。假设 $X\subset X_\Sigma$ 是\hyperref[subsec:nef-gkz]{第~\ref*{subsec:nef-gkz} 小节}意义下的 nef 完全交，即 $-(K_{X_\Sigma}+\sum_{i=1}^{\ell}\mathcal L_i)$ 在 $X_\Sigma$ 上为 nef。此时假设基 $T_1,\ldots,T_r$ 已选取为使
\SourcePage{405}
$-(K_{X_\Sigma}+\sum_{i=1}^{\ell}\mathcal L_i)$ 位于 $T_i$ 所生成的锥中。在 nef 情形总可作这种选择。于是所有 $i$ 均有 $\deg q_i\ge0$。

现在可陈述 Givental 版本的环面镜像定理。
\begin{theorem}\label{thm:givental-toric-mirror}
设 $X\subset X_\Sigma$ 是 nef 完全交，$I_{\mathcal V}$、$J_{\mathcal V}$ 分别由式~\eqref{eq:toric-complete-intersection-i-function} 和~\eqref{eq:toric-complete-intersection-j-function} 定义。则二者在加权齐次三角变量变换后相等：
\[
t_0\longmapsto t_0+f_0(q)\hbar+h(q),\qquad t_i\longmapsto t_i+f_i(q),\quad1\le i\le r,
\]
其中 $f_0,\ldots,f_r,h$ 为加权齐次幂级数，且 $\deg f_0=\deg f_i=0$、$\deg h=1$。
\end{theorem}
\begin{proof}
如可预料，证明的策略是构造 $I_{\mathcal V}$ 和 $J_{\mathcal V}$ 的等变版本。这正是环面簇的重要性所在，因为 $T\subset X_\Sigma$ 提供了所需的群作用。不过，式~\eqref{eq:toric-complete-intersection-j-function} 表明，还需要虚拟基本类的等变版本。局部化在证明中起着如此核心的作用，因此等变虚拟基本类的局部化公式不可缺少。这便联系到\hyperref[subsec:virtual-localization]{第~\ref*{subsec:virtual-localization} 小节}讨论的一些问题。

这里不证明此定理。\SourceCite{Givental4} 的证明在写作时尚不完整。特别地，如上所述，证明以存在适当的等变虚拟基本类局部化公式为前提。\SourceCite{GPa} 已建立了这样的公式。考虑到此定理的重要性，我们希望其余证明细节能被仔细写出。我们对 Givental 建立的概念框架完全有信心，相信这些细节将会出现，就如 \SourceCite{BDPP,Pandharipande3} 对命题~\ref{prop:projective-space-i-equals-j} 和定理~\ref{thm:givental-projective-mirror} 所做的一样。
\end{proof}
还可问：式~\eqref{eq:toric-complete-intersection-j-function} 定义的形式函数 $J_{\mathcal V}$，与\hyperref[subsec:giventals-j-function]{第~\ref*{subsec:giventals-j-function} 小节}的 Givental 函数 $J_X$ 有什么关系？这里 $i:X\hookrightarrow X_\Sigma$ 是 $\mathcal V$ 的一般截面的零点集的包含映射。基本思想是，$J_{\mathcal V}$ 应计算来自环境空间的上同调类所对应的亏格零\IndexLocation{indexloc:0459:405}引力关联数。因此，预期 $J_{\mathcal V}$ 包含 $J_X$ 的部分信息。有一个情形下二者关系很简单：当 $\mathcal V=0$ 时，$X=X_\Sigma$，不难证明 $J_X=J_{\mathcal V}$。当 $\mathcal V$ 非零时，情况更复杂，因为 $J_{\mathcal V}$ 取值于 $H^*(X_\Sigma)$，而 $J_X$ 取值于 $H^*(X)$。对于 $X_\Sigma=\mathbb P^n$，命题~\ref{prop:complete-intersection-j-pushforward} 已证明 $J_{\mathcal V}=i_!(J_X)$。更一般地，当 $X_\Sigma$ 为凸簇且 $\mathcal V$ 丰沛时，也成立。不过，尚不知道在此定理的情形下是否仍有 $J_{\mathcal V}=i_!(J_X)$。一个问题是，$H^2(X_\Sigma)\to H^2(X)$ 未必为满射，故 $J_X$ 可能涉及比 $J_{\mathcal V}$ 更多的变量。若令 $H^*(\mathcal V)$ 为 $i^*:H^*(X_\Sigma)\to H^*(X)$ 的像，可能须先把 $J_X$ 适当地投影到 $H^*(\mathcal V)$ 上，把它称为 $J_X$ 的“环面”部分，再对该部分应用 $i_!$。一般而言，我们认为 $J_{\mathcal V}$ 与 $J_X$ 的确切关系值得进一步研究。

最后，还应评论此定理在 Calabi--Yau 完全交情形下的含义。如\hyperref[chap:mirror-constructions]{第四章}，设 $\Delta$ 为\IndexLocation{indexloc:1218:405}自反多胞形，$X_\Sigma$ 为对相应扇作\IndexLocation{indexloc:1096:405}极大射影细分得到的环面簇。进一步假设 $X_\Sigma$ 光滑。
\SourcePage{406}
为简便起见，只考虑 $V\subset X_\Sigma$ 为反典范除子的情形，并可假设 $V$ 光滑。于是定理~\ref{thm:givental-toric-mirror} 的 nef 假设当然成立，$J_{\mathcal V}$ 与 $I_{\mathcal V}$ 因而由坐标变换联系。如刚解释的，$J_{\mathcal V}$ 反映来自 $V$ 的“环面”部分的亏格零引力关联数。所以，形式函数 $I_{\mathcal V}$ 想必应告诉我们一些关于镜像 $V^\circ$ 的信息。

\hyperref[sec:hypergeometric]{第~\ref*{sec:hypergeometric} 节}研究过这一点，现在回顾当时的发现。首先，由\hyperref[subsec:nef-gkz]{第~\ref*{subsec:nef-gkz} 小节}，$I_{\mathcal V}$ 满足由式~\eqref{eq:nef-gkz-configuration} 中集合 $\mathcal A$ 确定的 GKZ 方程。因 $\Sigma(1)=\Delta^\circ\cap N-\{0\}$，可见 $\mathcal A=(\Delta^\circ\cap N)\times\{1\}$。令 $V^\circ$ 为 $V$ 的 Batyrev 镜像，它由自反多胞形 $\Delta^\circ$ 确定。再回顾\hyperref[subsec:toric-gkz]{第~\ref*{subsec:toric-gkz} 小节}，$\mathcal A=(\Delta^\circ\cap N)\times\{1\}$ 的 GKZ 方程是 $V^\circ$ 的 Picard--Fuchs 方程。看起来很好，但有两个困难。首先，GKZ 方程并不包含全部 Picard--Fuchs 方程，因此尚未证明 $I_{\mathcal V}$ 来自 $V^\circ$ 的周期，虽然 \SourceCite{Givental4} 未加证明地声称如此。其次，该小节并不处理 $V^\circ$ 的完整模空间，而只处理由环境环面簇所诱导的 $V^\circ$ 的形变给出的较小模空间。

因此，还未完全理解此定理与本书前几章所讨论镜像对称的关系。当 Calabi--Yau 超曲面 $V$ 为三维簇时，我们希望\hyperref[subsec:toric-mirror-conjecture]{第~\ref*{subsec:toric-mirror-conjecture} 小节}的猜想可由此定理推出。下面的例题~\ref{ex:octic-toric-mirror-theorem} 将说明这是怎样实现的。一般而言，镜像对称与此定理的关系，确实应进一步研究。

用两个例题结束本小节，展示此定理的效力。第一个是本书中已多次见过的例子，即 $\mathbb P(1,1,2,2,2)$ 的极小解奇异中的 Calabi--Yau 超曲面。
\SourceExampleNumber{11.2.5}{1}
\begin{example}\label{ex:octic-toric-mirror-theorem}
设 $\Sigma$ 是从 $\mathbb P(1,1,2,2,2)$ 的通常扇得到的扇：细分由 $(-1,-2,-2,-2)$ 与 $(1,0,0,0)$ 张成的锥，插入由 $(0,-1,-1,-1)$ 张成的射线。对于此扇，
\begin{equation}\label{eq:octic-toric-mirror-rays}
\begin{aligned}
\Sigma(1)=\{&(-1,-2,-2,-2),(1,0,0,0),(0,1,0,0),\\
&(0,0,1,0),(0,0,0,1),(0,-1,-1,-1)\},
\end{aligned}
\end{equation}
其中稍微滥用记号，把一维锥与其本原生成元等同。

此扇是 $\mathbb P(1,1,2,2,2)$ 标准扇的唯一极大射影细分。注意，式~\eqref{eq:octic-toric-mirror-rays} 中各点的凸包，正是 Batyrev 构造中使用的多胞形 $\Delta^\circ$。在此背景下，得到 Calabi--Yau 超曲面 $V\subset X_\Sigma$。例题~\ref{ex:weighted-extremal-transition} 证明，Kähler 锥由环面除子 $D_1,D_3$ 的类张成，这里按式~\eqref{eq:octic-toric-mirror-rays} 的次序为各环面除子编号。该例还指出线性等价关系 $D_1\sim D_2$、$D_3\sim D_4\sim D_5$、$D_6\sim-2D_1+D_3$，由此可用 $D_1,D_3$ 表示全部除子类。此外，$V\in|4D_3|$。按本小节的记号，令 $\mathcal V=\mathcal O_{X_\Sigma}(4D_3)$。

为与前几章的记号一致，将 Kähler 锥生成元排序为 $(D_3,D_1)$。令 $\{\beta_1,\beta_2\}$ 为 $M(X_\Sigma)$ 中与 $\{D_3,D_1\}$ 对偶的基，于是一般元素 $\beta\in M(X_\Sigma)$ 为 $\beta=d_1\beta_1+d_2\beta_2$。例题~\ref{ex:weighted-extremal-transition} 所计算的交数
\SourcePage{407}
给出 $\beta_1=D_3^2D_1$、$\beta_2=D_3^3-2D_3^2D_1$；这里通过 Poincaré 对偶，将 $\beta_i$ 看作 $H^6(X_\Sigma)$ 的元素。该例还指出，$|D_1|$ 定义映射 $X_\Sigma\to\mathbb P^1$，故 $D_1^2=0$。

为考察环面镜像定理在此例中的含义，先写出 $I_{\mathcal V}$。将式~\eqref{eq:toric-complete-intersection-i-function} 用于本例，得到
\begin{equation}\label{eq:octic-toric-i-function}
\begin{aligned}
I_{\mathcal V}={}&e^{(t_0+t_1D_3+t_2D_1)/\hbar}4D_3\sum_{d_1,d_2\ge0}z_1^{d_1}z_2^{d_2}\\
&\times\frac{\prod_{m=1}^{4d_1}(4D_3+m\hbar)\prod_{m=-\infty}^0(D_3-2D_1+m\hbar)}{\prod_{m=1}^{d_2}(D_1+m\hbar)^2\prod_{m=1}^{d_1}(D_3+m\hbar)^3\prod_{m=-\infty}^{d_1-2d_2}(D_3-2D_1+m\hbar)}.
\end{aligned}
\end{equation}
注意，对 $i=1,2$ 设 $z_i=e^{t_i}$。用 $z_i$ 而非 $q_i$，是为了与\hyperref[chap:hodge-yukawa]{第五章}、\hyperref[chap:moduli]{第六章}一致。这些坐标 $z_i$ 在例题~\ref{ex:weighted-complex-moduli} 中定义，并用于这两章中的多个例题。

环面镜像定理断言，适当变量变换后 $I_{\mathcal V}$ 等于 $J_{\mathcal V}$。此外，因
\[
J_{\mathcal V}=e^{(t_0+t_1D_3+t_2D_1)/\hbar}4D_3\bigl(1+o(\hbar^{-1})\bigr),
\]
可展开 $I_{\mathcal V}$，读出变量变换。由式~\eqref{eq:octic-toric-i-function}，
\begin{equation}\label{eq:octic-toric-i-period-expansion}
I_{\mathcal V}=e^{t_0/\hbar}4D_3\left(y_0+y_{1,1}\frac{D_3}{\hbar}+y_{1,2}\frac{D_1}{\hbar}
+y_{2,1}\frac{D_1D_3}{\hbar^2}+y_{2,2}\frac{D_3^2}{\hbar^2}+y_3\frac{D_1D_3^2}{\hbar^3}\right),
\end{equation}
其中
\[
y_0=\sum_{d_1\ge2d_2\ge0}z_1^{d_1}z_2^{d_2}\frac{(4d_1)!}{(d_2!)^2(d_1!)^3(d_1-2d_2)!},
\]
\[
\begin{aligned}
y_{1,1}={}&t_1y_0+\sum_{d_1\ge2d_2\ge0}z_1^{d_1}z_2^{d_2}\frac{(4d_1)!}{(d_2!)^2(d_1!)^3(d_1-2d_2)!}
\bigl(4\Psi(4d_1)-3\Psi(d_1)-\Psi(d_1-2d_2)\bigr)\\
&+\sum_{0\le d_1<2d_2}z_1^{d_1}z_2^{d_2}\frac{(-1)^{d_1+1}(2d_2-d_1-1)!(4d_1)!}{(d_2!)^2(d_1!)^3},
\end{aligned}
\]
\[
\begin{aligned}
y_{1,2}={}&t_2y_0+\sum_{d_1\ge2d_2\ge0}z_1^{d_1}z_2^{d_2}\frac{(4d_1)!}{(d_2!)^2(d_1!)^3(d_1-2d_2)!}
\bigl(2\Psi(d_1-2d_2)-2\Psi(d_2)\bigr)\\
&-2\sum_{0\le d_1<2d_2}z_1^{d_1}z_2^{d_2}\frac{(-1)^{d_1+1}(2d_2-d_1-1)!(4d_1)!}{(d_2!)^2(d_1!)^3}.
\end{aligned}
\]
这里使用了经典函数 $\Psi(n)=\sum_{i=1}^n1/i$，并约定 $\Psi(0)=0$。注意 $y_0,y_{1,1},y_{1,2}$ 是 $t_1,t_2$ 的函数。

由此容易看出，定理~\ref{thm:givental-toric-mirror} 中的 $f_0,f_1,f_2,h$ 为
\[
e^{f_0}=y_0,\qquad f_1=\frac{y_{1,1}}{y_0}-t_1,\qquad f_2=\frac{y_{1,2}}{y_0}-t_2,\qquad h=0,
\]
\SourcePage{408}
因此坐标变换为 $t_0\mapsto t_0+f_0\hbar$，以及 $t_i\mapsto t_i+f_i=y_{1,i}/y_0$，$i=1,2$。于是，对于
\[
(s_1,s_2)=\left(\frac{y_{1,1}(t_1,t_2)}{y_0(t_1,t_2)},\frac{y_{1,2}(t_1,t_2)}{y_0(t_1,t_2)}\right),
\]
环面镜像定理给出
\begin{equation}\label{eq:octic-toric-mirror-i-j-equality}
J_{\mathcal V}(s_1,s_2)=\frac{I_{\mathcal V}(t_1,t_2)}{y_0(t_1,t_2)}.
\end{equation}
令 $q_i=e^{s_i}$，上述变量变换成为
\[
(q_1,q_2)=\left(\exp\left(\frac{y_{1,1}}{y_0}\right),\exp\left(\frac{y_{1,2}}{y_0}\right)\right),
\]
这就是本例的镜像映射。例题~\ref{ex:weighted-hypergeometric-mirror-map} 用另一种方法计算了坐标 $q_2$，结果完全一致。

接着说明式~\eqref{eq:octic-toric-mirror-i-j-equality} 与用 A 模型、B 模型关联函数表述的“经典”镜像对称概念的关系。策略是证明 $I_{\mathcal V}/y_0$ 满足某些含 B 模型关联函数的微分方程。因此，$J_{\mathcal V}$ 满足相同的微分方程，进而得到所需的关联函数相等。

从 $I_{\mathcal V}$ 开始。由命题~\ref{prop:nef-i-gkz}，知道它满足来自 $V^\circ$ 的 GKZ 方程，但如前所述，这些并非 $V^\circ$ 的全部 Picard--Fuchs 方程。不过，例题~\ref{ex:weighted-gkz} 已计算了缺少的 Picard--Fuchs 方程，见式~\eqref{eq:weighted-gkz-second} 和~\eqref{eq:weighted-gkz-third}。直接计算表明 $I_{\mathcal V}$ 满足这些方程。（对此应当有一个一般定理，但我们尚不知道证明。）

例题~\ref{ex:weighted-gkz} 所构造的 Picard--Fuchs 方程采用了一个特定的全纯 $3$ 形式 $\Omega$。还知道，$y_0$ 是除常数倍外唯一在 $z_1=z_2=0$ 处全纯的周期。于是 $\widetilde\Omega=\Omega/y_0$ 是除常数倍外的归一化全纯 $3$ 形式，所以 $I_{\mathcal V}/y_0$ 满足 $\widetilde\Omega$ 的 Picard--Fuchs 方程。从变量 $z_i=e^{t_i}$ 变到 $q_i=e^{s_i}$ 后，仍然如此。

在坐标 $q_i$ 下，$V^\circ$ 的归一化 Yukawa 耦合导出一些漂亮的 Picard--Fuchs 方程。此时，归一化耦合为
\[
Y_{ijk}=-\frac1{(2\pi i)^3}\int_{V^\circ}\widetilde\Omega\wedge\nabla_{\delta_i}\nabla_{\delta_j}\nabla_{\delta_k}\widetilde\Omega,
\]
其中 $\delta_j=\partial/\partial s_j$，$\nabla$ 是镜像族的 Gauss--Manin 联络。这与\hyperref[subsec:toric-mirror-conjecture]{第~\ref*{subsec:toric-mirror-conjecture} 小节}的定义在适当的常数归一化下相对应；那里写 $q_j=e^{2\pi iu_j}$，而这里使用 $q_j=e^{s_j}$。\footnote{这在先验上可能不同于定义~\ref{def:normalized-yukawa} 中的归一化耦合。那里假设整性猜想，并使用\hyperref[subsec:mirror-map-definition]{第~\ref*{subsec:mirror-map-definition} 小节}定义的镜像映射。这里不假设整性猜想，而使用\hyperref[subsec:toric-mirror-map]{第~\ref*{subsec:toric-mirror-map} 小节}给出的环面版本镜像映射。利用 \SourceCite{CdFKM} 中对单值性的详细描述，可验证本例中二者相同。但我们更希望说明，即使不了解单值性的全部细节，仍能得到良好结果。}采用 $H^2(V)$ 的上同调基 $T_1=i^*(D_3)$、$T_2=i^*(D_1)$，目标是证明
\begin{equation}\label{eq:octic-toric-mirror-correlation-equality}
\langle T_i,T_j,T_k\rangle=Y_{ijk}\qquad\text{对所有 }i,j,k.
\end{equation}
\SourcePage{409}
回顾例题~\ref{ex:octic-toric-mirror-correlations}，可显式计算 $Y_{ijk}$。于是该式给出 $V$ 的全部 A 模型关联函数 $\langle T_i,T_j,T_k\rangle$ 的显式公式。

为证明此式，先注意可用 $Y_{ijk}$ 描述 Gauss--Manin 联络。更确切地，断言 $H^3(V^\circ,\mathbb C)$ 有一组基 $e_0,e_1,e_2,e^1,e^2,e^0$，使 $e_0=\widetilde\Omega$，且
\[
\nabla_{\delta_i}(e_0)=e_i,\qquad\nabla_{\delta_i}(e_j)=\sum_kY_{ijk}e^k,\qquad
\nabla_{\delta_i}(e^j)=\delta_{ij}e^0,\qquad\nabla_{\delta_i}(e^0)=0,
\]
与式~\eqref{eq:multimoduli-connection} 相同。对于\hyperref[subsec:normalized-yukawa-mirror]{第~\ref*{subsec:normalized-yukawa-mirror} 小节}定义的归一化 Yukawa 耦合，命题~\ref{prop:normalized-yukawa-pf} 及其证明后的讨论已给出这一结果。虽然证明要求整性猜想，但只要在 $\mathbb C$ 上而非 $\mathbb Z$ 上工作，所有内容都适用于此情形。

利用上述基，直接计算可得
\[
\nabla_{\delta_1}\left(\frac{Y_{112}\nabla_{\delta_1}^2-Y_{111}\nabla_{\delta_1}\nabla_{\delta_2}}{Y_{112}^2-Y_{111}Y_{122}}\right)\widetilde\Omega=0,
\]
因为 $\widetilde\Omega=e_0$。因此，这是一个 Picard--Fuchs 方程，而 $I_{\mathcal V}/y_0$ 满足 Picard--Fuchs 方程，故
\[
\frac\partial{\partial s_1}\left(\frac{Y_{112}\frac{\partial^2}{\partial s_1^2}-Y_{111}\frac{\partial^2}{\partial s_1\partial s_2}}{Y_{112}^2-Y_{111}Y_{122}}\right)\frac{I_{\mathcal V}}{y_0}=0.
\]
由环面镜像定理~\eqref{eq:octic-toric-mirror-i-j-equality}，得到
\begin{equation}\label{eq:octic-toric-j-yukawa-differential-equation}
\frac\partial{\partial s_1}\left(\frac{Y_{112}\frac{\partial^2}{\partial s_1^2}-Y_{111}\frac{\partial^2}{\partial s_1\partial s_2}}{Y_{112}^2-Y_{111}Y_{122}}\right)J_{\mathcal V}(s_1,s_2)=0,
\end{equation}
这里忽略了 $t_0$ 的依赖。还应指出，显式计算表明 $Y_{112}^2-Y_{111}Y_{122}$ 不恒为零。

下一步，将它与 $V$ 的 Givental 函数 $J_V$ 联系。包含映射 $i:V\subset X_\Sigma$ 诱导同构 $i^*:H^2(X_\Sigma)\to H^2(V)$，例题~\ref{ex:weighted-extremal-transition} 已说明，$V$ 的 Kähler 锥由 $T_1=i^*(D_3)$、$T_2=i^*(D_1)$ 张成。因此，$J_{\mathcal V}$ 与 $J_V$ 都可看作 $s_1,s_2$ 的函数，仍忽略 $t_0$。不过，尚不知道 $i_!(J_V)$ 是否等于 $J_{\mathcal V}$。环境空间为 $\mathbb P^n$ 时，命题~\ref{prop:complete-intersection-j-pushforward} 证明了二者相等。但证明用了 $\mathbb P^n$ 的凸性，此时的 $X_\Sigma$ 并不凸。考察该命题的证明，可验证：只要
\begin{equation}\label{eq:octic-virtual-class-pushforward-assumption}
j_*\bigl([\overline M_{0,2}(V,\beta)]^{\mathrm{virt}}\bigr)
=\operatorname{Euler}(\mathcal V_{\beta,2})\cap[\overline M_{0,2}(X_\Sigma,\beta)]^{\mathrm{virt}},
\end{equation}
仍有 $i_!(J_V)=J_{\mathcal V}$。我们推测此式为真，但尚无完整证明。因此，本例余下部分将\emph{假设}该式成立。特别地，这允许假设 $i_!(J_V)=J_{\mathcal V}$。

还可验证，本例中 $i_!$ 在偶次上同调上为单射。结合式~\eqref{eq:octic-toric-j-yukawa-differential-equation} 与 $i_!(J_V)=J_{\mathcal V}$，得到
\begin{equation}\label{eq:octic-ordinary-j-yukawa-quantum-equation}
\hbar\frac\partial{\partial s_1}\left(\frac{Y_{112}\hbar^2\frac{\partial^2}{\partial s_1^2}-Y_{111}\hbar^2\frac{\partial^2}{\partial s_1\partial s_2}}{Y_{112}^2-Y_{111}Y_{122}}\right)J_V(s_1,s_2)=0.
\end{equation}
加入各个 $\hbar$ 因子，是为了将此式写成\hyperref[sec:quantum-cohomology-relations]{第~\ref*{sec:quantum-cohomology-relations} 节}意义下的\IndexLocation{indexloc:1150:409}量子微分方程。为应用定理~\ref{thm:quantum-differential-relations}，展开此式，
\SourcePage{410}
得到
\[
\left(\frac{Y_{112}\hbar^3\frac{\partial^3}{\partial s_1^3}-Y_{111}\hbar^3\frac{\partial^3}{\partial s_1^2\partial s_2}}{Y_{112}^2-Y_{111}Y_{122}}
+\hbar P(\hbar\partial/\partial s,q,\hbar)\right)J_V(s_1,s_2)=0,
\]
其中 $P(\hbar\partial/\partial s,q,\hbar)$ 是某个微分算子。定理~\ref{thm:quantum-differential-relations} 表明，代入 $\hbar\partial/\partial s_i\mapsto T_i$、$\hbar\mapsto0$，得到 $V$ 的小量子上同调中的多项式关系。因此，$H^*(V)$ 中有
\[
\frac{Y_{112}T_1*T_1*T_1-Y_{111}T_1*T_1*T_2}{Y_{112}^2-Y_{111}Y_{122}}=0.
\]
由式~\eqref{eq:three-point-quantum-integral}，
\[
\int_V T_i*T_j*T_k=\langle T_i,T_j,T_k\rangle,
\]
所以上式给出
\[
Y_{112}\langle T_1,T_1,T_1\rangle-Y_{111}\langle T_1,T_1,T_2\rangle=0.
\]
刚才的论证适用于一般下标，故对任意 $1\le i,j,k,a,b,c\le2$，均有
\[
Y_{ijk}\langle T_a,T_b,T_c\rangle-Y_{abc}\langle T_i,T_j,T_k\rangle=0.
\]
由此可知，有一个共同因子 $c$，使
\[
\langle T_i,T_j,T_k\rangle=cY_{ijk}\qquad\text{对所有 }i,j,k.
\]
两侧都来自势函数，故各四阶混合导数的相容性给出 $(\partial_1c)Y_{12k}=(\partial_2c)Y_{11k}$，$k=1,2$。因为 $Y_{112}^2-Y_{111}Y_{122}$ 不恒为零，得到 $\partial_1c=\partial_2c=0$，所以 $c$ 为常数。

把 $\langle T_i,T_j,T_k\rangle$ 看作 $q_1,q_2$ 的形式函数，有 $\langle T_1,T_1,T_1\rangle(0,0)=\int_VT_1^3=8$。因此，为证明 $c=1$，须证明 $Y_{111}(0,0)=8$。令 $N_1$ 为 $V^\circ$ 的复模空间中边界除子 $z_1=0$ 周围单值性变换的对数，沿用例题~\ref{ex:weighted-yukawa} 的记号。\SourceCite{CdFKM} 的显式计算表明，有 $\mathcal H$ 的整平坦截面 $\beta^2,\alpha_2$，它们是一组 $\mathbb Z$ 基的部分元素，且满足 $\langle\alpha_2,\beta^2\rangle=1$、$N_1^3(\beta^2)=-8\alpha_2$。注意，\SourceCite{CdFKM} 中单值性的取向与这里相反。类似于例题~\ref{ex:quintic-normalized-yukawa} 的论证给出 $Y_{111}(0,0)=8$。故 $c=1$，式~\eqref{eq:octic-toric-mirror-correlation-equality} 得证。

例题~\ref{ex:octic-toric-mirror-correlations} 把 $\langle T_1,T_1,T_1\rangle$ 记为 $\langle H,H,H\rangle$，现在已证明该例中这个关联函数的公式。还可把式~\eqref{eq:octic-toric-mirror-correlation-equality} 与\hyperref[subsec:toric-mirror-conjecture]{第~\ref*{subsec:toric-mirror-conjecture} 小节}中镜像猜想的 Hodge 理论版本联系。如式~\eqref{eq:toric-mirror-correlations} 附近的讨论，A 模型与 B 模型关联函数在相差常数的意义下相等，意味着本例中猜想~\ref{conj:toric-hodge-mirror-surjective} 成立。更确切地，对于这里考虑的三维簇 $V,V^\circ$，已证明该猜想由定理~\ref{thm:givental-toric-mirror} 和式~\eqref{eq:octic-virtual-class-pushforward-assumption} 推出。此外，利用上一段，还可说明 $(V^\circ,V)$ 是定义~\ref{def:mathematical-mirror-pair} 意义下的数学镜像对。

不过，在充分发挥环面镜像定理的作用之前，还有更多内容可说。将看到，式~\eqref{eq:octic-toric-mirror-i-j-equality} 给出 $V$ 的 Gromov--Witten 势的一个漂亮公式。采用 $V$ 的偶次上同调的以下基：
\[
\begin{gathered}
T_0=1,\qquad T_1=i^*(D_3),\qquad T_2=i^*(D_1),\\
T^1=\frac{i^*(D_1D_3)}4,\qquad T^2=\frac{i^*(D_3^2-2D_1D_3)}4,\qquad T^0=\frac{i^*(D_1D_3^2)}4.
\end{gathered}
\]
\SourcePage{411}
可验证，杯积下 $T^1,T^2$ 与 $T_1,T_2$ 对偶，$T^0$ 与 $T_0$ 亦然。此时，变量为 $t_0,s_1,s_2$；由命题~\ref{prop:cy-threefold-j-potential}，$J_V$ 由公式
\begin{equation}\label{eq:octic-ordinary-j-potential}
\begin{aligned}
J_V=e^{t_0/\hbar}\Biggl(&1+\hbar^{-1}(s_1T_1+s_2T_2)
+\hbar^{-2}\left(\frac{\partial\Phi}{\partial s_1}T^1+\frac{\partial\Phi}{\partial s_2}T^2\right)\\
&+\hbar^{-3}\left(s_1\frac{\partial\Phi}{\partial s_1}+s_2\frac{\partial\Phi}{\partial s_2}-2\Phi\right)T^0\Biggr)
\end{aligned}
\end{equation}
给出，其中 $\Phi$ 是 $V$ 的 Gromov--Witten 势。

仍假设式~\eqref{eq:octic-virtual-class-pushforward-assumption}，故 $J_{\mathcal V}=i_!(J_V)$。因为 $i_!(i^*\alpha)=\alpha\cup4D_3$，对式~\eqref{eq:octic-ordinary-j-potential} 应用 $i_!$，得到
\[
\begin{aligned}
J_{\mathcal V}(s_1,s_2)=e^{t_0/\hbar}4D_3\Biggl(&1+s_1\frac{D_3}{\hbar}+s_2\frac{D_1}{\hbar}
+\frac{\partial\Phi}{\partial s_1}\frac{D_1D_3}{4\hbar^2}
+\frac{\partial\Phi}{\partial s_2}\frac{D_3^2-2D_1D_3}{4\hbar^2}\\
&+\left(s_1\frac{\partial\Phi}{\partial s_1}+s_2\frac{\partial\Phi}{\partial s_2}-2\Phi\right)\frac{D_1D_3^2}{4\hbar^3}\Biggr).
\end{aligned}
\]
为与 $I_{\mathcal V}/y_0$ 比较，注意式~\eqref{eq:octic-toric-i-period-expansion} 可写成
\[
\begin{aligned}
\frac{I_{\mathcal V}}{y_0}(t_1,t_2)=e^{t_0/\hbar}4D_3\Biggl(&1+\frac{y_{1,1}}{y_0}\frac{D_3}{\hbar}+\frac{y_{1,2}}{y_0}\frac{D_1}{\hbar}
+\frac{y_{2,1}+2y_{2,2}}{y_0}\frac{D_1D_3}{\hbar^2}\\
&+\frac{y_{2,2}}{y_0}\frac{D_3^2-2D_1D_3}{\hbar^2}+\frac{y_3}{y_0}\frac{D_1D_3^2}{\hbar^3}\Biggr).
\end{aligned}
\]
将镜像映射 $(s_1,s_2)=(y_{1,1}/y_0,y_{1,2}/y_0)$ 记为 $(s_1,s_2)=\Psi(t_1,t_2)$。如式~\eqref{eq:octic-toric-mirror-i-j-equality}，环面\IndexLocation{indexloc:0855:411}\IndexLocation{indexloc:1013:411}镜像定理表明，上述两个表达式在变量变换 $(s_1,s_2)=\Psi(t_1,t_2)$ 后相等。比较系数，得到 $V$ 的 Gromov--Witten 势的显式公式：
\begin{equation}\label{eq:octic-toric-mirror-potential}
\Phi(\Psi(t_1,t_2))=2\left(\frac{y_{1,1}}{y_0}\frac{y_{2,1}+2y_{2,2}}{y_0}
+\frac{y_{1,2}}{y_0}\frac{y_{2,2}}{y_0}-\frac{y_3}{y_0}\right),
\end{equation}
其中 $\Psi(t_1,t_2)=(y_{1,1}/y_0,y_{1,2}/y_0)$ 是镜像映射。这与\hyperref[sec:quintic-mirror-theorem]{第~\ref*{sec:quintic-mirror-theorem} 节}导出的\IndexLocation{indexloc:1190:411}五次三维簇的 Gromov--Witten 势公式十分相似。

最后一个观察是，比较上述 $J_{\mathcal V}$ 和 $I_{\mathcal V}/y_0$ 的公式，也给出 A \IndexLocation{indexloc:1003:411}模型关联函数的一些漂亮公式。例如，已知
\[
\langle T_1,T_1,T_1\rangle=\frac{\partial^3\Phi}{\partial s_1^3}.
\]
于是，把 $J_{\mathcal V}(s_1,s_2)=I_{\mathcal V}/y_0(t_1,t_2)$ 对 $s_1$ 求两次导数，再比较 $D_1D_3$ 的系数，得到
\begin{equation}\label{eq:octic-yukawa-period-second-derivative}
\langle T_1,T_1,T_1\rangle=4\frac{\partial^2}{\partial s_1^2}\left(\frac{y_{2,1}+2y_{2,2}}{y_0}\right).
\end{equation}
这个公式比例题~\ref{ex:octic-toric-mirror-correlations} 中考虑的公式\emph{容易推导得多}。也可将其看作归一化 Yukawa 耦合 $Y_{111}$ 的一个简洁公式。

此例很有启发性，因为其中许多计算应能推广到其他环面 Calabi--Yau 三维簇。特别地：
\SourcePage{412}
\begin{itemize}
\item 式~\eqref{eq:octic-toric-mirror-potential} 应能推广到任意作为环面簇 $X_\Sigma$ 中 nef 完全交的 Calabi--Yau 三维簇 $V$。由于 $H^2(X_\Sigma)\to H^2(V)$ 未必为满射，须定义 $V$ 的 Gromov--Witten 势的“环面”部分，再用类似于该式的公式，显式地以 $I_{\mathcal V}/y_0$ 表示它。
\item Gromov--Witten 势是 A 模型关联函数 $\langle T_i,T_j,T_k\rangle$ 的势函数。因此，该式表明，$I_{\mathcal V}/y_0$ 的系数给出 $V^\circ$ 的 B 模型关联函数 $Y_{ijk}$ 的一个显式势函数，含义如引理~\ref{lem:yukawa-potential}。应该可以直接证明这一点，而不使用环面镜像定理。
\item 式~\eqref{eq:octic-yukawa-period-second-derivative} 实质上给出了 $Y_{111}$ 的简单公式。一般而言，$V^\circ$ 的归一化 Yukawa 耦合应有类似的公式，以 $I_{\mathcal V}/y_0$ 表示。如上，也应该能够直接证明这一点。
\end{itemize}
我们认为，环面簇中的三维 nef 完全交还有许多有趣的未解问题。
\end{example}
作为定理~\ref{thm:givental-toric-mirror} 的第二个应用，研究 $\mathcal V=0$ 时的情形，此时 $X$ 就是环面簇 $X_\Sigma$。将看到，镜像定理与前面的结果很好地结合，描述 $X_\Sigma$ 的量子上同调。
\SourceExampleNumber{11.2.5}{2}
\begin{example}\label{ex:toric-fano-and-hirzebruch-quantum-ring}
本例将说明，若 $X=X_\Sigma$ 为光滑 Fano 环面簇，则例题~\ref{ex:batyrev-quantum-ring} 的 Batyrev 量子环 $H_\omega^*(X)$ 与小量子上同调环相同。将结合定理~\ref{thm:givental-toric-mirror} 与命题~\ref{prop:nef-i-gkz} 证明这一点。还将看到，当 $-K_X$ 为 nef 而非丰沛时，$H_\omega^*(X)$ 可能不同于量子上同调环，但该定理仍意味着两环之间存在关系。以 $X=\mathbb F_2$ 为例说明。

既然考虑环面簇 $X=X_\Sigma$ 本身，就在式~\eqref{eq:toric-complete-intersection-i-function} 中令 $\mathcal V$ 为零秩平凡丛，并分别将 $I_{\mathcal V}$、$J_{\mathcal V}$ 记为 $I_X$、$J_X$。注意，$X=\mathbb P^n$ 时，这与命题~\ref{prop:projective-space-i-equals-j} 的记号一致。有
\begin{equation}\label{eq:toric-zero-bundle-i-function}
I_X=e^{(t_0+\delta)/\hbar}\sum_{\beta\in M(X)}q^\beta
\frac{\prod_\rho\prod_{m=-\infty}^0(D_\rho+m\hbar)}{\prod_\rho\prod_{m=-\infty}^{D_\rho(\beta)}(D_\rho+m\hbar)},
\end{equation}
其中 $\delta=\sum_{i=1}^rt_iT_i$、$D_\rho(\beta)=\int_\beta D_\rho$。\hyperref[subsec:givental-nef-projective-mirror]{第~\ref*{subsec:givental-nef-projective-mirror} 小节}的式~\eqref{eq:complete-intersection-ordinary-j-function} 还给出了 $J_X$ 的公式。于是得到命题~\ref{prop:projective-space-i-equals-j} 的以下推广。
\begin{proposition}\label{prop:toric-fano-i-equals-j}
若 $X=X_\Sigma$ 为光滑 Fano 环面簇，则 $I_X=J_X$。
\end{proposition}
\begin{proof}
由式~\eqref{eq:j-function-asymptotic-normalization}，$J_X=e^{(t_0+\delta)/\hbar}(1+o(\hbar^{-1}))$。若知道 $X$ 为 Fano 时 $I_X$ 也如此，定理~\ref{thm:givental-toric-mirror} 的变量变换便是恒等变换，命题就由环面镜像定理立即推出。因此，只须证明
\begin{equation}\label{eq:toric-fano-i-normalization}
I_X=e^{(t_0+\delta)/\hbar}(1+o(\hbar^{-1})).
\end{equation}
为此，固定 Mori 锥 $M(X)$ 中的非零整类 $\beta$ 与 $\rho\in\Sigma(1)$，考察式~\eqref{eq:toric-zero-bundle-i-function} 中的分式
\[
\frac{\prod_{m=-\infty}^0(D_\rho+m\hbar)}{\prod_{m=-\infty}^{D_\rho(\beta)}(D_\rho+m\hbar)}.
\]
\SourcePage{413}
将此分式按 $\hbar^{-1}$ 的幂展开，研究首项。若 $D_\rho(\beta)>0$，首项容易看出为
\[
\frac1{D_\rho(\beta)!}\hbar^{-D_\rho(\beta)}.
\]
若 $D_\rho(\beta)=0$，分式就是 $1$；若 $D_\rho(\beta)=-1$，分式就是 $D_\rho$。最后，若 $D_\rho(\beta)<-1$，首项为
\[
(-1)^{-D_\rho(\beta)-1}D_\rho\,(-D_\rho(\beta)-1)!\,\hbar^{-D_\rho(\beta)-1}.
\]
所有情形下，首项中 $\hbar^{-1}$ 的幂次都至少为 $D_\rho(\beta)$，而且当且仅当 $D_\rho(\beta)<0$ 时严格大于它。

固定 $\beta$，对所有 $\rho$ 相乘，可见式~\eqref{eq:toric-zero-bundle-i-function} 展开中 $q^\beta$ 的系数的首项，其 $\hbar^{-1}$ 的幂次至少为 $\sum_\rho D_\rho(\beta)$。\hyperref[chap:toric]{第三章}已说明 $\sum_\rho D_\rho=-K_X$，因此该幂次至少为 $-\int_\beta K_X$。

$X$ 为 Fano，故 $-\int_\beta K_X$ 严格为正，所以首项中 $\hbar^{-1}$ 的幂次至少为一。但为证明式~\eqref{eq:toric-fano-i-normalization}，须证明至少为二。上述讨论表明，首项幂次为一的唯一可能是 $-\int_\beta K_X=1$，且所有 $\rho$ 均有 $D_\rho(\beta)\ge0$。因 $-K_X=\sum_\rho D_\rho$，可见存在 $\rho_0$，使 $D_{\rho_0}(\beta)=1$，而所有 $\rho\ne\rho_0$ 的 $D_\rho(\beta)=0$。须证明这不可能发生。

使用\hyperref[subsec:nef-gkz]{第~\ref*{subsec:nef-gkz} 小节}的引理~\ref{lem:nef-relation-lattice}，取 $k=0$。对于 $\beta\in H_2(X,\mathbb Z)$，该引理给出
\begin{equation}\label{eq:toric-curve-divisor-ray-relation}
\sum_\rho D_\rho(\beta)v_\rho=0,
\end{equation}
其中 $v_\rho$ 是 $\rho\cap N$ 的本原生成元，因为 $D_\rho(\beta)=D_\rho\cdot\beta$。上一段的情形会导致 $v_{\rho_0}=0$，矛盾。于是式~\eqref{eq:toric-fano-i-normalization} 得证，命题立即由定理~\ref{thm:givental-toric-mirror} 推出。
\end{proof}
接下来用 $I_X=J_X$ 研究光滑 Fano 环面簇 $X=X_\Sigma$ 的\IndexLocation{indexloc:1141:413}\IndexLocation{indexloc:1146:413}小量子上同调环。基本思想是：由\hyperref[subsec:nef-gkz]{第~\ref*{subsec:nef-gkz} 小节}，$I_X$ 满足一个 GKZ 方程组，故 $J_X$ 满足相同的方程；随后用\hyperref[sec:quantum-cohomology-relations]{第~\ref*{sec:quantum-cohomology-relations} 节}的\IndexLocation{indexloc:1150:413}量子微分方程理论，得到所需的量子上同调关系。

为实施此策略，须看看命题~\ref{prop:nef-i-gkz} 在这里的形式。这里 $\mathcal V=\bigoplus_{i=1}^k\mathcal L_i$ 为零丛，故 $k=0$。于是式~\eqref{eq:nef-gkz-configuration} 化为
\[
\mathcal A=\{v_\rho:\rho\in\Sigma(1)\},
\]
而引理~\ref{lem:nef-relation-lattice} 中的格 $\Lambda$，恰为 $v_\rho$ 之间整关系构成的格。此外，该引理通过将 $\beta$ 映到式~\eqref{eq:toric-curve-divisor-ray-relation} 的映射，把 $H_2(X,\mathbb Z)$ 与 $\Lambda$ 等同。

$\Lambda$ 同时也是相应 $\mathcal A$ 方程组的关系格，所以对于每个 $\beta\in\Lambda$，式~\eqref{eq:nef-gkz-box} 给出算子
\[
\Box_\beta=\prod_{D_\rho(\beta)>0}\partial_\rho^{D_\rho(\beta)}
-\prod_{D_\rho(\beta)<0}\partial_\rho^{-D_\rho(\beta)},
\]
\SourcePage{414}
因为 $\beta$ 对应于关系~\eqref{eq:toric-curve-divisor-ray-relation}。这里 $\partial_\rho=\partial/\partial\lambda_\rho$，由式~\eqref{eq:nef-gkz-q-coordinates}，用 $\lambda_\rho$ 定义 $q^\beta$ 如下：
\begin{equation}\label{eq:toric-quantum-monomial-lambda}
q^\beta=\prod_\rho\lambda_\rho^{D_\rho(\beta)}.
\end{equation}
若令
\[
\widetilde I=e^\delta\sum_{\beta\in M(X)}q^\beta
\frac{\prod_\rho\prod_{m=-\infty}^0(D_\rho+m)}{\prod_\rho\prod_{m=-\infty}^{D_\rho(\beta)}(D_\rho+m)},
\]
命题~\ref{prop:nef-i-gkz} 表明，所有 $\beta$ 均有 $\Box_\beta\widetilde I=0$。

注意，$\widetilde I$ 是在 $I_X$ 中置 $t_0=0$、$\hbar=1$ 所得。只须乘以 $e^{t_0/\hbar}$ 即可恢复 $t_0$，但恢复 $\hbar$ 须更小心：使用把 $\Box_\beta$ 中每个 $\partial_\rho$ 换成 $\hbar\partial_\rho$ 所得的算子 $\Box_{\beta,\hbar}$，即
\[
\Box_{\beta,\hbar}=\prod_{D_\rho(\beta)>0}(\hbar\partial_\rho)^{D_\rho(\beta)}
-\prod_{D_\rho(\beta)<0}(\hbar\partial_\rho)^{-D_\rho(\beta)}.
\]
容易证明 $\Box_\beta\widetilde I=0$ 蕴含 $\Box_{\beta,\hbar}I_X=0$，证明省略。

再定义
\[
\Box'_{\beta,\hbar}=\left(\prod_{D_\rho(\beta)>0}\lambda_\rho^{D_\rho(\beta)}\right)\Box_{\beta,\hbar},
\]
则 $\Box'_{\beta,\hbar}I_X=0$。令 $\delta_\rho=\lambda_\rho\partial_\rho$，将式~\eqref{eq:weighted-logarithmic-identities} 用于本情形，得到
\begin{equation}\label{eq:toric-quantum-gkz-operator}
\begin{aligned}
\Box'_{\beta,\hbar}={}&\prod_{D_\rho(\beta)>0}
\hbar\delta_\rho(\hbar\delta_\rho-\hbar)\cdots(\hbar\delta_\rho-(D_\rho(\beta)-1)\hbar)\\
&-q^\beta\prod_{D_\rho(\beta)<0}
\hbar\delta_\rho(\hbar\delta_\rho-\hbar)\cdots(\hbar\delta_\rho-(-D_\rho(\beta)-1)\hbar).
\end{aligned}
\end{equation}
$X$ 光滑且为 Fano，故命题~\ref{prop:toric-fano-i-equals-j} 给出 $I_X=J_X$，从而 $\Box'_{\beta,\hbar}J_X=0$。不过，在应用\hyperref[sec:quantum-cohomology-relations]{第~\ref*{sec:quantum-cohomology-relations} 节}的理论之前，须把 $\Box'_{\beta,\hbar}$ 写成关于 $\hbar\partial/\partial t_j$ 的微分算子。回顾 $t_j$ 是与 $H^2(X,\mathbb Z)$ 的基 $T_j$ 对应的变量。它们与 $\lambda_\rho$ 的联系如下：令 $\beta_j\in H_2(X,\mathbb Z)$ 为与 $T_j$ 对偶的基，再令
\[
q_j=\prod_\rho\lambda_\rho^{D_\rho(\beta_j)}.
\]
容易证明，$\prod_jq_j^{\int_\beta T_j}$ 等于式~\eqref{eq:toric-quantum-monomial-lambda} 所定义的 $q^\beta$。令 $q_j=e^{t_j}$，可见 $q^\beta$ 相对于 $t_j$ 具有通常的含义。

一般不能把式~\eqref{eq:toric-quantum-gkz-operator} 中的 $\hbar\delta_\rho$ 表示为 $\hbar\partial/\partial t_j$。关键观察是，对偶环面 $T^\circ=M\otimes\mathbb C^*$ 通过 $(m\otimes t)\cdot\lambda_\rho=t^{\langle m,v_\rho\rangle}\lambda_\rho$ 作用于这些变量。$q_j$ 在此作用下不变。如式~\eqref{eq:nef-logarithmic-identities} 所述，
\SourcePage{415}
对于 $T^\circ$ 不变的函数，有
\begin{equation}\label{eq:toric-quantum-gkz-coordinate-derivative}
\hbar\delta_\rho=\sum_j D_\rho(\beta_j)\hbar\frac\partial{\partial t_j}.
\end{equation}
此恒等式使式~\eqref{eq:toric-quantum-gkz-operator} 可用 $\hbar\partial/\partial t_j$ 写出，即 $\Box'_{\beta,\hbar}$ 是\hyperref[sec:quantum-cohomology-relations]{第~\ref*{sec:quantum-cohomology-relations} 节}意义下的量子微分算子。于是定理~\ref{thm:quantum-differential-relations} 表明，作代换
\[
\hbar\frac\partial{\partial t_j}\longmapsto T_j,\qquad\hbar\longmapsto0,
\]
就由该算子得到量子上同调关系。应用于式~\eqref{eq:toric-quantum-gkz-operator}，得到 $X$ 的小量子上同调中的关系
\begin{equation}\label{eq:toric-fano-quantum-stanley-reisner-relation}
\prod_{D_\rho(\beta)>0}D_\rho^{D_\rho(\beta)}
=q^\beta\prod_{D_\rho(\beta)<0}D_\rho^{-D_\rho(\beta)},
\end{equation}

可用它证明 Batyrev 量子环 $H_\omega^*(X)$ 同构于 $X$ 的小量子上同调环。令 $H^*(X)$ 采用小量子乘积 $*$，$\mathbb C[x_\rho]$ 为每个 $\rho\in\Sigma(1)$ 对应一个变量 $x_\rho$ 的多项式环。于是 $x_\rho\mapsto D_\rho$ 定义环同态 $\mathbb C[x_\rho]\to H^*(X)$。由上式，理想
\[
\left\langle\prod_{D_\rho(\beta)>0}x_\rho^{D_\rho(\beta)}
-q^\beta\prod_{D_\rho(\beta)<0}x_\rho^{-D_\rho(\beta)}:\beta\in H_2(X,\mathbb Z)\right\rangle
\]
包含于此映射的核。随着 $\beta$ 遍历 $H_2(X,\mathbb Z)$，关系~\eqref{eq:toric-curve-divisor-ray-relation}，即
\[
\sum_{D_\rho(\beta)>0}D_\rho(\beta)v_\rho
=\sum_{D_\rho(\beta)<0}-D_\rho(\beta)v_\rho,
\]
张成关系格 $\Lambda$。因此，上述理想就是式~\eqref{eq:quantum-sr-all-relations} 的理想，也就是 \SourceCite[定义 5.1]{Batyrev2} 给出的量子 Stanley--Reisner 理想 $SR_\omega(\Sigma)$ 的形式。于是 $SR_\omega(\Sigma)$ 包含于映射 $\mathbb C[x_\rho]\to H^*(X)$ 的核。此核还包含理想
\[
P(\Sigma)=\left\langle\sum_{i=1}^s\langle m,v_i\rangle x_i:m\in M\right\rangle,
\]
故有良定义的环同态
\[
\mathbb C[x_\rho]/(P(\Sigma)+SR_\omega(\Sigma))\longrightarrow H^*(X).
\]
例题~\ref{ex:batyrev-quantum-ring} 将左侧商环定义为 Batyrev 量子环 $H_\omega^*(X)$，所以得到环同态
\[
\varphi:H_\omega^*(X)\longrightarrow H^*(X),
\]
提醒读者，右侧采用小量子乘积。

现在断言 $\varphi$ 为同构。为此，回顾总可假设 $H^2(X)$ 的基 $T_j$ 位于 $X$ 的 Kähler 锥内部。这允许取极限 $q_j\to0$，此时 $\varphi$ 变为描述环面簇 $X$ 的通常上同调环的同构。因处理的是有限维向量空间，$\varphi$ 在 $q_j=0$ 的某个邻域中仍为同构。此外，若将 $q^\beta$ 看作形式变量，也得到同构。因此，$X$ 光滑且为 Fano 时，Batyrev 量子环与 $X$ 的通常小量子环相同，如所断言。
\SourcePage{416}
若削弱 $X$ 为 Fano 的假设，只假设 $-K_X$ 为 nef，则定理~\ref{thm:givental-toric-mirror} 仍适用，但结果更复杂，因为可能没有 $I_X=J_X$，而涉及非平凡变量变换。关键是，$-K_X$ 仅为 nef 时，可能有
\[
I_X\ne e^{(t_0+\delta)/\hbar}(1+o(\hbar^{-1})).
\]
于是该定理断言 $I_X$ 与 $J_X$ 由非平凡变量变换联系。这种情形下，仍可得到 $X$ 的小量子环中的有趣关系。命题~\ref{prop:nef-i-gkz} 仍适用，故 $\Box'_{\beta,\hbar}I_X=0$，其中算子由式~\eqref{eq:toric-quantum-gkz-operator} 给出。但不像 $\Box'_{\beta,\hbar}J_X=0$ 那样，直接用定理~\ref{thm:quantum-differential-relations} 得到式~\eqref{eq:toric-fano-quantum-stanley-reisner-relation}，这里须先改变变量，在 $H^*(X)$ 中得到该关系的一个“扭曲”版本。这正是 $H_\omega^*(X)$ 与 $X$ 的小量子上同调相联系的含义。

\begingroup\linespread{1.18}\selectfont
用 $X=\mathbb F_2$ 说明这一点。考虑 $\mathbb F_2$ 的标准扇 $\Sigma$，
\[
\Sigma(1)=\{(1,0),(0,1),(-1,2),(0,-1)\},
\]
如常把一维锥与其本原生成元等同。用 $D_1,D_2,D_3,D_4$ 表示相应环面除子，也用 $D_i$ 表示 $H^2(\mathbb F_2)$ 中的除子类。$\mathbb F_2$ 的 Kähler 锥由纤维类 $f$ 和满足 $H^2=2$ 的截面类 $H$ 张成。用环面除子表示，$D_1=D_3=f$、$D_2=H-2f$、$D_4=H$，因此 $-K_{\mathbb F_2}=\sum_iD_i=2H$，它为 nef。注意 $D_2$ 是 $\mathbb F_2$ 上的自交数为 $-2$ 的曲线。Mori 锥由这条曲线与纤维生成，即由 $H-2f$ 与 $f$ 生成。

须比较函数 $I_{\mathbb F_2}$ 与 $J_{\mathbb F_2}$。因为需要改变坐标，将沿用上一例题的约定，对 $I_{\mathbb F_2}$ 使用 $z_i=e^{t_i}$，对 $J_{\mathbb F_2}$ 使用 $q_i=e^{s_i}$。$H^2(\mathbb F_2)$ 的上同调基为 $T_1=f$、$T_2=H$。

先研究 $I_{\mathbb F_2}$，其公式为
\[
\begin{aligned}
I_{\mathbb F_2}(t_1,t_2)={}&e^{(t_0+t_1f+t_2H)/\hbar}\sum_{\beta\in M(\mathbb F_2)}z^\beta\\
&\times\frac{\prod_{m=-\infty}^0(f+m\hbar)^2\prod_{m=-\infty}^0(H-2f+m\hbar)\prod_{m=-\infty}^0(H+m\hbar)}{\prod_{m=-\infty}^{f\cdot\beta}(f+m\hbar)^2\prod_{m=-\infty}^{(H-2f)\cdot\beta}(H-2f+m\hbar)\prod_{m=-\infty}^{H\cdot\beta}(H+m\hbar)}.
\end{aligned}
\]
如命题~\ref{prop:toric-fano-i-equals-j} 的证明，特别关心其展开中 $\hbar^{-1}$ 的系数。该证明表明，此系数由满足 $-K_{\mathbb F_2}\cdot\beta=0$，且恰有一个 $i$ 满足 $D_i\cdot\beta<0$ 的有效类 $\beta$ 决定。容易计算，唯一的这些类为 $\beta=kD_2$，$k>0$。此时 $D_1\cdot\beta=D_3\cdot\beta=k$、$D_2\cdot\beta=-2k$、$D_4\cdot\beta=0$。

令 $\beta=kD_2=k(H-2f)$，$k\ge1$。计算得 $z^\beta=z_1^{f\cdot\beta}z_2^{H\cdot\beta}=z_1^k$，上述 $I_{\mathbb F_2}$ 公式中 $z_1^k$ 的系数为
\[
\frac{\prod_{m=-2k+1}^0(H-2f+m\hbar)}{\prod_{m=1}^k(f+m\hbar)^2}
=-(H-2f)\frac{(2k-1)!}{(k!)^2}\hbar^{-1}+o(\hbar^{-1}).
\]
因此，
\[
I_{\mathbb F_2}(t_1,t_2)=e^{(t_0+t_1f+t_2H)/\hbar}
\bigl(1-(H-2f)F(z_1)\hbar^{-1}+o(\hbar^{-1})\bigr),
\]
其中
\[
F(z_1)=\sum_{k=1}^{\infty}\frac{(2k-1)!}{(k!)^2}z_1^k.
\]
\SourcePage{417}
定理~\ref{thm:givental-toric-mirror} 要求
\[
f_0=0,\qquad f_1=2F(z_1),\qquad f_2=-F(z_1),\qquad h=0,
\]
随后环面镜像定理给出 $J_{\mathbb F_2}(s_1,s_2)=I_{\mathbb F_2}(t_1,t_2)$，其中
\[
(s_1,s_2)=(t_1+2F(e^{t_1}),t_2-F(e^{t_1})).
\]
置 $z_i=e^{t_i}$、$q_i=e^{s_i}$，即
\begin{equation}\label{eq:hirzebruch-two-mirror-coordinate-change}
(q_1,q_2)=(z_1e^{2F(z_1)},z_2e^{-F(z_1)}).
\end{equation}
令人意外的是，变量变换 $s_1=t_1+2F(z_1)$ 早已作为\hyperref[chap:moduli]{第六章}的式~\eqref{eq:weighted-boundary-log-series} 出现。与式~\eqref{eq:weighted-boundary-log-closed} 比较，得到
\[
s_1=\log\left(\frac{1-2z_1-\sqrt{1-4z_1}}{2z_1}\right),\qquad
q_1=\frac{1-2z_1-\sqrt{1-4z_1}}{2z_1}.
\]
于是 $z_1=q_1/(1+q_1)^2$，再容易得到上述坐标变换的逆为
\begin{equation}\label{eq:hirzebruch-two-inverse-mirror-coordinates}
(z_1,z_2)=\left(\frac{q_1}{(1+q_1)^2},q_2(1+q_1)\right),
\end{equation}
以及
\begin{equation}\label{eq:hirzebruch-two-inverse-log-coordinates}
(t_1,t_2)=\bigl(s_1-2\log(1+e^{s_1}),s_2+\log(1+e^{s_1})\bigr).
\end{equation}
这些坐标变换也出现在 \SourceCite{Givental4} 的一个例子中，只是变量名不同；那里考虑了 $\mathbb P^1$ 上的射影丛
$\mathbb P(\mathcal O_{\mathbb P^1}\oplus\mathcal O_{\mathbb P^1}(-1)\oplus\mathcal O_{\mathbb P^1}(-1))$ 与
$\mathbb P(\mathcal O_{\mathbb P^1}\oplus\mathcal O_{\mathbb P^1}(-2))$。

\looseness=-1 现在计算 $H_\omega^*(\mathbb F_2)$ 和 $\mathbb F_2$ 的小量子上同调中的一些关系。令 $\beta\in H_2(\mathbb F_2,\mathbb Z)$ 为 $f$ 的 Poincaré 对偶，计算算子 $\Box'_{\beta,\hbar}$。令 $\lambda_i$ 对应于 $D_i$，$\delta_i=\lambda_i\partial/\partial\lambda_i$。由式~\eqref{eq:toric-quantum-gkz-operator}，计算得
\[
\Box'_{\beta,\hbar}=\hbar\delta_2\hbar\delta_4-z_2,
\]
因为 $z^\beta=z_2$；回顾现在用 $z_i$ 而非 $q_i$。由式~\eqref{eq:toric-quantum-gkz-coordinate-derivative}，
\[
\hbar\delta_2=\hbar\frac\partial{\partial t_2}-2\hbar\frac\partial{\partial t_1},\qquad
\hbar\delta_4=\hbar\frac\partial{\partial t_2},
\]
随后命题~\ref{prop:nef-i-gkz} 给出
\begin{equation}\label{eq:hirzebruch-two-i-quantum-equation}
\left(\hbar\frac\partial{\partial t_2}\left(\hbar\frac\partial{\partial t_2}-2\hbar\frac\partial{\partial t_1}\right)-z_2\right)I_{\mathbb F_2}(t_1,t_2)=0.
\end{equation}
在 Batyrev 量子环 $H_\omega^*(\mathbb F_2)$ 中，这对应于关系
\begin{equation}\label{eq:hirzebruch-two-batyrev-quantum-relation}
H\mathbin{\cdot}(H-2f)=z_2.
\end{equation}
\looseness=-1 再应用环面镜像定理，考察它对 $J_{\mathbb F_2}$ 的含义。关键是，变量变换给出一个略有不同的微分方程。对式~\eqref{eq:hirzebruch-two-inverse-log-coordinates} 应用链式法则，得到
\[
\left(\hbar\frac\partial{\partial s_2}-2\hbar\frac\partial{\partial s_1}\right)
=\frac{1-q_1}{1+q_1}\left(\hbar\frac\partial{\partial t_2}-2\hbar\frac\partial{\partial t_1}\right),\qquad
\hbar\frac\partial{\partial s_2}=\hbar\frac\partial{\partial t_2}.
\]
\SourcePage{418}
结合这些等式，并利用式~\eqref{eq:hirzebruch-two-mirror-coordinate-change} 中 $q_1$ 与 $t_2$ 无关的事实，得到
\[
\begin{aligned}
&\frac{1-q_1}{1+q_1}\left(\hbar\frac\partial{\partial t_2}\left(\hbar\frac\partial{\partial t_2}-2\hbar\frac\partial{\partial t_1}\right)-z_2\right)\\
&\qquad=\frac{1-q_1}{1+q_1}\hbar\frac\partial{\partial t_2}\left(\hbar\frac\partial{\partial t_2}-2\hbar\frac\partial{\partial t_1}\right)-\frac{1-q_1}{1+q_1}z_2\\
&\qquad=\hbar\frac\partial{\partial s_2}\left(\hbar\frac\partial{\partial s_2}-2\hbar\frac\partial{\partial s_1}\right)-q_2(1-q_1),
\end{aligned}
\]
最后一行用了式~\eqref{eq:hirzebruch-two-inverse-mirror-coordinates} 中的 $z_2=q_2(1+q_1)$。

上式左侧算子消去 $I_{\mathbb F_2}(t_1,t_2)$，所以由定理~\ref{thm:givental-toric-mirror}，右侧算子消去 $J_{\mathbb F_2}(s_1,s_2)$。因此
\[
\hbar\frac\partial{\partial s_2}\left(\hbar\frac\partial{\partial s_2}-2\hbar\frac\partial{\partial s_1}\right)-q_2(1-q_1)
\]
是量子微分算子。由定理~\ref{thm:quantum-differential-relations}，代换
\[
\hbar\partial/\partial s_1\longmapsto T_1=f,\qquad\hbar\partial/\partial s_2\longmapsto T_2=H
\]
给出 $\mathbb F_2$ 的小量子上同调中的关系
\begin{equation}\label{eq:hirzebruch-two-small-quantum-relation}
H*(H-2f)=q_2(1-q_1).
\end{equation}
\looseness=-1 另一方面，式~\eqref{eq:hirzebruch-two-batyrev-quantum-relation} 表明，在 $\mathbb F_2$ 的 Batyrev 量子环中，$H\mathbin{\cdot}(H-2f)=z_2=q_2(1+q_1)$。因此，这两个环并不相同。

\looseness=-1 还应指出，计算 $\mathbb F_2$ 的小量子上同调有一种更初等的方法，因为 $\mathbb F_2$ 是 $\mathbb F_0=\mathbb P^1\times\mathbb P^1$ 的一个形变。再次用 $f$ 表示 $\mathbb F_0$ 的纤维，$s$ 表示由另一份 $\mathbb P^1$ 上的纤维化给出的截面。$\mathbb F_2$ 与 $\mathbb F_0$ 形变等价，所以由形变公理，它们的量子上同调相同。将它们的上同调等同，可见纤维类 $f$ 对应，且 $H=s+f$。$\mathbb P^1\times\mathbb P^1$ 为 Fano，可用 $\mathbb F_0$ 的 Batyrev 量子环计算其小量子上同调，得到关系
\[
s*s=q^f,\qquad f*f=q^s.
\]
再计算
\[
H*(H-2f)=(s+f)*(s-f)=s*s-f*f=q^f-q^s=q^f-q^{H-f}.
\]
因 $q^f=q_1^{f\cdot f}q_2^{H\cdot f}=q_2$，而 $q^{H-f}=q_1^{f\cdot(H-f)}q_2^{H\cdot(H-f)}=q_1q_2$，此计算与式~\eqref{eq:hirzebruch-two-small-quantum-relation} 一致。

\endgroup
\end{example}
